\section{Separation of rows by theta pairings}\label{sec:pf:fs}

\providecommand{\fsT}{\mathcal T}
\providecommand{\fsX}{\mathbb X}
\providecommand{\fsA}{\mathcal A}
\providecommand{\fsE}{\mathcal E}
\providecommand{\fsF}{\mathfrak F}
\providecommand{\fsB}{\mathcal B}
\providecommand{\fsP}[1]{\mathsf{#1}}
\providecommand{\fsc}{\mathrm{c}}
\providecommand{\fst}{\mathrm{t}}
\providecommand{\fsper}{\mathrm{per}}
\providecommand{\fsKF}{\mathrm{KF}}
\providecommand{\fsCols}{\mathrm{Cols}}
\providecommand{\fsReach}{\mathrm{Reach}}
\providecommand{\fsdn}{\mathrm{dn}}
\providecommand{\fslift}{\mathrm{lift}}
\providecommand{\fsRows}{\mathrm{Rows}}
\providecommand{\fsCore}{\mathrm{Core}}
\providecommand{\fsHi}{H^{\infty}}
\providecommand{\fsm}{\mathfrak m}

This section proves \cref{prop:families:fs}: the theta-local inputs $\sfname{ThetaLocalInputs}$, Howe's unramified
statement $\sfname{HoweUnram}$ and the Siegel--Weil formula in the range of absolute convergence
$\sfname{SiegelWeilUnitaryConvergent}$ imply the separation statement $\sfname{FunctionalsSeparate}$, which says that
every non-zero \emph{row form} is detected by a convergent pairing between an iterated theta lift of it and a
\emph{column form} on a totally definite unitary group.

\subsection{Introduction}\label{sec:pf:fs:intro}

\paragraph{The statement.}
Let $D$ be a global datum with CM extension $E/F$ and chain length $m$ (an even integer). A \emph{row form} is a cusp
form $u=t(w)$ on the adelic group $G_{Y_0}$ of a hermitian space $Y_0$ of $D$ (of signature $(r-m/2,m/2)$ at every
archimedean place), obtained by applying a multiplicity map $t$ (an intertwining map from a discrete series model at
the archimedean places to automorphic forms, in a fixed eigenspace $\fsE$) to a smooth vector $w$. Such a form is
lifted along the first row of a grid of hermitian spaces,
\[
Y_0=V_{(0,0)}\longrightarrow V_{(1,0)}\longrightarrow\cdots\longrightarrow V_{(m,0)},
\]
by $m$ successive theta lifts whose Schwartz functions are Fock vectors $\mathrm{fock}(\phi_i(\Phi_i),\rho_i)$ built from
\emph{coset vectors} $\phi_i(\Phi_i)$ and polynomials $\rho_i$; the result $\Theta_{\rho,\Phi}u$ is paired by a
Bessel period with the scalar form $\Psi_\xi$ of a vector-valued \emph{column form} $\Psi$ on the totally definite
group $G_{V_{\fst_0}}$, which embeds into $G_{V_{(m,0)}}$ as the stabiliser of a Bessel datum. The statement
$\sfname{FunctionalsSeparate}$ (\cref{pf:fs:FS}) asserts that, at all large scales, for every non-zero $w$ there are a
polynomial vector $\rho$ and a linear form $\xi$ \emph{independent of $t$} such that for every non-zero $t\in\fsE$ some
column form $\Psi$ and some coset datum $\Phi$ give a convergent and non-zero pairing
\[
f_{\Psi,\Phi}(\rho,\xi;t(w))=\fsB\bigl(\Psi_\xi,\Theta_{\rho,\Phi}\,t(w)\bigr)\neq0 .
\]

\paragraph{The proof.}
The formal proof (\cref{sec:pf:fs:assembly}) combines a dozen statements about the grid, each quantified over all
large-scale settings. In the order in which they are used:
\begin{enumerate}[label=(\arabic*)]
\item \emph{Row tower} (\cref{sec:pf:fs:tower}). At every node $(n,0)$ of the first row there is a discrete series
model, and the $n$-th partial row lift of a multiplicity form is again a multiplicity form, evaluated at a vector
which depends only on $(\rho,w)$. This comes from seesaw intertwiners between Fock vectors and discrete series
models. Consequently all partial lifts converge and decay rapidly, and all column--row pairings converge.
\item \emph{Periods} (\cref{sec:pf:fs:periods}). For every non-zero row form $u$ there is a core form $y(a)$ on the group
of $V_{\fsc_0}=Y_0\oplus\ell$ ($\ell$ a line) with $\int y(a)(\iota_0h)\,u(h)\,dh\ne0$. The core form is a Poincar\'e series of
a matrix coefficient of a discrete series of the larger group; for large level the identity term dominates. The
archimedean input is a statement $\fsP{RI}$ on discrete series of $\Uni(n+1)$ and $\Uni(n)$ with interlacing parameters
(orthogonality relations and decay of restricted matrix coefficients), proved in \cref{sec:pf:fs:real}.
\item \emph{Reachability} (\cref{sec:pf:fs:onto}). Every core form at $\fsc_0$ lies in the span of the functions obtained
from column forms at the top $\fsc_m=\fst_0$ by downward theta lifts along the core column $\fsc_m\to\cdots\to\fsc_0$. This
is a surjectivity statement for composites of theta lifts between finite-dimensional multiplicity spaces, proved by
an adjoint argument from the doubling method, the Rallis inner product, the non-vanishing of partial $L$-values and
Howe's unramified computation.
\item \emph{Ascent of periods} (\cref{sec:pf:fs:ascent}). If a function obtained by downward lifts along the core column
has non-zero period against $u$, then for some $(\rho,\Phi)$ the Bessel pairing at the top of the row lift of $u$ with
the top of the column chain is non-zero. Along the tower, Fourier--Jacobi periods at the nodes $(2k+1,0)$ and Bessel
periods at the nodes $(2k+2,0)$ alternate; each step is a seesaw transfer, and the Fock vectors built on the coset
vectors are numerous enough to realise the transfer.
\item \emph{Uniformity} (\cref{sec:pf:fs:weak,sec:pf:fs:bessel}). Steps (1)--(4) give a weak separation, in which
$(\rho,\xi)$ depend on $t$. Uniqueness of Bessel models for discrete series with a large gap between the
Harish-Chandra parameters, applied at each archimedean place, shows that
$f_{\Psi,\Phi}(\rho,\xi;t(w))=a(\rho,\xi)\,l(t,\Psi,\Phi)$ for a function $a$ independent of $(t,\Psi,\Phi)$; a pair
$(\rho,\xi)$ with $a(\rho,\xi)\ne0$ then works for all $t$.
\end{enumerate}
The individual mechanisms are classical: seesaw pairs and towers of theta lifts \cite{Kudla1994,HarrisKudlaSweet1996},
non-vanishing of periods of Poincar\'e series with a dominating identity term (in the spirit of \cite{Li1992}), the
doubling method and the Rallis inner product formula \cite{GelbartPiatetskiShapiroRallis1987,Rallis1984inj}, the ascent
of Bessel and Fourier--Jacobi periods along towers which underlies the Gan--Gross--Prasad framework and the descent method
\cite{GanGrossPrasad2012,GinzburgRallisSoudry2011}, and the uniqueness of Bessel models \cite{JiangSunZhu2010,SunZhu2012}.
The statement itself, with its grid, its explicit thresholds and its row and column data, does not appear in the
literature.

\paragraph{Organisation.}
\Cref{sec:pf:fs:setting} fixes notation for the objects of the statement, states the separation statement and the
main theorem. \Cref{sec:pf:fs:assembly} lists the intermediate statements and proves the main theorem from them;
\cref{sec:pf:fs:weak,sec:pf:fs:bessel} contain the two final steps in full. The remaining subsections prove the
intermediate statements: the row tower (\cref{sec:pf:fs:tower}), the periods (\cref{sec:pf:fs:periods}), the
reachability (\cref{sec:pf:fs:onto}), the ascent (\cref{sec:pf:fs:ascent}), the Bessel gap (\cref{sec:pf:fs:bur}) and the
archimedean input (\cref{sec:pf:fs:real}). \Cref{sec:pf:fs:aux} proves auxiliary results used in these subsections: the totality of Fock
vectors, the seesaw transfers, the existence of tower frames and of chains along the core column, the cuspidality and the periods of
Poincar\'e series, the existence of packs, and the branching of discrete series. \Cref{sec:pf:fs:remarks} collects remarks on the formal
statements.

The proof is long (about $73\,000$ lines of Lean in the formal development \cite{LeanSource2026}); we give the complete
argument at the top levels and state precisely the lower-level results that are used. Results of the formal
development proved elsewhere are cited by the label of their formal name, with a statement of the form in which they are
used.

% ---------------------------------------------------------------------------
\subsection{The setting and the statement}\label{sec:pf:fs:setting}

\subsubsection{Groups, discrete series models and multiplicity maps}

Fix a global datum $D$ (\cref{sec:families:theta}). It contains a CM extension $E/F$, integer ranks $(k,m,t,r,f,\dots)$ with
$m$ even and $t$ odd, and a finite set $\fsT$ of \emph{archimedean indices}, of cardinality $f$, with an embedding
$\tau\colon E\to\C$ for every $\tau\in\fsT$. For a hermitian space $V$ of $D$ (a hermitian space over $E$ with a standard
lattice) we write $V^\tau$ for its Gram matrix mapped by $\tau$, $\Uni(V^\tau)=\{g\in\GL_{\dim V}(\C):g^*V^\tau g=V^\tau\}$,
$G_V=\Uni(V)(\A_F)$ for the adelic group, $\Gamma_V$ for its rational points, $G_{V,\infty}$ and $G_{V,f}$ for its archimedean
and finite parts, and $-V$ for the space with the opposite form (so that $G_{-V}=G_V$). The objects below depend on a
\emph{family of archimedean embeddings} $\theta$ (a ring homomorphism $\theta_\tau$ from the archimedean part of $\A_E$ to
$\C$ for every $\tau$, with the properties recorded in the formal definition); it gives homomorphisms
$\iota_\tau\colon\Uni(V^\tau)\to G_{V,\infty}$ and components $g\mapsto g_\tau\in\Uni(V^\tau)$. A function on $G_V$ is
\emph{rapidly decreasing} if it is left $\Gamma_V$-invariant, continuous and $O(\mathrm{ht}(g)^{-N})$ for every $N$,
where $\mathrm{ht}$ is the quotient height of $G_V$.

For $\eps\in\{\pm1\}^n$ let $J_\eps$ be the diagonal matrix with entries $\eps$. A \emph{discrete series model}
$\mathcal M$ of $V$ with parameter $S=(S_\tau)_{\tau\in\fsT}$ (a family of finite sets of signed Harish-Chandra
parameters) consists, for every $\tau$, of a Hilbert space $H_\tau$, a matrix $\mathrm{fr}_\tau$ with
$\mathrm{fr}_\tau^*J_{\eps_\tau}\mathrm{fr}_\tau=V^\tau$ and a discrete series representation $\sigma_\tau$ of
$\Uni(J_{\eps_\tau})$ on $H_\tau$ with Harish-Chandra parameter $\lambda_\tau$, where $(\lambda_\tau,\eps_\tau)$ enumerates
$S_\tau$. The representation of $\Uni(V^\tau)$ is $\pi_\tau(g)=\sigma_\tau(\mathrm{fr}_\tau g\,\mathrm{fr}_\tau^{-1})$, and
$c_\tau(\mathcal M)=\mathrm{fr}_\tau^{-1}J_{\eps_\tau}\mathrm{fr}_\tau$ is the \emph{Cartan involution} of the model. We write
$\fsHi_\tau$ for the span of the smooth vectors of $\pi_\tau$, $\fsHi(\mathcal M)=\bigotimes_\tau\fsHi_\tau$ (algebraic tensor
product over $\C$) and $\fsKF(\mathcal M)\subseteq\fsHi(\mathcal M)$ for the span of the pure tensors of $K$-finite vectors.

\begin{definition}[Multiplicity maps]\label{pf:fs:mult}
Let $\mathcal M$ be a discrete series model of $V$, $J\subseteq G_V$ and $\chi\colon G_V\to\C$. A multilinear map $T$ from
$\prod_\tau\fsHi_\tau$ to functions on $G_V$ is a \emph{multiplicity map at level $(J,\chi)$} if for all $v$: $T(v)$ is a
smooth automorphic form and a cusp form; $T(v)(gk)=\chi(k)T(v)(g)$ for $k\in J$; $T$ is equivariant at every $\tau$, i.e.\
replacing $v_\tau$ by $\pi_\tau(h)v_\tau$, $h\in\Uni(V^\tau)$, replaces $T(v)$ by its right translate by $\iota_\tau(h)$;
and $T$ is continuous in each variable for a Sobolev norm of $\pi_\tau$. The space $\fsA(\mathcal M;J,\chi)$ is the span of
the multiplicity maps at level $(J,\chi)$, viewed as linear maps $\fsHi(\mathcal M)\to\{\text{functions on }G_V\}$; its
elements are again called multiplicity maps, and their values \emph{multiplicity forms}.
\formal{Formalism.IndefR.Mult.DSModel}
\formal{Formalism.IndefR.Links.SmT}
\formal{GlobalTheory.SepRows.KF}
\formal{Formalism.IndefR.Mult.IsMultMap}
\formal{Formalism.IndefR.Mult.MultT}
\formal{GlobalTheory.Grow.RapidDec}
\end{definition}

Every value $T(v)$ of a multiplicity map is a smooth automorphic cusp form of rapid decay
(\lref{GlobalTheory.ChainOnto.MultRapid.multT_apply_rapidDec}, which states this for all $T\in\fsA(\mathcal M;J,\chi)$
and all $v$, given the rapid decay of cusp forms $\sfname{CuspRapidDecay}$).

\subsubsection{Grid data}

The objects of the statement live on a \emph{grid datum}. We describe the components which enter.

\paragraph{Level-character data.} A level-character datum $L$ over a type $T$ of \emph{twist indices} consists of Hecke
characters $\mu_t$ ($t\in T$) and $\eta$, finite sets of places of $E$, lattices $\Lambda_Y(j)$ of the spaces
$Y_j=Y_{\min(j,m)}$ of $D$ and $\Lambda_X$ of the space $X'$, frames, and local data at the level places. For
$\mu_0\in T$, the condition $L.\mathrm{Is}(\mu_0)$ is a list of normalisations (the twist at $\mu_0$ is trivial, the
twists are injective, of finite order, tame and trivial on $\A_F^\times$, the lattices $\Lambda_Y(j)$, $j<m$, are self-dual
with respect to frames at level $g_0g_1$, and so on). The \emph{row level} $K_{\mathrm{row}}\subseteq G_{Y_0}$ is the group
of elements whose finite component lies in the principal congruence subgroup of $\Lambda_Y(0)$ modulo $g_0g_1$, where
$g_0,g_1$ are the type primes of $D$.

\paragraph{Nodes, pairs and arrows.} The \emph{nodes} are the grid nodes $(i,j)$, the core nodes $\fsc_i$ and the top nodes
$\fst_j$ ($i,j\in\N$; they are \emph{in range} if $i,j\le m$). The \emph{pairs} are $\chi(i,j)$, joining $(i,j)$ and
$(i+1,j)$; $\mu(s,i,j)$ ($s\in T$), joining $(i,j)$ and $(i,j+1)$; the core pairs $\fsc(i)$, joining $\fsc_i$ and
$\fsc_{i+1}$; and the top pairs, joining $\fst_j$ and $\fst_{j+1}$. The \emph{arrows} are the right arrows
$(i,j)\to(i+1,j)$, the down arrows $(i,j+1)\to(i,j)$, the arrows $\fsc_i\to\fsc_{i+1}$ and $\fsc_{i+1}\to\fsc_i$ along
the core column, and the column arrows $\fst_j\to\fst_{j+1}$; each arrow has an underlying pair.

\begin{definition}[Grid data]\label{pf:fs:grid}
Let $\mathbf c=(c_{\mathrm{gap}},c_{\mathrm{dep}},b)$ be real constants. A \emph{grid datum} $G$ over $L$ at $\mathbf c$
consists of:
\begin{enumerate}[label=(\alph*)]
\item a base place $v_0$ of $F$ above the prime $q_0$ of $D$ and, for $v$ in $S_0=\{v_0\}\cup S_1$ ($S_1$ a finite set of
places attached to $D$), a place $w(v)$ of $E$ above $v$ which splits (with further compatibilities at the level places);
the \emph{modulus} is $\fsm=\prod_{v\in S_0}w(v)^{d(v)}$, with $d(v_0)=d_0(c_{\mathrm{dep}},b)$ (an explicit depth) and
$d(v)=2$ for $v\in S_1$;
\item Hecke characters $\chi_1,\dots,\chi_m$ trivial on $\A_F^\times$, with prescribed infinity types, unramified outside
$S_0$, and of prescribed conductor at $S_0$;
\item hermitian lattices $(V_X,\Lambda_X)$ at all nodes, with $(V_{(0,j)},\Lambda_{(0,j)})=(Y_j,\Lambda_Y(j))$ for $j<m$,
$(V_{(0,m)},\Lambda_{(0,m)})=(X',\Lambda_X)$ and $V_{\fst_0}=V_{\fsc_m}$, such that $V_X$ has, at every $\tau$, $n_X$
negative and $\dim V_X-n_X$ positive directions, where
$n_{(i,j)}=\lfloor(m-j)/2\rfloor$, $n_{\fsc_i}=\lfloor(m-i+1)/2\rfloor$ and $n_{\fst_j}=0$;
\item orthogonal decompositions: $V_{\fsc_0}=Y_0\oplus\ell$ with a line $\ell$, compatible with the lattices; for
$1\le i\le m$, an embedding of $V_{\fsc_i}$ into $V_{(i,0)}$ whose complement has dimension $i-1$ and an antidiagonal Gram
matrix of sign $-1$, with the corresponding lattices; similar decompositions along the last row and at the top.
\end{enumerate}
The decomposition $V_{\fsc_0}=Y_0\oplus\ell$ gives the embedding $\iota_0\colon G_{Y_0}\to G_{V_{\fsc_0}}$ ($h$ acts on $Y_0$ and
trivially on $\ell$). The case $i=m$ of (d) gives the embedding $\iota_{\mathrm{top}}\colon G_{V_{\fst_0}}\to G_{V_{(m,0)}}$ and
a \emph{Bessel datum} at the top: an isotropic family $e=(e_1,\dots,e_{(m-2)/2})$ in the complement of $V_{\fst_0}$, a vector
$z$ of norm $-1$ orthogonal to it, the unipotent radical $N$ of the flag of $e$ and the character
$\psi_N\colon N\to\C^\times$ attached to the additive character $\psi$ below. For a node $X$, $K_X\subseteq G_{V_X}$ is the
group of elements whose finite component lies in the principal congruence subgroup of $\Lambda_X$ modulo $\fsm$, and
$K^+=K_{\fst_0}$. A \emph{coset datum} is a family $\Phi=(\Phi_i)_{i<m}$ with $\Phi_i$ in the finite set
$\mathcal C_i=\mathrm{Hom}(\Lambda_{(i+1,0)},\Lambda_{(i,0)})/\fsm\,\mathrm{Hom}(\Lambda_{(i+1,0)},\Lambda_{(i,0)})$; a \emph{polynomial vector} is a family $\rho=(\rho_i)_{i<m}$ of
polynomials in the entries of $\dim V_{(i,0)}\times\dim V_{(i+1,0)}$ matrices.
\formal{GlobalTheory.ChiGrid.Consts}
\formal{GlobalTheory.ChiGrid.Places}
\formal{GlobalTheory.ChiGrid.Chars}
\formal{GlobalTheory.ChiGrid.Vertices}
\formal{GlobalTheory.ChiGrid.Node}
\formal{GlobalTheory.ChiGrid.Pair}
\formal{GlobalTheory.ChiGrid.Arrow}
\formal{GlobalTheory.ChiGrid.GridData}
\formal{GlobalTheory.ChiGrid.GridData.besTop}
\formal{GlobalTheory.ChiGrid.GridData.levCong}
\formal{GlobalTheory.ChiGrid.GridData.Datum}
\formal{GlobalTheory.ChiGrid.GridData.RealVec}
\formal{GlobalTheory.SepRows.levRow}
\formal{GlobalTheory.SepRows.lev}
\formal{GlobalTheory.FunObj.stabilizerEmbedding}
\formal{GlobalTheory.FunObj.embTop}
\end{definition}

For a pair $p$ we write $(S_p,T_p)$ for the two spaces on which its Weil representation acts: $S_p=-V_{(i+1,j)}$ and
$T_p=V_{(i,j)}$ for $p=\chi(i,j)$, $S_p=-V_{\fsc_{i+1}}$ and $T_p=V_{\fsc_i}$ for $p=\fsc(i)$, and $(S_p,T_p)=(V_{\text{small}},
V_{\text{big}})$ for the other pairs. Theta functions on $\fsX(S_p,T_p)=\mathrm{Mat}_{\dim T_p\times\dim S_p}(\A_E)$, Weil data
and theta kernels $\theta_\omega(g,h;f)=(\omega(g,h)f)(0)$ are as in \cref{pf:liftinj:thetadef}.

\subsubsection{Weil data, volume data and lifts}

\begin{definition}[Weil data and coset vectors]\label{pf:fs:weil}
A \emph{Weil datum} $W$ of $G$ consists of a standard additive character $\psi$ of $\A_F/F$; a family of archimedean
embeddings $\theta$; Cartan involutions $c_X^\tau$ of $V_X^\tau$ and orthonormal frames adapted to them, for all nodes in
range; for every pair $p$ in range a Weil datum $\omega_p$ for $(S_p,T_p)$ whose splitting characters are attached to
the character of $p$ (that is $\chi_{m-i}$ for $p=\chi(i,j)$ and $p=\fsc(i)$, and the twisting characters of $L$ for the
other pairs) and to $\eta$; for $i<m$ a \emph{vacuum vector} $\phi_i^0$ for the pair $\chi(i,0)$, which is a Fock vacuum
for the Cartan involutions and a lattice vacuum for $\Lambda_{(i+1,0)},\Lambda_{(i,0)}$; and for $\kappa\in\mathcal C_i$ a
\emph{coset vector} $\phi_i(\kappa)$: for some idempotent $e$ of $\A_E$ supported at the places $w(v)$, $v\in S_0$, and every
local model at $e$ which carries a multiple of the characteristic function of the lattice to $\phi_i^0$, the coset vector is
the image of the characteristic function of the coset $\kappa$. For a theta function $\phi$ of a pair $p$ and a polynomial $P$, $\mathrm{fock}_p(\phi,P)$ is the
Fock vector obtained from $\phi$ by the creation operators given by $P$ (with respect to the Cartan involutions and frames
of $W$); it is linear in $P$. The \emph{row datum} of $(\rho,\Phi)$ is
\[
\phi_i(\rho,\Phi)=\mathrm{fock}_{\chi(i,0)}\bigl(\phi_i(\Phi_i),\rho_i\bigr),\qquad 0\le i<m .
\]
For a real datum $R$ (\cref{pf:fs:rowform} below), $W$ is \emph{over} $R$ if $\psi$ and $\theta$ are those of $R$, the Cartan
involutions and frames at the nodes $(0,j)$ are those of the discrete series models of $R$, and the Weil representations of
the pairs $\mu(s,0,j)$ are those of $R$.
\formal{GlobalTheory.ChiGrid.GridData.Weil}
\formal{GlobalTheory.ChiGrid.GridData.Weil.fock}
\formal{GlobalTheory.ChiGrid.GridData.Weil.datum}
\formal{GlobalTheory.ChiGrid.GridData.Weil.Over}
\end{definition}

\begin{definition}[Volume data, lifts and the Bessel pairing]\label{pf:fs:vol}
A \emph{volume datum} $M$ of $G$ consists of measures $\nu_X$ on $G_{V_X}$ and subsets $\fsF_X\subseteq G_{V_X}$ for all nodes, and
a measure $\nu_N$ and a subset $\fsF_N$ of $N$. It is a \emph{row} if for $i\le m$ the measure $\nu_{(i,0)}$ is a Haar measure,
$\fsF_{(i,0)}$ is a fundamental domain for $\Gamma_{V_{(i,0)}}$ and the pair is normalised with respect to $K_{(i,0)}$; the same
holds at $\fst_0$ with respect to $K^+$; and $\nu_N$ is a Haar measure on $N$, $\fsF_N$ a fundamental domain for $N(F)$, and
$\nu_N(\fsF_N)=1$.
\begin{enumerate}[label=(\alph*)]
\item For an arrow $a$ with source $X$ and target $Y$, the \emph{theta kernel} $k_a(f;g,h)$, $g\in G_{V_X}$, $h\in G_{V_Y}$, is
the theta kernel of $\omega_p$ ($p$ the pair of $a$) at the pair of group elements determined by $(g,h)$: for a right arrow
$(i,j)\to(i+1,j)$ it is $\theta_{\omega_p}(h,g;f)$ ($h$ viewed in $G_{-V_{(i+1,j)}}$), and for the arrow
$\fsc_{i+1}\to\fsc_i$ it is $\theta_{\omega_p}(g,h;f)$. The \emph{lift} along $a$ is
\begin{equation}\label{eq:pf:fs:lift}
\fslift_a(f,F)(h)=\int_{\fsF_X}k_a(f;g,h)\,F(g)\,d\nu_X(g),
\end{equation}
and the lift \emph{converges} if the integrand is $\nu_X$-integrable on $\fsF_X$ for every $h$.
\item The \emph{partial row lifts} of a function $u$ on $G_{Y_0}$ are
$\Theta^{(0)}_{\rho,\Phi}u=u$ and $\Theta^{(n+1)}_{\rho,\Phi}u=\fslift_{(n,0)\to(n+1,0)}(\phi_n(\rho,\Phi),\Theta^{(n)}_{\rho,\Phi}u)$
for $n<m$; the \emph{row lift} is $\Theta_{\rho,\Phi}u=\Theta^{(m)}_{\rho,\Phi}u$, a function on $G_{V_{(m,0)}}$. The row lift
\emph{converges} if each of the $m$ lifts converges.
\item The \emph{Bessel pairing} of functions $\varphi$ on $G_{V_{\fst_0}}$ and $F$ on $G_{V_{(m,0)}}$ is
\begin{equation}\label{eq:pf:fs:bessel}
\fsB(\varphi,F)=\int_{\fsF_{\fst_0}}\int_{\fsF_N}\varphi(h)\,F\bigl(n\,\iota_{\mathrm{top}}(h)\bigr)\,\overline{\psi_N(n)}\,
d\nu_N(n)\,d\nu_{\fst_0}(h),
\end{equation}
and it \emph{converges} if the inner integrand is integrable for every $h$ and the inner integral is integrable on
$\fsF_{\fst_0}$.
\end{enumerate}
\formal{GlobalTheory.FunObj.Vol}
\formal{GlobalTheory.FunObj.Vol.IsRow}
\formal{GlobalTheory.ChiGrid.GridData.Weil.kerA}
\formal{GlobalTheory.FunObj.liftA}
\formal{GlobalTheory.FunObj.LiftConv}
\formal{GlobalTheory.SepRows.part}
\formal{GlobalTheory.FunObj.thetaChi}
\formal{GlobalTheory.FunObj.ThetaConv}
\formal{GlobalTheory.FunObj.Bplus}
\formal{GlobalTheory.FunObj.BplusConv}
\end{definition}

\begin{definition}[Column forms and the column--row pairing]\label{pf:fs:cols}
The \emph{top coefficient datum} $c^+$ of $G$ and $\theta$ is a finite-dimensional complex vector space $W^+$ with a
representation $\rho^+$ of $G_{V_{\fst_0}}$ which factors through the archimedean components; it is built from dominant
weights at the archimedean indices and the dual Gram matrices of $V_{\fst_0}$, which are positive definite. A
\emph{column form} is a function $\Psi\colon G_{V_{\fst_0},f}\to W^+$ with $\Psi(\gamma gk)=\rho^+(\gamma)\Psi(g)$ for
$\gamma\in\Gamma_{V_{\fst_0}}$ and $k\in K^+$; the column forms form a space $\fsCols$. For a linear form $\xi$ on $W^+$ the
\emph{scalar form} of $\Psi$ is the function on $G_{V_{\fst_0}}$
\[
\Psi_\xi(h)=\xi\bigl(\rho^+(h)^{-1}\Psi(h_f)\bigr).
\]
For a polynomial vector $\rho$, a linear form $\xi$, $\Psi\in\fsCols$, a coset datum $\Phi$ and a function $u$ on $G_{Y_0}$, the
\emph{column--row pairing} is
\begin{equation}\label{eq:pf:fs:pairing}
f_{\Psi,\Phi}(\rho,\xi;u)=\fsB\bigl(\Psi_\xi,\Theta_{\rho,\Phi}u\bigr),
\end{equation}
and $\mathrm{FConv}(\rho,\xi,\Psi,\Phi;u)$ means that the row lift $\Theta_{\rho,\Phi}u$ converges and the Bessel pairing in
\eqref{eq:pf:fs:pairing} converges.
\formal{GlobalTheory.IndefiniteRowSideFullHecke.LevelPlusSubgroup.coefPlus}
\formal{GlobalTheory.IndefiniteRowSideFullHecke.LevelPlusSubgroup.Cols}
\formal{GlobalTheory.FunObj.pairVec}
\formal{GlobalTheory.FunObj.fPsiPhi}
\formal{GlobalTheory.FunObj.FConv}
\end{definition}

\subsubsection{Real data, row forms and the statement}

\begin{definition}[Row forms]\label{pf:fs:rowform}
Let $R$ be a real datum over $(D,\mathcal A,K_s)$ (the real data of \cref{prop:families:liftinj}, described in
\cref{sec:pf:liftinj:realdata}). It contains a level-character datum $L_R$, a family
$\theta_R$ and a character $\psi_R$, and discrete series models $\mathcal M_Y(j)$ of $Y_j$. The \emph{eigenspace} $\fsE_R$ is the
space of multiplicity maps $t$ of $\mathcal M_Y(0)$ at the level of $R$ for the trivial level datum which are fixed by the
group-algebra idempotent $e_0$ of $R$. For $t\in\fsE_R$ and $w\in\fsHi(\mathcal M_Y(0))$, the \emph{row form} $t(w)$ is the value
at $w$ of the linear map on the tensor product induced by $t$.
\formal{GlobalTheory.IndefiniteRowSideFullHecke.LevelPlusSubgroup.rowForm}
\formal{Formalism.IndefR.Links.RealDat.EEat}
\formal{Formalism.IndefR.Links.LevCh.triv}
\end{definition}

The statement concerns size data $L_{\mathrm{sz}}$ of $D$ (\cref{sec:families:theta}) with scale $(x,\eps)$ and $b=e^{x^\eps}$,
subject to the size conditions of rows $\sfname{SizeRows}(L_{\mathrm{sz}})$: a conjunction of thirteen explicit conditions on the
ranks of $D$, its host and the scale, among them $44\le m$, $m^4\le f$, $fm\le(r+1)^3$, $mfN(m)^2d\le b$ and the existence of a
parameter choice of the prescribed kind; and the condition $\sfname{ESerTie}(D)$ on the type data at $g_0,g_1$.

\begin{definition}[Separation by functionals]\label{pf:fs:FS}
Let $R$ be a real datum, $G$ a grid datum over $L_R$, $W$ a Weil datum of $G$, $M$ a volume datum of $G$ and
$w\in\fsHi(\mathcal M_Y(0))$. A polynomial vector $\rho$ and a linear form $\xi$ on $W^+$ \emph{separate} at $w$ if for every
non-zero $t\in\fsE_R$ there are $\Psi\in\fsCols$ and a coset datum $\Phi$ with
\[
\mathrm{FConv}(\rho,\xi,\Psi,\Phi;t(w))\quad\text{and}\quad f_{\Psi,\Phi}(\rho,\xi;t(w))\neq0 .
\]
$\sfname{FunctionalsSeparate}$ is the statement: for all real $\eps$ and $C_F\ge6510$ there is $x_0$ such that for every global
datum $D$ and size data $L_{\mathrm{sz}}$ of $D$ with exponent $\eps$, scale $x\ge x_0$ and $\sfname{SizeRows}(L_{\mathrm{sz}})$, if
$\sfname{ESerTie}(D)$ holds, then for every finite group $\Gamma$, types $\mathrm{Pl}_0,T$, family input $\mathcal A$, number field
$K_s$ and real datum $R$ over $(D,\mathcal A,K_s)$ with $L_R.\mathrm{Is}(\mu_0)$ ($\mu_0$ the base tuple of $\mathcal A$) and the group
record of $R$, every grid datum $G$ over $L_R$ at $(1,C_F,b)$, every Weil datum $W$ of $G$ over $R$, every volume datum $M$ of $G$
which is a row, and every non-zero $w\in\fsHi(\mathcal M_Y(0))$, some $(\rho,\xi)$ separates at $w$.
\formal{GlobalTheory.IndefiniteRowSideFullHecke.LevelPlusSubgroup.Separates}
\formal{GlobalTheory.IndefiniteRowSideFullHecke.LevelPlusSubgroup.FunctionalsSeparate}
\formal{GlobalTheory.IndefiniteRowSideFullHecke.LevelPlusSubgroup.SizeRows}
\formal{GlobalTheory.GLq.ESerTie}
\formal{Formalism.IndefR.Links.RealDat.grp_record}
\end{definition}

The order of the quantifiers matters: $(\rho,\xi)$ is chosen before $t$, while $(\Psi,\Phi)$ may depend on $t$. If $\fsE_R=0$
the statement holds trivially. If $\fsE_R\ne0$ the form $\xi$ must be non-zero: the formal development records, as an example
attached to the definition, that a separating pair with $\xi=0$ forces $\fsE_R=0$ (the scalar form of $\Psi$ for $\xi=0$
vanishes).

\begin{theorem}[\cref{prop:families:fs}]\label{pf:fs:main}
Assume $\sfname{ThetaLocalInputs}$, $\sfname{HoweUnram}$ and $\sfname{SiegelWeilUnitaryConvergent}$. Then
$\sfname{FunctionalsSeparate}$ holds.
\formallabel{UpperTheory.Lit.functionalsSeparate_of_howeUnram_of_siegelWeil}
\end{theorem}

The theorem is deduced in \cref{sec:pf:fs:assembly} from \cref{pf:fs:mainRI} and \cref{pf:fs:RIholds}.

% ---------------------------------------------------------------------------
\subsection{The intermediate statements and the proof of the main theorem}\label{sec:pf:fs:assembly}

\subsubsection{Settings}
All intermediate statements quantify over \emph{settings}. The primes $p_{de},p'_{de}$ of the host of $D$ are \emph{off $Q$} if
both are prime to $Q_Sq$ (the product of the small primes of the host and the distinguished prime $q$); $\sfname{Sizes}(D)$ is
the condition
\begin{equation}\label{eq:pf:fs:sizes}
1800<r\quad\text{and}\quad b_\tau(j+1)-b_\tau(j)>2N^3\quad(\tau\in\fsT,\ j+1<m),
\end{equation}
where $b_\tau$ are the nodes of the Harish-Chandra data of $D$ at $\tau$ and $N=k+m(t+r+7)$.

\begin{definition}[Settings and three quantifiers]\label{pf:fs:settings}
A \emph{setting} is a tuple $(D,T,L,\mu_0,\mathbf c,G,W,M)$ where $D$ is off $Q$, $L$ is a level-character datum over $T$ with
$L.\mathrm{Is}(\mu_0)$, $G$ is a grid datum over $L$ at $\mathbf c$, $W$ a Weil datum of $G$ and $M$ a volume datum of $G$ which
is a row. For a property $Q(W,M)$ of settings:
\begin{enumerate}[label=(\alph*)]
\item $\forall_{\mathrm{all}}\,Q$ means that $Q$ holds in every setting;
\item $\forall_{\mathrm{gap}}\,Q$ means that $Q$ holds in every setting with $\sfname{Sizes}(D)$ and $c_{\mathrm{gap}}=1$;
\item $\forall_{\mathrm{ev}}\,Q$ means: for all $\eps$ and $C_F\ge6510$ there is $x_0$ such that $Q$ holds for every $D$ with size
data of exponent $\eps$, scale $x\ge x_0$ and $\sfname{SizeRows}$, and every $(T,L,\mu_0,G,W,M)$ with $L.\mathrm{Is}(\mu_0)$, $G$
a grid datum at $(1,C_F,b)$, and $M$ a row.
\end{enumerate}
\formal{GlobalTheory.SepRows.OffQ}
\formal{GlobalTheory.IndefiniteRowSideFullHecke.CuspAlongSmoothLifts.Sizes}
\formal{GlobalTheory.SepRows.All}
\formal{GlobalTheory.SepRows.Gap}
\formal{GlobalTheory.SepRows.Ev}
\end{definition}

The condition $\sfname{SizeRows}$ implies that $D$ is off $Q$ (\cref{pf:fs:offQ} below). Hence $\forall_{\mathrm{all}}Q$ implies
$\forall_{\mathrm{gap}}Q$, and $\forall_{\mathrm{gap}}Q$ implies $\forall_{\mathrm{ev}}Q$ once $\sfname{SizeRows}$ is known to imply
$\sfname{Sizes}$ at large scale (statement $\fsP{SG}$ below): given $\eps,C_F$, take the threshold of $\fsP{SG}$ for $\eps$.

\begin{lemma}\label{pf:fs:offQ}
\begin{enumerate}[label=(\alph*)]
\item If $\sfname{SizeRows}$ holds for size data of $D$, then $D$ is off $Q$.
\item For every constant $c$, the row parameter $S^{\mathrm{row}}_0$ at $(0,0)$ (defined below) equals the dual twisted
parameter $S^Y$ at $j=0$.
\end{enumerate}
\formallabel{GlobalTheory.SepRows.offQ_of_sizeRows}
\formallabel{GlobalTheory.SepRows.parRow_zero}
\end{lemma}

\begin{proof}
(a) The parameter choice contained in $\sfname{SizeRows}$ asserts that $p_{de}$ and $p'_{de}$ are primes of the class
\emph{split in $K$}; such a prime is prime to $Q_Sq$ (\lref{UpperTheory.Ram.Fun.coprime_QSq_of_splitK}). (b) The grid
parameter at $(0,j)$ is the twisted parameter at $j$ (\lref{GlobalTheory.ChiGrid.parV_zero}); both sides of (b) are its
dual.
\end{proof}

\subsubsection{The objects of the intermediate statements}

Fix a setting. Let $S^{\mathrm{row}}_i$ be the parameter of the node $(i,0)$ (the dual of the grid parameter at $(i,0)$ for the
constant $c_{\mathrm{gap}}$), $S^{\fsc}_i$ the parameter of the core node $\fsc_i$, and $S^Y=S^{\mathrm{row}}_0$ the dual twisted
parameter at $j=0$ (\cref{pf:fs:offQ}). A \emph{row model} is a discrete series model $\mathcal M_0$ of
$Y_0$ with parameter $S^Y$ and Cartan involution $c_{(0,0)}$ (that of $W$). The \emph{rows} of $\mathcal M_0$ are the multiplicity
maps
\[
\fsRows(\mathcal M_0)=\fsA(\mathcal M_0;K_{\mathrm{row}},1),
\]
and the \emph{core forms} at $\fsc_i$ of a model $\mathcal M$ of $V_{\fsc_i}$ are the elements of $\fsA(\mathcal M;K_{\fsc_i},1)$.

\begin{definition}[Depth vacua, column chains and periods]\label{pf:fs:chain}
Fix a setting.
\begin{enumerate}[label=(\alph*)]
\item For a pair $p$, the \emph{depth lattice} $\Gamma_p$ is the group of $u\in\fsX(S_p,T_p)$ in the lattice of
$\mathrm{Hom}(\Lambda_{S_p},\Lambda_{T_p})$ such that $eu$ lies in the closure of $\fsm\,\mathrm{Hom}(\Lambda_{S_p},\Lambda_{T_p})$
for every idempotent $e$ of $\A_E$ supported at the places $w(v)$, $v\in S_0$. A \emph{depth vacuum} for $p$ is a theta function
for $\omega_p$ which is fixed by the Heisenberg translations by $\Gamma_p$.
\item A \emph{core volume} $v$ is a family $(\mu_i,\fsF_i)_{i}$ of Haar measures on $G_{V_{\fsc_i}}$ and fundamental domains for
$\Gamma_{V_{\fsc_i}}$. The \emph{downward lift} along the core column of $F$ on $G_{V_{\fsc_{i+1}}}$ by a theta function $f$ of
$\omega_{\fsc(i)}$ is
\[
\fsdn_i(f,F)(g)=\int_{\fsF_{i+1}}k_{\fsc_{i+1}\to\fsc_i}(f;h,g)\,F(h)\,d\mu_{i+1}(h).
\]
\item A \emph{column chain} for $v$ is a family $(\Psi_i)_{i\in\N}$ of functions on $G_{V_{\fsc_i}}$ such that $\Psi_i$ is rapidly
decreasing for $i\le m$ and, for $i<m$, $\Psi_i=\fsdn_i(f_i,\Psi_{i+1})$ for some depth vacuum $f_i$ for $\fsc(i)$. The set
$\fsReach(v)$ of \emph{reachable} functions consists of the $\Psi_0$ of the column chains whose top $\Psi_m$ is a scalar form
$\Psi'_\xi$ of a column form $\Psi'\in\fsCols$ (recall $V_{\fsc_m}=V_{\fst_0}$).
\item The \emph{core period} of $F$ on $G_{V_{\fsc_0}}$ and $u$ on $G_{Y_0}$ is
$\fsper_0(F,u)=\int_{\fsF_{(0,0)}}F(\iota_0h)\,u(h)\,d\nu_{(0,0)}(h)$.
\item A function $u$ on $G_{Y_0}$ is \emph{row-rapid} if all partial row lifts $\Theta^{(n)}_{\rho,\Phi}u$, $n\le m$, are
rapidly decreasing.
\end{enumerate}
\formal{GlobalTheory.SepRows.depLat}
\formal{GlobalTheory.SepRows.Dep}
\formal{GlobalTheory.SepRows.CoreVol}
\formal{GlobalTheory.SepRows.dn}
\formal{GlobalTheory.SepRows.IsColChain}
\formal{GlobalTheory.SepRows.Reach}
\formal{GlobalTheory.SepRows.per0}
\formal{GlobalTheory.SepRows.RapidRow}
\formal{GlobalTheory.SepRows.Rows}
\formal{GlobalTheory.SepRows.CoreForms}
\formal{GlobalTheory.SepRows.parRow}
\end{definition}

\begin{definition}[The intermediate statements]\label{pf:fs:statements}
In each item, $\mathcal M_0$ is a row model, and the quantifier is the one indicated.
\begin{enumerate}[label=(\arabic*),leftmargin=*]
\item $\fsP{SG}$ (sizes give gaps): for every $\eps$ there is $x_0$ such that $\sfname{SizeRows}$ at scale $x\ge x_0$ implies
$\sfname{Sizes}(D)$.
\item $\fsP{RM}$ (row models), $\forall_{\mathrm{gap}}$: for $i\le m$ there is a model of $V_{(i,0)}$ with parameter
$S^{\mathrm{row}}_i$ and Cartan involution $c_{(i,0)}$.
\item $\fsP{CM}$ (core models), $\forall_{\mathrm{all}}$: for $i\le m$ there is a model of $V_{\fsc_i}$ with parameter $S^{\fsc}_i$
and Cartan involution $c_{\fsc_i}$.
\item $\fsP{PE}$ (seesaw pairs): for every $D$, every pair of families of parameters $(S,S')$ and exponents $(a_S,a_T)$ which are
matched at every $\tau$, all spaces $V,W'$ with models $\mathcal M$ of $V$ (parameter $S$), $\mathcal M'$ of $W'$ (parameter $S'$),
every $\theta$ and $\psi$, there is a bilinear map $I$ from polynomials and $\fsHi(\mathcal M')$ to $\fsHi(\mathcal M)$ such that for
every Weil datum $\omega$ for $(V,W')$ with exponents $(a_S,a_T)$ and every vacuum vector $f$ of $(\mathcal M,\mathcal M')$, the map
$I$ is an intertwiner for $\omega$ and the Fock vectors over $f$ (it realises the theta pairing of the Fock vectors over $f$).
\item $\fsP{RS}$ (row step), $\forall_{\mathrm{gap}}$: for $i<m$ and models $\mathcal M_A$ of $V_{(i,0)}$, $\mathcal M_B$ of
$V_{(i+1,0)}$ with parameters $S^{\mathrm{row}}_i,S^{\mathrm{row}}_{i+1}$ and the Cartan involutions of $W$, there is a map $I$
from polynomials and $\fsHi(\mathcal M_A)$ to $\fsHi(\mathcal M_B)$ such that for every $F\in\fsA(\mathcal M_A;\emptyset,1)$ and
$\kappa\in\mathcal C_i$ there is $F'\in\fsA(\mathcal M_B;\emptyset,1)$ with: for all polynomials $P$ and all $v$, the lift of $F(v)$
along $(i,0)\to(i+1,0)$ by $\mathrm{fock}(\phi_i(\kappa),P)$ converges and equals $F'(I(P,v))$.
\item $\fsP{RT}$ (row tower), $\forall_{\mathrm{gap}}$: for $n\le m$ there are a model $\mathcal M_n$ of $V_{(n,0)}$ with
parameter $S^{\mathrm{row}}_n$ and Cartan involution $c_{(n,0)}$ and a map $\mathrm{vec}_n(\rho,w)\in\fsHi(\mathcal M_n)$
($\rho$ a polynomial vector, $w\in\fsHi(\mathcal M_0)$) such that for every $t\in\fsA(\mathcal M_0;\emptyset,1)$ and every coset
datum $\Phi$ there is $F'\in\fsA(\mathcal M_n;\emptyset,1)$ with: for all $\rho,w$, the first $n$ row lifts of $t(w)$ converge and
$\Theta^{(n)}_{\rho,\Phi}\,t(w)=F'(\mathrm{vec}_n(\rho,w))$.
\item $\fsP{RC}$ (row convergence), $\forall_{\mathrm{gap}}$: $\mathrm{FConv}(\rho,\xi,\Psi,\Phi;t(w))$ for all
$t\in\fsRows(\mathcal M_0)$, $w$, $\rho$, $\xi$, $\Psi\in\fsCols$, $\Phi$.
\item $\fsP{RR}$ (rapid rows), $\forall_{\mathrm{gap}}$: for all $t\in\fsRows(\mathcal M_0)$ and $w$, the row form $t(w)$ is
row-rapid and all its row lifts converge.
\item $\fsP{PN}$ (period non-vanishing), $\forall_{\mathrm{ev}}$: for every model $\mathcal M'_0$ of $Y_0$ with parameter $S^Y$ (no
condition on its Cartan involution), every model $\mathcal M_{\fsc}$ of $V_{\fsc_0}$ with parameter $S^{\fsc}_0$, every non-zero
$t\in\fsRows(\mathcal M'_0)$ and every non-zero $w$, there are a core form $y$ of $\mathcal M_{\fsc}$ and $a\in\fsKF(\mathcal M_{\fsc})$
with $\fsper_0(y(a),t(w))\neq0$.
\item $\fsP{CO}$ (core forms are reachable), $\forall_{\mathrm{gap}}$: for every core volume $v$, every model $\mathcal M_{\fsc}$ of
$V_{\fsc_0}$ with parameter $S^{\fsc}_0$ and Cartan involution $c_{\fsc_0}$, every core form $y$ of $\mathcal M_{\fsc}$ and every
$a\in\fsKF(\mathcal M_{\fsc})$, the function $y(a)$ lies in the span of $\fsReach(v)$.
\item $\fsP{DF}$ (depth vacua are fixed), $\forall_{\mathrm{all}}$: for $i<m$, $p\in\{\chi(i,0),\fsc(i)\}$, every depth vacuum $f$
for $p$ and all $k\in K_{\mathrm{small}(p)}$, $k'\in K_{\mathrm{big}(p)}$: $\omega_p(k,k')f=f$ (with $(k,k')$ placed in
$G_{S_p}\times G_{T_p}$ as for the theta kernel).
\item $\fsP{AS}$ (ascent), $\forall_{\mathrm{all}}$: for every core volume $v$, every column chain $\Psi$ and every row-rapid $u$
with $\fsper_0(\Psi_0,u)\ne0$, there are $\rho,\Phi$ with $\fsB(\Psi_m,\Theta_{\rho,\Phi}u)\neq0$.
\item $\fsP{WS}$ (weak separation), $\forall_{\mathrm{ev}}$: for every non-zero $t\in\fsRows(\mathcal M_0)$ and every non-zero $w$
there are $\rho,\xi,\Psi\in\fsCols,\Phi$ with $f_{\Psi,\Phi}(\rho,\xi;t(w))\neq0$.
\item $\fsP{GR}$ (good rows): for every real datum $R$ with $L_R.\mathrm{Is}(\mu_0)$, every $t\in\fsE_R$ defines an element of
$\fsRows(\mathcal M_Y(0))$ (for $\theta_R$ and $L_R$).
\end{enumerate}
\formal{GlobalTheory.SepRows.SizeGap}
\formal{GlobalTheory.SepRows.RowModels}
\formal{GlobalTheory.SepRows.CoreModels}
\formal{GlobalTheory.SepRows.PairEx}
\formal{GlobalTheory.SepRows.RowStep}
\formal{GlobalTheory.SepRows.RowTower}
\formal{GlobalTheory.SepRows.RowConv}
\formal{GlobalTheory.SepRows.RowRapid}
\formal{GlobalTheory.SepRows.PeriodNe}
\formal{GlobalTheory.SepRows.CoreOnto}
\formal{GlobalTheory.SepRows.DepFixed}
\formal{GlobalTheory.SepRows.Ascent}
\formal{GlobalTheory.SepRows.WeakSep}
\formal{GlobalTheory.SepRows.GoodRows}
\end{definition}

In $\fsP{RS}$ and $\fsP{RT}$ the maps $I$ and $\mathrm{vec}_n$ do not depend on $F,\kappa$ resp.\ $t,\Phi$, while $F'$ does not
depend on $(P,v)$ resp.\ $(\rho,w)$: the partial lift of a row form is a multiplicity form evaluated at a vector depending only
on the polynomial data and on $w$. This separation of variables is what makes the final factorisation possible.

The archimedean input is the following statement about pairs of discrete series. For $n\ge1$ and a bijection
$e\colon\{\ast\}\sqcup\{1,\dots,n\}\to\{1,\dots,n+1\}$, let $h\mapsto\tilde h$ be the embedding of $n\times n$ into
$(n+1)\times(n+1)$ matrices placing $h$ on the coordinates $e(1),\dots,e(n)$ and $1$ at $e(\ast)$, $\|g\|$ the Hilbert--Schmidt
norm and $\mathrm{col}(g)$ the norm of the column $e(\ast)$ of $g$. For signs $\eps\in\{\pm1\}^{n+1}$, $\eps'\in\{\pm1\}^n$,
parameters $\lambda\in\R^{n+1}$, $\mu\in\R^n$ and $B\ge0$, $\mathrm{Par}(e,\eps,\lambda,\eps',\mu,B)$ is the list of conditions:
$\eps_{e(\ast)}=+1$ and $\eps_{e(i)}=\eps'_i$; $n\ge3$; $\lambda,\mu$ are Harish-Chandra parameters; $(\eps,\lambda)$ and
$(\eps',\mu)$ are \emph{interlaced} (between two consecutive elements of $\lambda\cup\mu$ the signs agree exactly when the two
elements belong to different families), also after lowering any one entry of $\lambda$ by $1$; consecutive entries of $\lambda$
and of $\mu$ differ by at least $2$; the gaps of $(\eps,\lambda)$ and $(\eps',\mu)$ are at least $1000(n+1)^3$; $\eps'$ has a
negative entry and at most as many negative as positive entries; and $|\lambda_a|,|\mu_i|\le B$. Put
$\mathrm{big}(C,n,B)=\exp(C(n+1)^4\log(8(n+1)B))$.

\begin{definition}[The archimedean input]\label{pf:fs:RI}
$\fsP{RI}$ is the statement: there is $C>0$ such that for all $(n,e,\eps,\lambda,\eps',\mu,B)$ with
$\mathrm{Par}(e,\eps,\lambda,\eps',\mu,B)$ and all discrete series $\pi$ of $\Uni(J_\eps)$ with parameter $\lambda$ on $H$ and
$\sigma$ of $\Uni(J_{\eps'})$ with parameter $\mu$ on $H'$, there are a unit vector $v_0$ in the lowest $K$-type of $\pi$,
$g_0\in\Uni(J_\eps)$, a smooth unit vector $w'\in H'$ and $c>0$ such that, with $v=\pi(g_0)v_0$ (a smooth vector) and the Haar
measure $dh$ of $\Uni(J_{\eps'})$ fixed in the formal statement:
\begin{enumerate}[label=(\roman*)]
\item $\int\overline{\langle v,\pi(\tilde h)v\rangle}\,\langle b,\sigma(h)w'\rangle\,dh=c\,\langle w',w'\rangle\langle b,w'\rangle$ for
all $b\in H'$;
\item for $g\in\Uni(J_\eps)$, $h\mapsto\langle v,\pi(g\tilde h)v\rangle$ is integrable and
\[
\int|\langle v,\pi(g\tilde h)v\rangle|\,dh\le c\,\mathrm{big}(C,n,B)\,\mathrm{col}(g)^{-(2n+7)};
\]
\item for $X$ in the Lie algebra of $\Uni(J_{\eps'})$ with $\|X\|\le1$, $h\mapsto\langle d\pi(\tilde X)v,\pi(\tilde h)v\rangle$ is
integrable with integral of the absolute value at most $c\,\mathrm{big}(C,n,B)$;
\item there is $A$ with $|\langle v,\pi(g)v\rangle|\le A\|g\|^{-20(n+1)^2}$ for $g\in\Uni(J_\eps)$.
\end{enumerate}
\formal{UpperTheory.Lit.Branch.RealInputE}
\formal{GlobalTheory.SepReal.Par}
\formal{GlobalTheory.SepReal.ForPair}
\formal{GlobalTheory.SepReal.big}
\end{definition}

\subsubsection{Proof of the main theorem}

\begin{theorem}[Separation from the archimedean input]\label{pf:fs:mainRI}
$\sfname{ThetaLocalInputs}$, $\sfname{HoweUnram}$, $\sfname{SiegelWeilUnitaryConvergent}$ and $\fsP{RI}$ imply
$\sfname{FunctionalsSeparate}$.
\formallabel{UpperTheory.Lit.real_functionals_separate_of_hypothesis}
\end{theorem}

\begin{proof}\sloppy
Let $P$ denote the bundle $\sfname{ThetaLocalInputs}$. We establish the intermediate statements in the following order, each
from the ones before.
\begin{enumerate}[label=(\alph*)]
\item $\fsP{RM}$ holds by $P$ (\cref{pf:fs:models}).
\item $\fsP{PE}$ holds by $P$ (\cref{pf:fs:pairex}), and $\fsP{RS}$ follows from $P$, $\fsP{PE}$ and $\fsP{RM}$ (\cref{pf:fs:rowstep}).
\item $\fsP{RT}$ follows from $\fsP{RM}$ and $\fsP{RS}$ (\cref{pf:fs:tower}).
\item $\fsP{RC}$ and $\fsP{RR}$ follow from $P$ and $\fsP{RT}$ (\cref{pf:fs:conv}).
\item $\fsP{PN}$ follows from the fields $\sfname{HarishChandraModules}$, $\sfname{CasselmanWallach}$ and $\sfname{CuspRapidDecay}$ of
$P$ and from $\fsP{RI}$ (\cref{pf:fs:periodNe}).
\item $\fsP{CO}$ follows from $P$, its fields $\sfname{HarishChandraFinite}$, $\sfname{CuspRapidDecay}$,
$\sfname{DiscreteSeriesCover}$, $\sfname{HarishChandraModules}$, $\sfname{CasselmanWallach}$, $\sfname{WeilRepresentationUnitaryPair}$,
$\sfname{WeilSiegelFormulaUnitary}$, $\sfname{DoublingMethodLocalZetaIntegrals}$, $\sfname{UnramifiedDoublingGL}$, and from
$\sfname{HoweUnram}$ and $\sfname{SiegelWeilUnitaryConvergent}$ (\cref{pf:fs:coreOnto}).
\item $\fsP{DF}$ holds (\cref{pf:fs:depFixed}), hence $\fsP{AS}$ (\cref{pf:fs:ascent}).
\item $\fsP{SG}$ (\cref{pf:fs:sizeGap}) and $\fsP{CM}$ (\cref{pf:fs:models}, from $P$) hold.
\item $\fsP{WS}$ follows from $\fsP{SG}$, $\fsP{CM}$, $\fsP{RR}$, $\fsP{PN}$, $\fsP{CO}$ and $\fsP{AS}$ (\cref{pf:fs:weakSep}).
\item $\fsP{GR}$ holds (\cref{pf:fs:goodRows}), and $\sfname{FunctionalsSeparate}$ follows from $P$, $\fsP{RT}$, $\fsP{SG}$,
$\fsP{GR}$, $\fsP{RC}$ and $\fsP{WS}$ (\cref{pf:fs:besselSep}).\qedhere
\end{enumerate}
\end{proof}

\begin{theorem}[The archimedean input holds]\label{pf:fs:RIholds}
$\fsP{RI}$ holds.
\formallabel{UpperTheory.Lit.realInputE_holds}
\end{theorem}

This is proved in \cref{sec:pf:fs:real}.

\begin{proof}[Proof of \cref{pf:fs:main}]
Apply \cref{pf:fs:mainRI} with the archimedean input of \cref{pf:fs:RIholds}.
\end{proof}

% ---------------------------------------------------------------------------
\subsection{Weak separation}\label{sec:pf:fs:weak}

\begin{lemma}[Integrability of core periods]\label{pf:fs:perint}
\begin{enumerate}[label=(\alph*)]
\item Let $V$ be a hermitian space of $D$, $\mu$ a Haar measure on $G_V$, $\fsF$ a fundamental domain for $\Gamma_V$, $\Phi$ a
rapidly decreasing function and $f$ a measurable function with $|f(g)|\le C\,\mathrm{ht}(g)^A$. Then $f\Phi$ is integrable on $\fsF$.
\item In a setting, let $F$ be a rapidly decreasing function on $G_{V_{\fsc_0}}$ and $u$ a rapidly decreasing function on $G_{Y_0}$.
Then $h\mapsto F(\iota_0h)u(h)$ is $\nu_{(0,0)}$-integrable on $\fsF_{(0,0)}$. Consequently, if $F=\sum_jc_jF_j$ is a finite linear
combination of rapidly decreasing functions, then $\fsper_0(F,u)=\sum_jc_j\fsper_0(F_j,u)$.
\end{enumerate}
\formallabel{GlobalTheory.Grow.RapidDec.integrableOn_mul_of_le}
\formallabel{GlobalTheory.SepRows.WeakSepPf.integrableOn_corePeriod_rapidlyDecreasing}
\end{lemma}

\begin{proof}
(a) By rapid decay, $|\Phi(g)|\mathrm{ht}(g)^A\le B$ for some $B$, so $|f\Phi|\le|C|B$; the fundamental domain has finite volume
(\lref{GlobalTheory.Mem.VolFin.finVol}: the quotient of the adelic points of a unitary group by its rational points has
finite volume, cf.\ \cite{Borel1969}), so the bounded measurable function $f\Phi$ is integrable on $\fsF$.
(b) Since $M$ is a row, $\nu_{(0,0)}$ is a Haar measure and $\fsF_{(0,0)}$ a fundamental domain. The function $F\circ\iota_0$ is
continuous and bounded (rapid decay of $F$ with exponent $0$), so (a) applies with $A=0$. The second assertion follows by
linearity of the integral, each term being integrable by the first assertion.
\end{proof}

\begin{lemma}[A reachable function with non-zero period]\label{pf:fs:reachspan}
In a setting, let $v$ be a core volume, $u$ rapidly decreasing on $G_{Y_0}$ and $F$ in the span of $\fsReach(v)$ with
$\fsper_0(F,u)\ne0$. Then some $F'\in\fsReach(v)$ has $\fsper_0(F',u)\neq0$.
\formallabel{GlobalTheory.SepRows.WeakSepPf.exists_reach_of_span}
\formallabel{GlobalTheory.SepRows.WeakSepPf.rapidDec_of_reach}
\end{lemma}

\begin{proof}
Every element of $\fsReach(v)$ is the bottom $\Psi_0$ of a column chain, hence rapidly decreasing. If $\fsper_0(F',u)=0$ for all $F'\in\fsReach(v)$, then
by \cref{pf:fs:perint} $\fsper_0(\cdot,u)$ vanishes on the span, a contradiction.
\end{proof}

\begin{proposition}[Weak separation]\label{pf:fs:weakSep}
$\fsP{SG}$, $\fsP{CM}$, $\fsP{RR}$, $\fsP{PN}$, $\fsP{CO}$ and $\fsP{AS}$ imply $\fsP{WS}$.
\formallabel{GlobalTheory.SepRows.weakSep_of_steps}
\formallabel{GlobalTheory.T.rowSeparation_by_realFunctionals_of_chain_steps}
\formallabel{GlobalTheory.SepRows.WeakSepPf.nonempty_coreVol}
\formallabel{GlobalTheory.SepRows.WeakSepPf.rapidDec_of_rapidRow}
\end{proposition}

\begin{proof}
Let $\eps$ and $C_F\ge6510$ be given. Let $x_1$ be the threshold of $\fsP{PN}$ for $(\eps,C_F)$ and $x_2$ that of $\fsP{SG}$ for
$\eps$; put $x_0=\max(x_1,x_2)$. Fix $D$, size data with exponent $\eps$, scale $x\ge x_0$ and $\sfname{SizeRows}$, and
$(T,L,\mu_0,G,W,M)$ as in $\forall_{\mathrm{ev}}$, a row model $\mathcal M_0$, a non-zero $t\in\fsRows(\mathcal M_0)$ and a non-zero
$w$. By $\sfname{SizeRows}$, $D$ is off $Q$, and by $\fsP{SG}$, $\sfname{Sizes}(D)$ holds; since $c_{\mathrm{gap}}=1$, the statements
quantified by $\forall_{\mathrm{gap}}$ and $\forall_{\mathrm{all}}$ apply to the present setting.

By $\fsP{CM}$ there is a model $\mathcal M_{\fsc}$ of $V_{\fsc_0}$ with parameter $S^{\fsc}_0$ and Cartan involution $c_{\fsc_0}$. By
$\fsP{PN}$ there are a core form $y$ of $\mathcal M_{\fsc}$ and $a\in\fsKF(\mathcal M_{\fsc})$ with $\fsper_0(y(a),t(w))\ne0$. Choose a
core volume $v$ (Haar measures and fundamental domains exist on every $G_{V_{\fsc_i}}$:
\lref{GlobalTheory.Unf.Inhab.nonempty_vol} provides one for each hermitian space). By $\fsP{RR}$, $u=t(w)$ is row-rapid; in
particular $u=\Theta^{(0)}u$ is rapidly decreasing. By $\fsP{CO}$, $y(a)$ lies in the span of $\fsReach(v)$, so by \cref{pf:fs:reachspan} there is a reachable $F$ with
$\fsper_0(F,u)\ne0$: there are a column chain $\Psi$ with $\Psi_0=F$, a column form $\Psi'\in\fsCols$ and $\xi$ with
$\Psi_m=\Psi'_\xi$. By $\fsP{AS}$ there are $\rho,\Phi$ with $\fsB(\Psi_m,\Theta_{\rho,\Phi}u)\ne0$, that is
$f_{\Psi',\Phi}(\rho,\xi;t(w))=\fsB(\Psi'_\xi,\Theta_{\rho,\Phi}u)\neq0$.
\end{proof}

% ---------------------------------------------------------------------------
\subsection{From weak separation to separation}\label{sec:pf:fs:bessel}

The passage from $\fsP{WS}$ to $\sfname{FunctionalsSeparate}$ rests on the uniqueness of Bessel models at the archimedean
places. We first recall the local notions.

\begin{definition}[Bessel frames, generic characters and Bessel forms]\label{pf:fs:besselforms}
Let $n=m'+2r'+1$ with $m'\ge0$, and let $J$ ($n\times n$), $J_0$ ($m'\times m'$), $F$ ($n\times m'$), $e$ ($n\times r'$) be complex
matrices and $z\in\C^n$; write $\langle x,y\rangle_J=x^*Jy$.
\begin{enumerate}[label=(\alph*)]
\item $(J,J_0,F,e,z)$ is a \emph{Bessel frame} if $J$ is hermitian and invertible, $J_0$ is invertible, $F^*JF=J_0$, the columns of $e$
are linearly independent, $e^*Je=0$, $F^*Je=0$, $F^*Jz=0$, $e^*Jz=0$ and $\langle z,z\rangle_J\neq0$. The \emph{unipotent radical}
$N(J,e)$ is the set of $u\in\Uni(J)$ with $(u-1)e_i\in\mathrm{span}(e_1,\dots,e_{i-1})$ and $(u-1)x\in\mathrm{span}(e)$ for
$x\perp_J e$. An element $g\in\Uni(J)$ \emph{extends} $h\in\Uni(J_0)$ if $gF=Fh$ and $g$ fixes the $J$-orthogonal of the
image of $F$.
\item A function $\chi$ on $N(J,e)$ is a \emph{generic character} if there are $c_1,\dots,c_{r'}\in\C^\times$ with
$\chi(u)=\exp(i\,\Re\sum_jc_ja_{jj})$ whenever $(u-1)f_{j+1}=\sum_ka_{jk}e_k$, where $(f_1,\dots,f_{r'+1})=(e_1,\dots,e_{r'},z)$.
\item Let $\sigma$ be a representation of $\Uni(J)$ on a Hilbert space $H$, $V\subseteq H$ a subspace and $\pi_0$ a
representation of $\Uni(J_0)$ on a finite-dimensional space $W$. A \emph{Bessel form} is a bilinear map
$\beta\colon W\times V\to\C$ such that for every $w$, $\beta(w,\cdot)$ is the restriction of a linear form continuous for a
Sobolev norm of $\sigma$; $\beta(\pi_0(h)w,y)=\beta(w,x)$ whenever $y=\sigma(g)x$ and $g\in\Uni(J)$ extends $h\in\Uni(J_0)$; and
$\beta(w,\sigma(u)x)=\chi(u)\beta(w,x)$ for $u\in N(J,e)$.
\item A family $\lambda\in\R^n$ is a \emph{Harish-Chandra parameter} if it is strictly decreasing and
$\lambda_a-(n-1)/2\in\Z$; the \emph{gap} of $(\eps,\lambda)$ is at least $h$ if $|\lambda_a-\lambda_b|\ge h$ whenever
$\eps_a\ne\eps_b$. A model $\mathcal M$ of $V$ has \emph{large gap} if at every $\tau$ its parameter $\lambda_\tau$ is a
Harish-Chandra parameter and the gap of $(\eps_\tau,\lambda_\tau)$ is at least $120(\dim V)^3$.
\end{enumerate}
\formal{GlobalTheory.RealBessel.IsFrame}
\formal{GlobalTheory.RealBessel.unipRad}
\formal{GlobalTheory.RealBessel.IsEmb}
\formal{GlobalTheory.RealBessel.IsGeneric}
\formal{GlobalTheory.RealBessel.IsBesselForm}
\formal{GlobalTheory.RealBessel.besselForms}
\formal{Formalism.QuantumDeterminants.IsHCParam}
\formal{Formalism.QuantumDeterminants.GapAtLeast}
\formal{UpperTheory.Lit.Bur.HasGap}
\end{definition}

\begin{definition}[Three statements on the top of the grid]\label{pf:fs:topstatements}
\begin{enumerate}[label=(\arabic*)]
\item $\fsP{BG}$ (Bessel gap): for every Bessel frame $(J,J_0,F,e,z)$ with $m'\ge1$, $\langle z,z\rangle_J<0$ and $J_0$ positive
definite, every generic character $\chi$, every irreducible unitary representation $\sigma$ of $\Uni(J)$ on a Hilbert space which
is, after conjugation by a frame of $J$, a discrete series with a Harish-Chandra parameter whose gap is at least
$h\ge120n^3$, and every continuous irreducible finite-dimensional representation $\pi_0$ of $\Uni(J_0)$, the space of Bessel
forms on the smooth vectors of $\sigma$ has dimension at most $1$.
\item $\fsP{CT}$ (coefficient tensor): for every $D$, $V$, $\theta$ and every coefficient datum $c$ of $V$ with its standing
properties, there are finite-dimensional spaces $W_\tau$ ($\tau\in\fsT$) with continuous irreducible representations $\pi_\tau$
of $\Uni(V^\tau)$ and an isomorphism $e\colon\bigotimes_\tau W_\tau\to(W^c)^*$ such that
$e(\otimes_\tau y_\tau)=e(\otimes_\tau x_\tau)\circ\rho^c(\iota_\tau(g))^{-1}$ whenever $y_\tau=\pi_\tau(g)x_\tau$ and
$y_{\tau'}=x_{\tau'}$ for $\tau'\ne\tau$.
\item $\fsP{TS}$ (top separation), $\forall_{\mathrm{all}}$: for every model $\mathcal M_B$ of $V_{(m,0)}$, every
$F'\in\fsA(\mathcal M_B;\emptyset,1)$ and every $\Psi\in\fsCols$: (i) $(\xi,v)\mapsto\fsB(\Psi_\xi,F'(v))$ is bilinear on
$(W^+)^*\times\fsHi(\mathcal M_B)$; (ii) for every decomposition $(W_\tau,\pi_\tau,e)$ as in $\fsP{CT}$ for $c^+$, the function
$\beta(x,v)=\fsB(\Psi_{e(\otimes x)},F'(\otimes v))$ of $x\in\prod_\tau W_\tau$ and $v\in\prod_\tau\fsHi_\tau$ is
\emph{separately Bessel}: for every $\tau$ and all fixed other coordinates, it is a Bessel form at $\tau$ for the top frame of
$G$, the character $\psi_N$ at $\tau$, $\pi_\tau$ and $\mathcal M_B$ at $\tau$.
\end{enumerate}
\formal{UpperTheory.Lit.Bur.BurGapPos}
\formal{GlobalTheory.SepRows.CoefTensor}
\formal{GlobalTheory.SepRows.SepBes}
\formal{GlobalTheory.SepRows.TopSep}
\end{definition}

The three statements are proved in \cref{sec:pf:fs:bur}: $\fsP{BG}$ in \cref{pf:fs:burGapPos}, $\fsP{CT}$ in \cref{pf:fs:coefTensor}
and $\fsP{TS}$ in \cref{pf:fs:topSep}.

\begin{lemma}[Configuration at the top]\label{pf:fs:topCfg}
In a setting, at every $\tau$: the top frame of $G$ (the frame of the decomposition $V_{(m,0)}=V_{\fst_0}\oplus U$, the isotropic
family $e$ and the vector $z$, mapped by $\tau$) is a Bessel frame; $\psi_N$ at $\tau$ is a generic character; $\dim V_{\fst_0}\ge1$;
$\langle z,z\rangle=-1$; and $V_{\fst_0}^\tau$ is positive definite.
\formallabel{UpperTheory.Lit.Bur.topCfg}
\formallabel{UpperTheory.Lit.Bur.lastVertex_besselDatum_isFrame_isGeneric}
\formallabel{UpperTheory.Lit.Bur.besselDatum_middleVector_norm_eq}
\formallabel{UpperTheory.Lit.Bur.posDef_of_hasSignature_zero_negative}
\end{lemma}

\begin{proof}
The frame property at a real place follows from the orthogonal decomposition data and the isometry property of the embedding
(\lref{GlobalTheory.T.isBesselFrame_of_besselDatum_realPlace}, which states that the Bessel datum of an orthogonal
extension gives a Bessel frame at every $\tau$). For genericity take $c_j=-4\pi\,\tau(c^{\mathrm{B}}_j)$, where $c^{\mathrm B}_j\ne0$
are the coefficients of the Bessel datum: by \lref{GlobalTheory.T.besselCharacter_realPlace_unipotent_eq_trace} the value of
$\psi_N$ at the image of $u\in N(J,e)$ is $\psi$ of a trace-type expression $t$ with $\sum_jc^{\mathrm B}_ja_{jj}$ at $\tau$ and $0$
elsewhere, and by \lref{GlobalTheory.T.standardChar_trace_single_realPlace} a standard character takes the value
$\exp(i\,\Re(-4\pi\,\cdot))$ on such elements; this is the required form. The middle vector $z$ of the Bessel datum is the
middle column of the antidiagonal complement, whose Gram entry is $-1$; mapping by $\tau$ gives $\langle z,z\rangle=-1$. The
node $\fst_0=\fsc_m$ has $n_{\fsc_m}=\lfloor1/2\rfloor=0$ negative directions, so $V_{\fst_0}^\tau$ has signature $(\dim,0)$ and
is positive definite; and $\dim V_{\fsc_m}\ge1$ by the rank formula.
\end{proof}

\begin{lemma}[Rank one on a model]\label{pf:fs:burGapModel}
Assume $\fsP{BG}$. Let $\mathcal M$ be a model of $V$ and $\tau\in\fsT$ such that $\lambda_\tau$ is a Harish-Chandra parameter
with gap at least $h\ge120(\dim V)^3$. For every Bessel frame $(V^\tau,J_0,F,e,z)$ with $J_0$ positive definite, $m'\ge1$ and
$\langle z,z\rangle<0$, every generic $\chi$ and every continuous irreducible finite-dimensional $\pi_0$ of $\Uni(J_0)$, the space
of Bessel forms for $\pi_\tau$ on $\fsHi_\tau$ has dimension at most $1$.
\formallabel{UpperTheory.Lit.Bur.burGap_model}
\end{lemma}

\begin{proof}
$\pi_\tau(g)=\sigma_\tau(\mathrm{fr}_\tau g\,\mathrm{fr}_\tau^{-1})$ is unitary and irreducible on $\Uni(V^\tau)$ because $\sigma_\tau$ is
(both properties transport along conjugation by $\mathrm{fr}_\tau$, \lref{GlobalTheory.SmVec.isURepOn_trRep} and
\lref{GlobalTheory.SmVec.isIrredOn_trRep}); conjugating back by $\mathrm{fr}_\tau$ shows that $\pi_\tau$ is a discrete series with
gap at least $h$ in the sense of $\fsP{BG}$; and $\fsHi_\tau$ is the space of smooth vectors of $\pi_\tau$
(\lref{GlobalTheory.SmVec.sm_eq}). Apply $\fsP{BG}$.
\end{proof}

\begin{lemma}[Products of separately proportional functions]\label{pf:fs:separately}
Let $(X_\tau)_{\tau\in\fsT}$ be sets, $b_\tau\colon X_\tau\to\C$ and $\beta\colon\prod_\tau X_\tau\to\C$. Suppose that for every $\tau$
and every $x$ there is $c$ with $\beta(x_{\tau\leftarrow y})=c\,b_\tau(y)$ for all $y\in X_\tau$ ($x_{\tau\leftarrow y}$ is $x$
with the $\tau$-coordinate replaced by $y$). Then there is $l$ with $\beta(x)=l\prod_\tau b_\tau(x_\tau)$ for all $x$.
\formallabel{GlobalTheory.SepRows.mul_prod_of_separately}
\end{lemma}

\begin{proof}
If some $b_\tau$ vanishes identically, then so does $\beta$ (apply the hypothesis at $\tau$), and $l=0$ works. Otherwise choose
$y_\tau$ with $b_\tau(y_\tau)\ne0$ and change one coordinate at a time: by the hypothesis at $\tau$,
$\beta(x_{\tau\leftarrow x'_\tau})\,b_\tau(x_\tau)=\beta(x)\,b_\tau(x'_\tau)$ for all $x,x'_\tau$ (both sides equal
$c\,b_\tau(x'_\tau)b_\tau(x_\tau)$). Moving from $(y_\tau)_\tau$ to $x$ in $|\fsT|$ steps gives
$\beta(x)\prod_\tau b_\tau(y_\tau)=\beta(y)\prod_\tau b_\tau(x_\tau)$; put $l=\beta(y)/\prod_\tau b_\tau(y_\tau)$.
\end{proof}

\begin{lemma}[Separately Bessel functions are proportional]\label{pf:fs:sepLine}
Assume $\fsP{BG}$. Let $\mathcal M$ be a model of $V$ with large gap, $A$ a hermitian space and, at every $\tau$, a Bessel frame
$(V^\tau,A^\tau,F_\tau,e_\tau,z_\tau)$ with $\dim A\ge1$, $\langle z_\tau,z_\tau\rangle<0$ and $A^\tau$ positive definite, a
generic character $\chi_\tau$ and a continuous irreducible finite-dimensional representation $\pi_\tau$ of $\Uni(A^\tau)$ on
$W_\tau$. Then there is a function $\ell$ on $\prod_\tau W_\tau\times\prod_\tau\fsHi_\tau$ such that every separately Bessel function
$\beta$ satisfies $\beta=l\,\ell$ for some constant $l$.
\formallabel{UpperTheory.Lit.Bur.sepLine_gap}
\end{lemma}

\begin{proof}
By \cref{pf:fs:burGapModel}, at every $\tau$ the space of Bessel forms has dimension at most $1$; choose $b_\tau$ spanning it
(possibly $0$), so that every Bessel form at $\tau$ is a multiple of $b_\tau$. Put $\ell(x,v)=\prod_\tau b_\tau(x_\tau,v_\tau)$.
If $\beta$ is separately Bessel, then for every $\tau$ and every point, the function of the $\tau$-coordinates
$(x_\tau,v_\tau)$ is a Bessel form, hence a multiple of $b_\tau$. \Cref{pf:fs:separately}, applied to $X_\tau=W_\tau\times\fsHi_\tau$,
gives $l$ with $\beta=l\,\ell$.
\end{proof}

\begin{proposition}[The top line]\label{pf:fs:topLine}
Assume $\fsP{BG}$, $\fsP{CT}$ and $\fsP{TS}$. In every setting, for every model $\mathcal M_B$ of $V_{(m,0)}$ with large gap there is
a function $\ell(\xi,v)$ such that for every $F'\in\fsA(\mathcal M_B;\emptyset,1)$ and every $\Psi\in\fsCols$ there is a constant $l$
with
\[
\fsB\bigl(\Psi_\xi,F'(v)\bigr)=\ell(\xi,v)\,l\qquad\text{for all }\xi\in(W^+)^*,\ v\in\fsHi(\mathcal M_B).
\]
\formallabel{UpperTheory.Lit.Bur.topLine_gap}
\formallabel{GlobalTheory.TopLineProof.bilin_prop}
\formallabel{GlobalTheory.TopLineProof.coefPlus_is}
\end{proposition}

\begin{proof}
The top coefficient datum has its standing properties (a short verification, recorded in the formal proof), so $\fsP{CT}$ gives
$(W_\tau,\pi_\tau,e)$ for $c^+$. By \cref{pf:fs:topCfg}, the hypotheses of \cref{pf:fs:sepLine} hold for the top frame, the characters
$\psi_N$ at $\tau$, the representations $\pi_\tau$ and $\mathcal M_B$; let $\ell_0$ be the function it provides. Pure tensors
$e(\otimes x)$ span $(W^+)^*$ (as $e$ is an isomorphism) and pure tensors $\otimes v$ span $\fsHi(\mathcal M_B)$.

If $\fsB(\Psi_\xi,F'(v))=0$ for all $F',\Psi,\xi,v$, then $\ell=0$ works. Otherwise fix $F_0,\Psi_0,\xi_0,v_0$ with
$\fsB(\Psi_{0,\xi_0},F_0(v_0))\ne0$ and put $\ell(\xi,v)=\fsB(\Psi_{0,\xi},F_0(v))$. Given $F',\Psi$, let $b$ and $b_0$ be the bilinear
forms $(\xi,v)\mapsto\fsB(\Psi_\xi,F'(v))$ and $(\xi,v)\mapsto\fsB(\Psi_{0,\xi},F_0(v))$ of $\fsP{TS}$(i). By $\fsP{TS}$(ii) and
\cref{pf:fs:sepLine} there are constants $l,l_0$ with $b(e(\otimes x),\otimes v)=\ell_0(x,v)\,l$ and
$b_0(e(\otimes x),\otimes v)=\ell_0(x,v)\,l_0$ for all $x,v$. Since $b_0\ne0$ and pure tensors span, $\ell_0\ne0$ and $l_0\ne0$. Hence
$b=(l/l_0)\,b_0$ on pairs of pure tensors, and by bilinearity and spanning, $b=(l/l_0)b_0$ (an elementary step of linear algebra). Thus $\fsB(\Psi_\xi,F'(v))=\ell(\xi,v)\cdot l/l_0$.
\end{proof}

\begin{proposition}[Factorisation of the column--row pairing]\label{pf:fs:factor}
Assume $\fsP{BG}$, $\fsP{CT}$, $\fsP{TS}$ and $\fsP{RT}$. Let $D$ carry size data with $\sfname{SizeRows}$ and $\sfname{Sizes}(D)$,
$(T,L,\mu_0,\mathbf c,G,W,M)$ a setting with $c_{\mathrm{gap}}=1$, $\mathcal M_0$ a row model and $w\in\fsHi(\mathcal M_0)$. There is a
function $a(\rho,\xi)$ such that for every $t\in\fsRows(\mathcal M_0)$, every $\Psi\in\fsCols$ and every coset datum $\Phi$ there is
a constant $l$ with
\[
f_{\Psi,\Phi}(\rho,\xi;t(w))=a(\rho,\xi)\,l\qquad\text{for all }\rho,\xi .
\]
\formallabel{UpperTheory.Lit.Bur.factor_gap}
\formallabel{GlobalTheory.SepRows.factor_multT_le_empty}
\end{proposition}

\begin{proof}
By $\fsP{RT}$ with $n=m$ there are a model $\mathcal M_B$ of $V_{(m,0)}$ with parameter $S^{\mathrm{row}}_m$ and a map
$\mathrm{vec}_m$. By \cref{pf:fs:gap} below (which uses $\sfname{SizeRows}$ and $c_{\mathrm{gap}}=1$), $\mathcal M_B$ has large gap, so
\cref{pf:fs:topLine} gives $\ell$. Put $a(\rho,\xi)=\ell(\xi,\mathrm{vec}_m(\rho,w))$. Let $t,\Psi,\Phi$ be given. Since
$K_{\mathrm{row}}$-invariance with trivial character is a stronger condition than none,
$\fsRows(\mathcal M_0)\subseteq\fsA(\mathcal M_0;\emptyset,1)$; so $\fsP{RT}$ gives $F'\in\fsA(\mathcal M_B;\emptyset,1)$ with
$\Theta_{\rho,\Phi}t(w)=\Theta^{(m)}_{\rho,\Phi}t(w)=F'(\mathrm{vec}_m(\rho,w))$ for all $\rho$. Let $l$ be the constant of
\cref{pf:fs:topLine} for $(F',\Psi)$. Then
$f_{\Psi,\Phi}(\rho,\xi;t(w))=\fsB(\Psi_\xi,F'(\mathrm{vec}_m(\rho,w)))=\ell(\xi,\mathrm{vec}_m(\rho,w))\,l=a(\rho,\xi)\,l$.
\end{proof}

\begin{proposition}[Separation]\label{pf:fs:besselSep}
$\fsP{BG}$, $\fsP{CT}$, $\fsP{TS}$, $\fsP{RT}$, $\fsP{SG}$, $\fsP{GR}$, $\fsP{RC}$ and $\fsP{WS}$ imply
$\sfname{FunctionalsSeparate}$. In particular, $\sfname{ThetaLocalInputs}$, $\fsP{RT}$, $\fsP{SG}$, $\fsP{GR}$, $\fsP{RC}$ and
$\fsP{WS}$ imply $\sfname{FunctionalsSeparate}$.
\formallabel{UpperTheory.Lit.Bur.separate_of_burGapPos}
\formallabel{UpperTheory.Lit.bessel_separate_of_step_gap}
\end{proposition}

\begin{proof}
The second assertion follows from the first, since $\fsP{BG}$, $\fsP{CT}$ and $\fsP{TS}$ hold (the last one given
$\sfname{ThetaLocalInputs}$), by \cref{pf:fs:burGapPos,pf:fs:coefTensor,pf:fs:topSep}.

Let $\eps$ and $C_F\ge6510$ be given; let $x_1$ be the threshold of $\fsP{SG}$ for $\eps$ and $x_2$ that of $\fsP{WS}$ for
$(\eps,C_F)$, and put $x_0=\max(x_1,x_2)$. Fix $D$, size data with $x\ge x_0$ and $\sfname{SizeRows}$, the remaining data
$\Gamma,\mathrm{Pl}_0,T,\mathcal A,K_s,R$ with $L_R.\mathrm{Is}(\mu_0)$, a grid datum $G$ at $(1,C_F,b)$, a Weil datum $W$ over $R$, a row
$M$ and a non-zero $w\in\fsHi(\mathcal M_Y(0))$. (The hypotheses $\sfname{ESerTie}(D)$ and the group record of $R$ are not used.)
Then $\sfname{Sizes}(D)$ holds by $\fsP{SG}$, $D$ is off $Q$ by \cref{pf:fs:offQ}, and $m\ge44>0$ by $\sfname{SizeRows}$. Let
$\mathcal M_0=\mathcal M_Y(0)$; since $W$ is over $R$, its Cartan involution at $(0,0)$ is that of $\mathcal M_0$, so $\mathcal M_0$ is a
row model for the setting. By $\fsP{GR}$ (and $\theta=\theta_R$, as $W$ is over $R$), every $t\in\fsE_R$ defines a row
$\tilde t\in\fsRows(\mathcal M_0)$ with $\tilde t(w)=t(w)$, and $\tilde t\ne0$ if $t\ne0$ (multilinear maps correspond injectively to
linear maps on the tensor product). Let $a$ be the function of \cref{pf:fs:factor} for $w$.

If $\fsE_R=0$, any $\rho$ and $\xi=0$ separate at $w$. Otherwise fix $t_0\ne0$ in $\fsE_R$. By $\fsP{WS}$ there are
$\rho_0,\xi_0,\Psi_0,\Phi_0$ with $f_{\Psi_0,\Phi_0}(\rho_0,\xi_0;t_0(w))\ne0$; writing this value as $a(\rho_0,\xi_0)\,l_0$
(\cref{pf:fs:factor}) shows $a(\rho_0,\xi_0)\ne0$. We claim that $(\rho_0,\xi_0)$ separates at $w$. Let $t\ne0$ in $\fsE_R$. By
$\fsP{WS}$ there are $\rho,\xi,\Psi,\Phi$ with $f_{\Psi,\Phi}(\rho,\xi;t(w))\ne0$; with the constant $l$ of \cref{pf:fs:factor} for
$(t,\Psi,\Phi)$ this value is $a(\rho,\xi)\,l$, so $l\ne0$. Hence
$f_{\Psi,\Phi}(\rho_0,\xi_0;t(w))=a(\rho_0,\xi_0)\,l\neq0$, and $\mathrm{FConv}(\rho_0,\xi_0,\Psi,\Phi;t(w))$ holds by $\fsP{RC}$.
\end{proof}

% ---------------------------------------------------------------------------
\subsection{The row tower}\label{sec:pf:fs:tower}

\subsubsection{Models}

\begin{lemma}[Models with a prescribed Cartan involution]\label{pf:fs:modelenum}
Assume $\sfname{DiscreteSeriesUpqExists}$ (existence of discrete series of $\Uni(p,q)$ for regular integral parameters, a field of
$\sfname{ThetaLocalInputs}$). Let $V$ be a hermitian space of $D$ with consistent signature data, $S$ a family of parameters and,
at every $\tau$, matrices $O_\tau$, $C_\tau$ with $O_\tau^*V^\tau O_\tau=J_{\eps}$ and $C_\tau O_\tau=O_\tau J_\eps$ for some
$\eps\in\{\pm1\}^{\dim V}$. Suppose that at every $\tau$ the set $S_\tau$ has an enumeration $(\lambda,\eps')$ with $\lambda$ a
Harish-Chandra parameter and with as many negative signs as $V^\tau$ has negative directions. Then there is a discrete series
model of $V$ with parameter $S$ and Cartan involution $C$.
\formallabel{GlobalTheory.SepRows.exists_model_of_enum}
\formallabel{GlobalTheory.SepRows.exists_dsModel_cart}
\end{lemma}

\begin{proof}
At every $\tau$ the sign vectors $\eps$ (of the frame) and $\eps'$ (of the enumeration) have the same number of negative entries
(both equal the number of negative directions of $V^\tau$), so there is a permutation $s$ with $\eps\circ s=\eps'$
(\lref{GlobalTheory.GridJoin.exists_perm_of_card_false}); permuting the columns of $O_\tau$ by $s$ gives a frame $O'_\tau$ with
$O'^*_\tau V^\tau O'_\tau=J_{\eps'}$ and $C_\tau O'_\tau=O'_\tau J_{\eps'}$. Given such orthonormal frames and enumerations by
Harish-Chandra parameters, $\sfname{DiscreteSeriesUpqExists}$ provides the discrete series $\sigma_\tau$ and hence the model with
Cartan involution $C$ (\lref{GlobalTheory.T.exists_discreteSeriesModel_of_orthonormalBasis_cartan}). The second formal name
is the variant in which the enumeration is obtained from the regularity of $S_\tau$, its cardinality $\dim V$, its number of
negative signs, and the fact that its enumerations are Harish-Chandra parameters.
\end{proof}

\begin{proposition}[Row and core models]\label{pf:fs:models}
$\sfname{ThetaLocalInputs}$ implies $\fsP{RM}$ and $\fsP{CM}$.
\formallabel{GlobalTheory.SepRows.separating_models_of_separating_data}
\formallabel{GlobalTheory.T.exists_models_of_separating_data}
\formallabel{GlobalTheory.SepRows.coreModels_of_upq}
\formallabel{GlobalTheory.SepRows.separating_family_has_models}
\formallabel{GlobalTheory.T.exists_separating_models}
\formallabel{GlobalTheory.SepRows.isHCParam_dual_parV}
\formallabel{GlobalTheory.SepRows.isHCParam_parPi}
\end{proposition}

\begin{proof}
In both cases we apply \cref{pf:fs:modelenum} with the frames and Cartan involutions of $W$ at the node (they satisfy its
hypothesis by the definition of a Weil datum), so it suffices to produce the enumerations.

$\fsP{RM}$, node $(0,0)$: here $S^{\mathrm{row}}_0=S^Y$ (\cref{pf:fs:offQ}) and $Y_0$ has consistent signature data; the twisted
parameter $S^Y_\tau$ has an enumeration by a Harish-Chandra parameter whose number of negative signs is the number of negative
directions of $Y_0^\tau$ (\lref{GlobalTheory.ChainModels.goodPar_parB}). Node $(i,0)$, $1\le i\le m$: the grid parameter at
$(i,0)$ is regular (\lref{GlobalTheory.GridParDCount.isRegular_parD}); the dual of its shift by $-\frac12$ or $0$ (according to
the parity of $i$) is regular, has the same cardinality $\dim V_{(i,0)}$ (\lref{GlobalTheory.GridParDCount.card_parD}) and
$n_{(i,0)}$ negative signs (\lref{GlobalTheory.GridJoin.dualPar_shift_counts} and \lref{GlobalTheory.SignedHC.Model.nMinus_parD});
hence it has an enumeration (\lref{GlobalTheory.GridJoin.exists_isEnum_of_regular}), and every enumeration of the dual grid
parameter is a Harish-Chandra parameter: its entries are strictly decreasing by definition and are congruent to
$(\dim V_{(i,0)}-1)/2$ modulo $\Z$ by inspection of the entries of the grid parameter.

$\fsP{CM}$, node $\fsc_i$: the core parameter is regular, has cardinality $\dim V_{\fsc_i}$ and $n_{\fsc_i}$ negative signs
(\lref{GlobalTheory.PiCounts.isRegular_parPi}, \lref{GlobalTheory.PiCounts.card_parPi}, and \lref{GlobalTheory.PiCounts.nMinus_parPi}),
and its enumerations are Harish-Chandra parameters (the integrality condition is checked as above).
\end{proof}

\subsubsection{Seesaw pairs and the row step}

\begin{definition}[Pair conditions, vacuum vectors and intertwiners]\label{pf:fs:pairdefs}
For families of parameters $(S,S')$ and exponents $(a_S,a_T)$, the \emph{pair condition} at $\tau$ says that $S_\tau,S'_\tau$ are the
duals of two parameter sets $P,P'$ which are in theta correspondence for a segment of some length $t'$ with the exponents
$(a_T,a_S)$, in one of the two orders, with a non-proximity condition in the first case. A Weil datum $\omega$ for $(V,W')$ has
\emph{exponents} $(a_S,a_T)$ if its splitting characters have infinity types $a_S(\tau),a_T(\tau)$. A \emph{vacuum vector} of
$(\mathcal M,\mathcal M')$ is a non-zero theta function which at every $\tau$ is a Fock vacuum for the complex structure defined by the
Cartan involutions of $\mathcal M$ and $\mathcal M'$; the \emph{Fock vectors} over it are $\mathrm{pt}_f(P)$, $P$ a polynomial. An
\emph{intertwiner} for $(\omega,\mathrm{pt}_f)$ is a non-zero bilinear map $I(P,v)$ from polynomials and $\fsHi(\mathcal M')$ to
$\fsHi(\mathcal M)$, bounded on pure tensors, such that $I$ is equivariant for $\Uni(V^\tau)$ and its Lie algebra (when
$\mathrm{pt}_f(P')$ is the translate of $\mathrm{pt}_f(P)$ by $g\in\Uni(V^\tau)$, $I(P',v)$ is the translate of $I(P,v)$ by
$\pi_\tau(g)$) and invariant for $\Uni(W'^\tau)$ and its Lie algebra. The statement $\fsP{PE}$ (\cref{pf:fs:statements}) asserts the
existence of an intertwiner for all $(\omega,\mathrm{pt}_f)$ simultaneously.
\formal{GlobalTheory.PairLift.PairPar}
\formal{GlobalTheory.PairLift.IsWeilExp}
\formal{GlobalTheory.Chain.IsVacVec}
\formal{GlobalTheory.Chain.ptOf}
\formal{Formalism.IndefR.Links.IsItW}
\end{definition}

\begin{proposition}[Seesaw pairs]\label{pf:fs:pairex}
$\sfname{ThetaLocalInputs}$ implies $\fsP{PE}$.
\formallabel{GlobalTheory.PairIt.exists_seesaw_pair}
\formallabel{GlobalTheory.T.exists_seesaw_pair_of_theta_data}
\formallabel{GlobalTheory.PairIt.exists_locIt_pair}
\end{proposition}

\begin{proof}
From the fields $\sfname{DiscreteSeriesCover}$ and $\sfname{HarishChandraModules}$ one obtains a record $X$ of archimedean theta data
(the chosen discrete series and the first two archimedean Rallis statements; \lref{GlobalTheory.Chain.ItPairObj.exists_record})
and the identification of the theta models with the discrete series of the cover
(\lref{GlobalTheory.T.thetaModels_eq_coverDiscreteSeries}). At every $\tau$, the pair condition gives positive ranks and a local
intertwiner $A_\tau$ between the Fock model at $\tau$ and $(\pi_\tau,\pi'_\tau)$ (writing $S_\tau,S'_\tau$ as duals of parameters in theta correspondence, the local intertwiner is provided by the record $X$). The
local intertwiners are glued into a global bilinear map $I$ by contraction over the archimedean indices
(\lref{GlobalTheory.T.exists_globalFockIntertwiner_tensor_realPlaces}). Given $\omega$ with exponents $(a_S,a_T)$ and a vacuum
vector $f$, the vacuum has the archimedean characters prescribed by the exponents
(\lref{GlobalTheory.T.thetaVacuumCharacters_atRealPlaces}) and the Fock vectors over $f$ are an isometric image of the polynomial
Fock space (\lref{GlobalTheory.T.fockVectorsOverVacuumIsometry_of_skeletonDatum}); these two facts show that $I$ is an intertwiner
for $(\omega,\mathrm{pt}_f)$ (\lref{GlobalTheory.Chain.RealIt.isItW_of_globIt}).
\end{proof}

The lift of multiplicity maps by Fock vectors is controlled by the following statement, proved elsewhere.

\begin{lemma}[Lifts of multiplicity maps along pairs; \lref{GlobalTheory.PairLift.RuleLift.pairLifts}]\label{pf:fs:pairLifts}
Let $X$ be a record of archimedean theta data with the lower-member properties. Let $(S,S',a_S,a_T)$ satisfy the rule condition of
$X$ at every $\tau$, $\mathcal M,\mathcal M'$ models of $V,W'$, $\omega$ a Weil datum with exponents $(a_S,a_T)$, $f_0$ a vacuum vector
and $I$ an intertwiner for $(\omega,\mathrm{pt}_{f_0})$. For every Haar measure $\nu$ and fundamental domain $\fsF$ on $G_{W'}$, every
Fock vacuum $f$ for the frames of the models and every $F\in\fsA(\mathcal M';\emptyset,1)$ there is $F'\in\fsA(\mathcal M;\emptyset,1)$
such that for all $P$, $v$ and $h$: $g\mapsto\theta_\omega(g,h;\mathrm{pt}_f(P))F(v)(g)$ is integrable on $\fsF$ and its integral is
$F'(I(P,v))(h)$.
\end{lemma}

\begin{lemma}[Lifts towards the negative space]\label{pf:fs:liftsneg}
Under the hypotheses of \cref{pf:fs:pairLifts} for the pair $(-B,A)$ formed with models $\mathcal M_A$ of $A$ and $\mathcal M_B$ of $B$,
frames $F_A,F_B$ with the Cartan involutions of the models, the pair condition for the flipped parameters, a Weil datum
$\omega$ with the exponents, and a non-zero Fock vacuum $f_0$, and given $\sfname{PairEx}$: there is a map $I$ such that for every
Fock vacuum $f$ the lift from $G_A$ to $G_B$ by $\theta_\omega(\cdot;\mathrm{pt}_f(P))$ carries every $F\in\fsA(\mathcal M_A;\emptyset,1)$
to some $F'\in\fsA(\mathcal M_B;\emptyset,1)$ evaluated at $I(P,v)$.
\formallabel{GlobalTheory.RowArrow.lifts_neg}
\end{lemma}

\begin{proof}
Transport $\mathcal M_B$ to a model $N$ of $-B$ with the opposite Cartan involution, together with an isomorphism $\kappa$ of
multiplicity spaces (\lref{GlobalTheory.negModel_tr}); a change of frames relates the Fock vectors for the two frame systems by a
linear map $j$ on polynomials (\lref{GlobalTheory.FrameChange.ptF_frame_change}). $\fsP{PE}$ gives an intertwiner for
$(N,\mathcal M_A)$; \cref{pf:fs:pairLifts} (with the rule condition derived from the pair condition and the third to fifth archimedean
Rallis statements, \lref{GlobalTheory.StepRule.rulePar_of_pairPar}) lifts $F$ to a multiplicity map of $N$, which $\kappa$
carries to $\fsA(\mathcal M_B;\emptyset,1)$. Put $I(P,v)=\kappa(\mathrm{It}(j(P),v))$.
\end{proof}

\begin{proposition}[The row step]\label{pf:fs:rowstep}
$\sfname{ThetaLocalInputs}$, $\fsP{PE}$ and $\fsP{RM}$ imply $\fsP{RS}$.
\formallabel{GlobalTheory.RowArrow.separatingFamily_extends_by_twisted_row}
\formallabel{GlobalTheory.T.separating_family_extend_by_twist}
\end{proposition}

\begin{proof}
Fix a setting with $\sfname{Sizes}(D)$ and $c_{\mathrm{gap}}=1$, $i<m$ and models $\mathcal M_A,\mathcal M_B$ as in $\fsP{RS}$. From
$\sfname{DiscreteSeriesCover}$, $\sfname{HarishChandraModules}$, $\sfname{CuspRapidDecay}$ and $\sfname{CasselmanWallach}$ one obtains
a record $X$ with the archimedean Rallis statements and the lower-member properties needed in \cref{pf:fs:pairLifts}
(\lref{GlobalTheory.RealPlaceTheta.exists_record_theta_all}, \lref{GlobalTheory.PairLift.inputs_of_chosenSeries}, and \lref{GlobalTheory.LowerMembersRule.lowFock}).
Choose frames $F_A,F_B$ adapted to the Cartan involutions and orthonormal frames of $W$ at $(i,0)$ and $(i+1,0)$. The parameters
$S^{\mathrm{row}}_i,S^{\mathrm{row}}_{i+1}$ satisfy the pair condition (after flipping) for the exponents attached to $\chi_{m-i}$ and
the ranks (\lref{GlobalTheory.GridParD.pairPar_X}, which uses $\sfname{Sizes}(D)$); the Weil representation $\omega_{\chi(i,0)}$ has
these exponents, because $\chi_{m-i}$ has the prescribed infinity type and $\eta$ has infinity type $\frac12$. The vacuum vector
$\phi_i^0$ is a non-zero Fock vacuum for $(F_B,F_A)$, and every coset vector $\phi_i(\kappa)$ is a Fock vacuum too
(\lref{GlobalTheory.DatumFock.isFockVacAt_of_isThetaAt}: a theta function obtained from a Fock vacuum by a local model at
non-archimedean places is again a Fock vacuum). Since $M$ is a row, $\nu_{(i,0)},\fsF_{(i,0)}$ are a Haar measure and a fundamental
domain. \Cref{pf:fs:liftsneg} with $f_0=\phi_i^0$ gives $I$; for $F\in\fsA(\mathcal M_A;\emptyset,1)$ and $\kappa$, applying it to the
Fock vacuum $f=\phi_i(\kappa)$ gives $F'$, and the Fock vectors over $\phi_i(\kappa)$ for the frames $(F_B,F_A)$ are the vectors
$\mathrm{fock}_{\chi(i,0)}(\phi_i(\kappa),P)$ of $W$. This is $\fsP{RS}$.
\end{proof}

\subsubsection{The tower and its consequences}

\begin{proposition}[The row tower]\label{pf:fs:tower}
$\fsP{RM}$ and $\fsP{RS}$ imply $\fsP{RT}$.
\formallabel{GlobalTheory.SepRows.rowTower_of_step}
\formallabel{GlobalTheory.T.firstRowTower_multiplicitySpaces_of_arrow}
\end{proposition}

\begin{proof}
Fix a setting with $\sfname{Sizes}(D)$ and $c_{\mathrm{gap}}=1$ and a row model $\mathcal M_0$. We prove the assertion of $\fsP{RT}$ by
induction on $n\le m$. For $n=0$ take $\mathcal M_0$ itself (its parameter is $S^{\mathrm{row}}_0$ by \cref{pf:fs:offQ}),
$\mathrm{vec}_0(\rho,w)=w$ and $F'=t$: there is nothing to converge and $\Theta^{(0)}t(w)=t(w)$. Suppose the assertion holds for $n<m$,
with $\mathcal M_n$ and $\mathrm{vec}_n$. By $\fsP{RM}$ there is a model $\mathcal M_{n+1}$ of $V_{(n+1,0)}$ with parameter
$S^{\mathrm{row}}_{n+1}$ and Cartan involution $c_{(n+1,0)}$, and $\fsP{RS}$ for $(\mathcal M_n,\mathcal M_{n+1})$ gives a map $I$. Put
\[
\mathrm{vec}_{n+1}(\rho,w)=I\bigl(\rho_n,\mathrm{vec}_n(\rho,w)\bigr).
\]
Given $t$ and $\Phi$, the induction hypothesis gives $F$ with $\Theta^{(n)}_{\rho,\Phi}t(w)=F(\mathrm{vec}_n(\rho,w))$ for all $\rho,w$, and
$\fsP{RS}$ applied to $F$ and $\kappa=\Phi_n$ gives $F'$ such that the lift of $F(\mathrm{vec}_n(\rho,w))$ by
$\mathrm{fock}(\phi_n(\Phi_n),\rho_n)=\phi_n(\rho,\Phi)$ converges and equals $F'(\mathrm{vec}_{n+1}(\rho,w))$. As
$\Theta^{(n+1)}_{\rho,\Phi}t(w)$ is that lift, the assertion holds for $n+1$ (the first $n$ lifts converge by induction).
\end{proof}

\begin{proposition}[Convergence and rapid decay]\label{pf:fs:conv}
$\sfname{ThetaLocalInputs}$ and $\fsP{RT}$ imply $\fsP{RC}$ and $\fsP{RR}$.
\formallabel{GlobalTheory.RowConvAll.separating_family_converges}
\formallabel{GlobalTheory.T.convergence_of_separating_family}
\formallabel{GlobalTheory.SepRows.RowRapidPf.separating_family_rapid_decay}
\formallabel{GlobalTheory.T.rapid_decay_of_separating_family}
\formallabel{GlobalTheory.RowConvAll.multT_le_empty}
\formallabel{GlobalTheory.SepRows.RowRapidPf.multT_le_empty}
\end{proposition}

\begin{proof}
Fix a setting with $\sfname{Sizes}(D)$ and $c_{\mathrm{gap}}=1$, a row model $\mathcal M_0$, $t\in\fsRows(\mathcal M_0)$ and $w$. As
$\fsRows(\mathcal M_0)\subseteq\fsA(\mathcal M_0;\emptyset,1)$ (a multiplicity map at level $(K_{\mathrm{row}},1)$ is one at level
$(\emptyset,1)$), $\fsP{RT}$ applies to $t$. For $n\le m$ and $\Phi$ it gives $F'$ with
$\Theta^{(n)}_{\rho,\Phi}t(w)=F'(\mathrm{vec}_n(\rho,w))$, the first $n$ lifts converging. By
\lref{GlobalTheory.ChainOnto.MultRapid.multT_apply_rapidDec} (using $\sfname{CuspRapidDecay}$) this value is a smooth automorphic
form of rapid decay. With $n=m$ this gives the convergence of the row lift and $\fsP{RR}$. For $\fsP{RC}$ it remains to see that the
Bessel pairing $\fsB(\Psi_\xi,F'(\mathrm{vec}_m(\rho,w)))$ converges: by \lref{GlobalTheory.RowBplus.bplusConv_col}, the Bessel
pairing of the scalar form of any column form with any automorphic function on $G_{V_{(m,0)}}$ converges when $M$ is a row (the
proof uses that $V_{\fst_0}$ is anisotropic, so that $\Gamma\backslash G_{V_{\fst_0}}$ is compact, and the normalisation of
$\nu_N$).
\end{proof}

\begin{proposition}[Sizes and good rows]\label{pf:fs:sizeGap}
\begin{enumerate}[label=(\alph*)]
\item $\fsP{SG}$ holds, with threshold $x_0=0$ for every $\eps$.
\item $\fsP{GR}$ holds.
\end{enumerate}
\formallabel{GlobalTheory.SepRows.sizeGap_of_sizes}
\formallabel{GlobalTheory.T.largeScaleSizes_imply_gapConditions}
\formallabel{GlobalTheory.SepRows.goodRows}
\formallabel{GlobalTheory.T.goodRows_principalLevel_are_rows}
\formallabel{GlobalTheory.SepRows.Good.levRow_sub}
\formallabel{GlobalTheory.SepRows.Good.principalRowLevelCharacter_trivial}
\label{pf:fs:goodRows}
\end{proposition}

\begin{proof}
(a) $\sfname{SizeRows}$ implies $\sfname{Sizes}(D)$ (\lref{GlobalTheory.RowSizes.sizes_of_sizeRows}: the inequality $1800<r$
follows from $44\le m$, $m^4\le f\le k$ and $r=k+m-1$; the gap inequality follows from the gap conditions of the parameter choice
contained in $\sfname{SizeRows}$).

(b) Let $R$ be a real datum with $L_R.\mathrm{Is}(\mu_0)$ and $t\in\fsE_R$. Since $\fsE_R$ is contained in the span of the multiplicity
maps of $\mathcal M_Y(0)$ at level $(K_R,\chi_R)$, where $K_R$ and $\chi_R$ are the level group and the level character of $R$ for the
trivial level datum, it suffices to see that every multiplicity map $T$ at level $(K_R,\chi_R)$ is one at level $(K_{\mathrm{row}},1)$.
The conditions of automorphy, cuspidality, equivariance and continuity are the same. The row level is contained in $K_R$: the
lattice $\Lambda_Y(0)$ is self-dual for the frame of $R$ at level $g_0g_1$ ($L_R.\mathrm{Is}(\mu_0)$), so the principal congruence
subgroup modulo $g_0g_1$ lies in the level group. And $\chi_R$ is identically $1$
for the trivial level datum, the extension of the trivial character being trivial. Hence $T(v)(gk)=\chi_R(k)T(v)(g)=T(v)(g)$ for
$k\in K_{\mathrm{row}}$.
\end{proof}

% ---------------------------------------------------------------------------
\subsection{Periods of core forms}\label{sec:pf:fs:periods}

This subsection proves $\fsP{PN}$. Fix a setting and a model $\mathcal M_{\fsc}$ of $V_{\fsc_0}=Y_0\oplus\ell$. The core forms
of $\mathcal M_{\fsc}$ will be Poincar\'e series of matrix coefficients of the discrete series of $\mathcal M_{\fsc}$, and the period
of such a series against a cusp form $u$ on $G_{Y_0}$ will be shown to be non-zero for large level, by isolating the identity
term. Throughout, $z_\ell$ denotes the basis vector of the line $\ell$ given by the decomposition, and for a lattice $\Lambda$ of a
space $V$ and $x\in G_V$, $\Lambda_x$ is the lattice of the vectors $y$ with $x_f^{-1}y$ in the closure of $\Lambda$ (its completions
are the translates of those of $\Lambda$ by the finite part of $x$).

\begin{definition}[Coefficients, Poincar\'e series and packs]\label{pf:fs:pack}
Let $\mathcal M$ be a model of $V$, $K$ a congruence subgroup of $G_{V,f}$ and $v,a\in\prod_\tau H_\tau$.
\begin{enumerate}[label=(\alph*)]
\item The \emph{coefficient} is $\mathrm{coef}_K(v,a)(g)=\prod_\tau\langle v_\tau,\pi_\tau(g_\tau)a_\tau\rangle$ if $g_f\in K$ and $0$
otherwise; the \emph{Poincar\'e series} with base point $c\in G_V$ is
$P_K(v,c,a)(g)=\sum_{\gamma\in\Gamma_V}\mathrm{coef}_K(v,a)(c\gamma g)$.
\item A \emph{frame} $A$ of $V$ is a family of matrices $A_\tau$ with signs (an orthonormalisation of $V^\tau$); it defines the norms
$\mathrm{nrm}_\tau(A,x;y)=|A_\tau(x^{-1})_\tau\,\tau(y)|$ of vectors $y\in E^{\dim V}$ at $x\in G_V$. For $s>0$, $x$ is
\emph{$s$-long} for $(A,\Lambda)$ if for every non-zero ideal $\mathfrak a$ and every non-zero $y\in\mathfrak a\Lambda_x$,
$s^{f}\sqrt{\Nm\mathfrak a}\le\prod_\tau\mathrm{nrm}_\tau(A,x;y)$. A function $\phi$ on $G_V$ is \emph{$c$-Lipschitz} for $A$ if
$|\phi(x\exp(sX))-\phi(x)|\le c|s|\,\|A_\tau XA_\tau^{-1}\|$ for all $x$, $\tau$, $X$ in the Lie algebra at $\tau$ and $s\in\R$.
\item Let $\mathcal M_0$ be a model of $Y_0$. A \emph{pack} for $(\mathcal M_0,\mathcal M_{\fsc})$ consists of frames $A$ of $Y_0$ and
$B$ of $V_{\fsc_0}$, vectors $v_{0,\tau}\in H_\tau(\mathcal M_{\fsc})$, a point $g_0\in G_{V_{\fsc_0}}$ and a smooth vector
$w=(w_\tau)$ of $\mathcal M_0$. Put $v_\tau=\pi_\tau((g_0)_\tau)v_{0,\tau}$. For $Q\in\R$ the pack is \emph{$Q$-good} if: $A$ and $B$
are adapted to the decomposition $V_{\fsc_0}=Y_0\oplus\ell$; $g_0$ is archimedean; every $v_{0,\tau}$ is $K$-finite with
$\langle v_{0,\tau},\pi_\tau(g)v_{0,\tau}\rangle(1+\|g\|)^{2\dim V_{\fsc_0}^2+1}$ bounded; every $w_\tau\ne0$; and for every Haar
measure $dh$ on $G_{Y_0,\infty}$ there is $c\ne0$ with
\begin{enumerate}[label=(\roman*)]
\item $\int\mathrm{coef}(v,v)(\iota_0h)\,\prod_\tau\langle b_\tau,\pi^0_\tau(h_\tau)w_\tau\rangle\,dh=c\prod_\tau\langle b_\tau,w_\tau\rangle$ for all
$b$ (an orthogonality relation for the restriction to $G_{Y_0,\infty}$);
\item for every $g\in G_{V_{\fsc_0}}$, $h\mapsto\mathrm{coef}(v,v)(g\iota_0h)$ is integrable and
\[
\Bigl(\int|\mathrm{coef}(v,v)(g\iota_0h)|\,dh\Bigr)\Bigl(\prod_\tau\mathrm{nrm}_\tau(B,g^{-1};z_\ell)\Bigr)^{2\dim V_{\fsc_0}+5}\le e^Q|c|;
\]
\item a Lipschitz bound with constant $e^Q|c|$ for the integral of
$|\mathrm{coef}(v,v)(\iota_0(\exp(sX)h))-\mathrm{coef}(v,v)(\iota_0h)|$.
\end{enumerate}
\item The \emph{loss} is $\mathrm{loss}(D,C,b)=Cf(r+1)^4\log b$. $\fsP{PX}$ (packs exist) is the statement: there is $C\ge0$ such
that $\forall_{\mathrm{ev}}$, for all models $\mathcal M_0$ of $Y_0$ with parameter $S^Y$ and $\mathcal M_{\fsc}$ of $V_{\fsc_0}$ with
parameter $S^{\fsc}_0$, there is a $\mathrm{loss}(D,C,b)$-good pack.
\item $\mathrm{Pays}(D,Q,N)$ means: there is $s\in(0,\frac12]$ with
$e^Q\,r^2f\,g_0g_1\,\mathrm{rd}_E\,(2s)^{1/m}\le\frac12$, $2\exp(Q+(6(r+1)+9)f)\le s^f\sqrt N$ and $1<s^f\sqrt N$, where $\mathrm{rd}_E$ is
the root discriminant of $E$. $\fsP{DP}$ (depth pays) is the statement: for every $C\ge0$, $\forall_{\mathrm{ev}}$,
$\mathrm{Pays}(D,\mathrm{loss}(D,C,b),\Nm\fsm)$.
\end{enumerate}
\formal{GlobalTheory.SepPer.coefR}
\formal{GlobalTheory.SepPer.coef}
\formal{GlobalTheory.SepPer.poin}
\formal{GlobalTheory.SepPer.nrm}
\formal{GlobalTheory.SepPer.Long}
\formal{GlobalTheory.SepPer.Lip}
\formal{GlobalTheory.SepPer.Pack}
\formal{GlobalTheory.SepPer.Pack.Is}
\formal{GlobalTheory.SepPer.loss}
\formal{GlobalTheory.SepPer.Pays}
\formal{GlobalTheory.SepPer.PackEx}
\formal{GlobalTheory.SepPer.DepthPays}
\end{definition}

The archimedean input gives packs (\cref{pf:fs:packEx} below: $\fsP{RI}$ implies $\fsP{PX}$), and the depth of the grid pays for their loss.

\begin{proposition}[Depth pays; \lref{GlobalTheory.T.gridDepth_pays_every_loss}]\label{pf:fs:packdepth}
$\fsP{DP}$ holds.
\end{proposition}

This is a computation with the explicit depth: the depth $d_0(C_F,b)$ at the base place makes $\Nm\fsm\ge q_0^{d_0}$ large enough to absorb
$e^{\mathrm{loss}}$; the threshold $C_F\ge6510$ enters here, and the scale must exceed $C+1$. We do not reproduce it.

\begin{theorem}[Non-vanishing of core periods]\label{pf:fs:cuspPer}
Assume $\sfname{HarishChandraModules}$, $\sfname{CasselmanWallach}$, $\sfname{CuspRapidDecay}$, $\fsP{PX}$ and $\fsP{DP}$. Then
$\forall_{\mathrm{ev}}$: for all models $\mathcal M_0$ of $Y_0$ with parameter $S^Y$ and $\mathcal M_{\fsc}$ of $V_{\fsc_0}$ with
parameter $S^{\fsc}_0$, and every non-zero $t\in\fsRows(\mathcal M_0)$, there are a core form $y$ of $\mathcal M_{\fsc}$ and
$a\in\fsHi(\mathcal M_{\fsc})$, $w'\in\fsHi(\mathcal M_0)$ with $\fsper_0(y(a),t(w'))\ne0$.
\formallabel{UpperTheory.Lit.Cusp.exists_cusp_form_theta_period_ne_zero}
\end{theorem}

\begin{proof}
Let $C$ be the constant of $\fsP{PX}$. Given $\eps$ and $C_F\ge6510$, let $x_1$ be the threshold of $\fsP{PX}$ and $x_2$ that of
$\fsP{DP}$ for $C$; put $x_0=\max(x_1,x_2)$. Fix the data of $\forall_{\mathrm{ev}}$, $\mathcal M_0$, $\mathcal M_{\fsc}$ and $t\ne0$. Write
$Q=\mathrm{loss}(D,C,b)$, $K=K_{\fsc_0}$ (the congruence subgroup of $\Lambda_{\fsc_0}$ modulo $\fsm$), $\nu=\nu_{(0,0)}$ and
$\fsF=\fsF_{(0,0)}$ (a Haar measure and a fundamental domain, $M$ being a row). We use the following facts about the bottom of the
grid (\lref{GlobalTheory.T.thetaRowGrid_bottom_selfDual_lattices}): $\Lambda_Y(0)$ and $\Lambda_{\fsc_0}$ are self-dual;
$\Lambda_{\fsc_0}$ is the lattice of the decomposition $Y_0\oplus\ell$ built on $\Lambda_Y(0)$ and contains $z_\ell$; if $k\in G_{Y_0,f}$
and $\iota_0(k)\in K$ then $k\in K_{\mathrm{row}}$; $\dim Y_0=r$, $\dim V_{\fsc_0}=r+1$; and every frame of $Y_0$ has at most $m$
negative signs at every $\tau$.

\emph{Step 1: a pack and a large value.} By $\fsP{PX}$ there is a $Q$-good pack $p=(A,B,v_0,g_0,w)$; by $\fsP{DP}$ there is $s$ as
in $\mathrm{Pays}(D,Q,\Nm\fsm)$. Let $\phi=t(\otimes_\tau w_\tau)$, a non-zero cusp form on $G_{Y_0}$, right invariant under
$K_{\mathrm{row}}$ (as $t$ is a row) and of rapid decay. Since $t\ne0$ and every $w_\tau\ne0$, there are $b\ge0$ and $x\in G_{Y_0}$ with
$|\phi|\le b$ everywhere and $|\phi(x)|>b/2$ (\lref{GlobalTheory.T.exists_bound_multiplicity_value_exceeds_half}: a non-zero
multiplicity map is non-zero on every pure tensor of non-zero vectors, and a bounded continuous function approaches its supremum).
By the Lipschitz clause of the pack, $\phi$ is $e^Qb$-Lipschitz for $A$
(\lref{GlobalTheory.T.unipotentPeriod_value_lipschitz_realOneParameter}).

\emph{Step 2: $x$ is long.} Suppose $x$ is not $s$-long for $(A,\Lambda_Y(0))$. A $K_{\mathrm{row}}$-invariant cusp form which is
$c$-Lipschitz satisfies, at a point which is not $s$-long,
\begin{equation}\label{eq:pf:fs:shortline}
|\phi(x)|\le c\,r^2f\,g_0g_1\,\mathrm{rd}_E\,(2s)^{1/m}
\end{equation}
(\cref{pf:fs:shortLine}, applied to $V=Y_0$ with $\dim Y_0=r$, the level $g=g_0g_1$ and $q=m$, every frame of $Y_0$ having at most $m$
negative signs at every $\tau$). With $c=e^Qb$ and the first inequality of $\mathrm{Pays}$,
\eqref{eq:pf:fs:shortline} gives $|\phi(x)|\le b/2$, a contradiction. Hence $x$ is $s$-long.

\emph{Step 3: the adelic clauses.} The real clauses of the pack and the facts on the bottom lattices give a constant $c\ne0$ with
the \emph{adelic clauses} for the lattice $\Lambda_{\fsc_0}$, the modulus $\fsm$, the measure $\nu$ and $\phi$
(\lref{GlobalTheory.T.unipotentPeriod_adelic_clauses_of_real}): the identity term reproduces $\phi$,
\begin{equation}\label{eq:pf:fs:idterm}
\int\mathrm{coef}_K(v,v)\bigl(\iota_0(x)^{-1}\iota_0(h)\bigr)\,\phi(h)\,d\nu(h)=c\,\phi(x)\qquad(x\in G_{Y_0}),
\end{equation}
and for $\gamma\in\Gamma_{V_{\fsc_0}}$ with $\gamma z_\ell=z$, the term of $\gamma$ is integrable, its integral is bounded by
\[
e^Q|c|\Bigl(\prod_\tau\mathrm{nrm}_\tau(B,\iota_0x;z)\Bigr)^{-(2r+7)},
\]
and it vanishes identically unless $z-z_\ell\in\fsm\Lambda_{\iota_0x}$.

\emph{Step 4: the identity term dominates.} Since $x$ is $s$-long and, by $\mathrm{Pays}$,
$2\exp(Q+(6(r+1)+9)f)\le s^f\sqrt{\Nm\fsm}$ and $1<s^f\sqrt{\Nm\fsm}$, the sum over $\gamma\ne1$ of the bounds of Step~3 is less
than $|c|\,b/2<|c\phi(x)|$; hence
\begin{equation}\label{eq:pf:fs:poinne}
\int_{\fsF}P_K\bigl(v,\iota_0(x)^{-1},v\bigr)(\iota_0h)\,\phi(h)\,d\nu(h)\neq0
\end{equation}
(\cref{pf:fs:idDom}, applied with $Y=Y_0$, $V=V_{\fsc_0}$, $\Lambda=\Lambda_Y(0)$ and the adelic clauses of Step~3).

\emph{Step 5: the Poincar\'e series is a core form.} The vector $v_0$ is $K$-finite and its coefficients decay fast
(\lref{GlobalTheory.T.coefficientDecay_kFinite_second_vector}, from $\sfname{HarishChandraModules}$). By
\cref{pf:fs:poinmult} below, applied with the base point $c=g_0^{-1}\iota_0(x)^{-1}$, there is $y\in\fsA(\mathcal M_{\fsc};K,1)$ (a
core form) with $y(\otimes a)=P_K(v_0,c,a)$ for all families $a$ of $K$-finite vectors. Let $a'_\tau=\pi_\tau((g_0)_\tau)v_{0,\tau}=v_\tau$;
equivariance of $y$ at the archimedean indices gives $y(\otimes a')(g)=y(\otimes v_0)(g\,g_0)$, and moving the archimedean point $g_0$
through the Poincar\'e series (\lref{GlobalTheory.T.coefficientPoincareSeries_move_real_point}:
$P_K(v_0,g_0^{-1}c',v_0)(gg_0)=P_K(v,c',v)(g)$ for archimedean $g_0$) gives $y(\otimes a')=P_K(v,\iota_0(x)^{-1},v)$. By
\eqref{eq:pf:fs:poinne}, $\fsper_0(y(\otimes a'),\phi)\ne0$, which is the assertion with $a=\otimes a'$ and $w'=\otimes w$.
\end{proof}

\begin{lemma}[Poincar\'e series are core forms]\label{pf:fs:poinmult}\sloppy
Assume $\sfname{HarishChandraModules}$, $\sfname{CasselmanWallach}$ and $\sfname{CuspRapidDecay}$. Let $\mathcal M$ be a model of $V$,
$\Lambda$ a self-dual lattice, $\fsm\ne0$ an ideal, $K$ the principal congruence subgroup of $\Lambda$ modulo $\fsm$,
$v=(v_\tau)$ $K$-finite with coefficients decaying faster than $(1+\|g\|)^{-(2\dim V^2+1)}$, and $c\in G_V$. If $\fsT\ne\emptyset$,
there is $y\in\fsA(\mathcal M;K,1)$ with $y(\otimes a)=P_K(v,c,a)$ for every family $a$ of $K$-finite vectors.
\formallabel{UpperTheory.Lit.Cusp.exists_multT_eq_poin}
\end{lemma}

\begin{proof}
The map $a\mapsto P_K(v,c,a)$ on $K$-finite vectors is multilinear; each value is an automorphic form, right $K$-invariant and
equivariant at every $\tau$ (\lref{GlobalTheory.T.coefficientPoincareSeries_isAutomorphic}: absolute convergence of the series and
invariance), smooth of moderate growth together with its derivatives
(\lref{GlobalTheory.T.poincareSeries_smooth_moderateGrowth_deriv}, using $\sfname{HarishChandraModules}$), and cuspidal
(\cref{pf:fs:poinCusp}). By $\sfname{CuspRapidDecay}$ the values are
of rapid decay. These properties make the map a $(\mathfrak g,K_\infty)$-module map from the $K$-finite vectors to cusp forms with
Petersson bounds, which extends to the smooth vectors as a multiplicity map with the same values on $K$-finite vectors
(\lref{GlobalTheory.GKExtensionToSmoothVectors.gkExtension}, an application of the Casselman--Wallach globalisation, using
$\sfname{HarishChandraModules}$ and $\sfname{CasselmanWallach}$).
\end{proof}

\begin{lemma}[From one vector to all vectors]\label{pf:fs:kf}\sloppy
Assume $\sfname{CuspRapidDecay}$ and $\sfname{HarishChandraModules}$. In a setting, let $\mathcal M_0$ be a model of $Y_0$ with
parameter $S^Y$, $\mathcal M_{\fsc}$ a model of $V_{\fsc_0}$ with parameter $S^{\fsc}_0$, $t\in\fsRows(\mathcal M_0)$ and $y$ a core form
of $\mathcal M_{\fsc}$. If $\fsper_0(y(a'),t(w'))\neq0$ for some $a',w'$, then for every non-zero $w\in\fsHi(\mathcal M_0)$ there is
$a\in\fsKF(\mathcal M_{\fsc})$ with $\fsper_0(y(a),t(w))\neq0$.
\formallabel{GlobalTheory.PeriodNePf.exists_kf}
\end{lemma}

\begin{proof}
Put $F=y(a')\circ\iota_0$, a bounded continuous left $\Gamma_{Y_0}$-invariant function (values of multiplicity maps are rapidly
decreasing). If $\int_{\fsF}F(g)\,t(w)(gh)\,d\nu(g)=0$ for all archimedean $h\in G_{Y_0,\infty}$, then
$\int_\fsF F\,t(w'')\,d\nu=0$ for all $w''$ (\lref{GlobalTheory.T.integral_eq_zero_of_orthogonal_real_translates}: the
functional $w''\mapsto\int_\fsF F\,t(w'')\,d\nu$ vanishes on all archimedean translates of $w$, and integrating against matrix
coefficients, the Schur orthogonality relations of the discrete series at the archimedean indices show that it vanishes on all
pure tensors); with $w''=w'$ this contradicts the hypothesis. So there is archimedean $h$ with
$\int_\fsF F(g)\,t(w)(gh)\,d\nu(g)\ne0$. The substitution $g\mapsto gh^{-1}$ (invariance of the integral over a fundamental domain
of a left-invariant function under right translation by $h$) turns this into
$\int_\fsF y(a')(\iota_0(g)\iota_0(h)^{-1})\,t(w)(g)\,d\nu(g)\ne0$. By equivariance of $y$ at the archimedean indices,
$y(a')(\cdot\,\iota_0(h)^{-1})=y(a'')$ for the vector $a''$ obtained from $a'$ by the action of $\iota_0(h)^{-1}$, so
$\fsper_0(y(a''),t(w))\ne0$. Finally, if $\fsper_0(y(a),t(w))=0$ for all $a\in\fsKF(\mathcal M_{\fsc})$, then
$\fsper_0(y(a),t(w))=0$ for all $a$ (\lref{GlobalTheory.T.period_eq_zero_of_vanishing_on_kFinite_values}: the functional
$a\mapsto\int y(a)(\iota_0g)\,t(w)(g)\,d\nu$ is continuous for a Sobolev norm and the $K$-finite vectors are dense,
$\sfname{HarishChandraModules}$), contradicting $\fsper_0(y(a''),t(w))\ne0$.
\end{proof}

\begin{theorem}[Period non-vanishing]\label{pf:fs:periodNe}
$\sfname{HarishChandraModules}$, $\sfname{CasselmanWallach}$, $\sfname{CuspRapidDecay}$ and $\fsP{RI}$ imply $\fsP{PN}$.
\formallabel{UpperTheory.Lit.Branch.periodNe_of_eigen}
\end{theorem}

\begin{proof}
$\fsP{RI}$ gives $\fsP{PX}$ (\cref{pf:fs:packEx}) and $\fsP{DP}$ holds (\cref{pf:fs:packdepth}), so \cref{pf:fs:cuspPer} applies. Given $\eps,C_F$, take its
threshold. For the data of $\fsP{PN}$, \cref{pf:fs:cuspPer} gives $y,a',w'$ with $\fsper_0(y(a'),t(w'))\ne0$, and \cref{pf:fs:kf} gives
$a\in\fsKF(\mathcal M_{\fsc})$ with $\fsper_0(y(a),t(w))\ne0$.
\end{proof}

% ---------------------------------------------------------------------------
\subsection{Core forms are reachable}\label{sec:pf:fs:onto}

This subsection proves $\fsP{CO}$: every core form at $\fsc_0$ is a combination of downward lifts, along the core column, of
column forms at the top. The heart of the matter is a surjectivity statement for composites of downward theta lifts between
finite-dimensional spaces of multiplicity maps, proved by an adjoint argument in the style of Rallis: the upward lifts are adjoint to
the downward ones, they transfer Hecke parameters, and they do not vanish on eigenvectors whose partial $L$-value at the relevant
point is non-zero.

\subsubsection{Chains of stages}

\begin{definition}[Common data, stages and steps]\label{pf:fs:stages}
\begin{enumerate}[label=(\alph*)]
\item \emph{Common data} $c$ consist of $\theta$, $\psi$, a finite set $S$ of places of $F$, integers $n,t',m'$, Hecke characters
$\chi_k$ and exponents $a_k(\tau)$ ($k<m'$). They are \emph{good} if $\theta$ is a family of archimedean embeddings, $\psi$ is standard,
each $\chi_k$ is a Hecke character with trivial restriction to $\A_F^\times$ and infinity type $a_k$, unramified outside $S$, the
exponents $a_i(\tau),a_k(\tau)$ are distinct for $i<k$, $1\le n$ and $n+2\le t'$.
\item A \emph{stage} $s$ is a hermitian space $V_s$ (with a unimodular lattice $\Lambda_s$), a parameter, a model $\mathcal M_s$, an open
compact subgroup $J_s$, a Haar measure $\nu_s$ and a fundamental domain $\fsF_s$. Its multiplicity space is
$\mathrm{CC}(s)=\fsA(\mathcal M_s;J_s,1)$. The stage is \emph{good of height $j$} for $c$ if $\dim V_s=n+jt'$, $V_s$ has consistent
signatures and Witt index $<t'$, $\Lambda_s$ is self-dual, $J_s$ stabilises $\Lambda_s$ and contains the elements which are integral
outside $S$ and trivial at $S$, $J_s$ is spherical off $S$ and splits at the places off $S$, and $(\nu_s,\fsF_s)$ is a Haar measure
with a fundamental domain.
\item A \emph{step} $p$ from $s$ to $s'$ consists of a Weil datum $\omega$ for $(V_s,V_{s'})$, a theta function $f$, a bilinear map
$\mathrm{It}$ from polynomials and $\fsHi(\mathcal M_{s'})$ to $\fsHi(\mathcal M_s)$ and a linear map
$\mathrm{io}\colon\mathrm{CC}(s')\to\mathrm{CC}(s)$. It is a \emph{good step of index $k$} if $\omega$ has the splitting characters
$(\chi_k,\chi_k)$, the parameters of $s,s'$ are related by the upward pair condition for a segment of length $t'$ with the exponents
$a_k$, $f$ is a vacuum vector of $(\mathcal M_s,\mathcal M_{s'})$ spherical off $S$, the Fock vectors over $f$ are fixed by
$J_s\times J_{s'}$, their Weil coefficients vanish on the integral elements off $S$ outside $G_{V_s,\infty}J_s$, $\mathrm{It}$ is an
intertwiner, and $\mathrm{io}$ is the downward lift: $\mathrm{io}(y)(\mathrm{It}(P,v))$ is the theta lift of $y(v)$ by the Fock
vector over $f$ attached to $P$.
\item \emph{Reach} is the relation generated by: $s_0$ reaches itself at height $0$ by the identity; if $s_0$ reaches $s$ at height
$k$ by $\Theta\colon\mathrm{CC}(s)\to\mathrm{CC}(s_0)$, $s'$ is good of height $k+1$ and $p$ is a good step of index $k$ from $s$ to
$s'$, then $s_0$ reaches $s'$ at height $k+1$ by $\Theta\circ\mathrm{io}_p$.
\item $\fsP{OA}(D)$ is the statement: for all good common data $c$, every stage $s_0$ good of height $0$ and every chain by which $s_0$
reaches a stage $s$ at height $k\le m'$, the composite $\Theta\colon\mathrm{CC}(s)\to\mathrm{CC}(s_0)$ is surjective.
\end{enumerate}
\formal{GlobalTheory.EvenOnto.Com}
\formal{GlobalTheory.EvenOnto.Com.Is}
\formal{GlobalTheory.EvenOnto.PiLocalization₂}
\formal{GlobalTheory.EvenOnto.PiLocalization₂.Is}
\formal{GlobalTheory.EvenOnto.PreStep}
\formal{GlobalTheory.EvenOnto.Step}
\formal{GlobalTheory.EvenOnto.Step.Is}
\formal{GlobalTheory.EvenOnto.Reach}
\formal{GlobalTheory.EvenOnto.OntoAll}
\end{definition}

The proof of $\fsP{OA}$ uses the following properties of good stages and good steps, each a statement of the development for all good
common data, stages and steps of the relevant heights. For an element $u$ of $\mathrm{CC}(s)$ and a family $B$ of Satake parameters
off $S$, $u$ \emph{has parameters} $B$ if $u\ne0$ and every non-zero value $u(v)$ has Satake parameters $B$ off $S$ (in the sense of
\cref{pf:liftinj:coefdef}). The upward lift of $F$ on $G_{V_s}$ by a step $p$ and a polynomial $P$ is
$\mathrm{up}_p(P,F)(h)=\int_{\fsF_s}\overline{\theta_\omega(g,h;\mathrm{pt}_f(P))}F(g)\,d\nu_s(g)$, and the transferred parameters
$\mathrm{step}_k(B)$ add a segment of length $t'$ twisted by $\chi_k$; $\mathrm{chain}_k(B)$ is the $k$-fold transfer.
\begin{enumerate}[label=(O\arabic*)]\sloppy
\item\label{pf:fs:O1} (finite dimension and metrics; \lref{GlobalTheory.T.evenThetaChain_finiteDimensional_and_peterssonProducts},
from $\sfname{HarishChandraFinite}$ and $\sfname{CuspRapidDecay}$) $\mathrm{CC}(s)$ is finite-dimensional and carries a positive
definite hermitian form $\langle\cdot,\cdot\rangle_s$ such that $\int_{\fsF_s}x(v)\overline{y(v')}\,d\nu_s=\langle x,y\rangle_s\langle v,v'\rangle$
(a Schur relation for the Petersson product).
\item\label{pf:fs:O2} (eigenvectors; \lref{GlobalTheory.T.even_theta_chain_spanned_params_of_value_independent}) $\mathrm{CC}(s)$ is
spanned by elements with parameters; an element whose value has parameters $B$ has parameters $B$; elements with pairwise distinct
parameters (at some place off $S$) are linearly independent.
\item\label{pf:fs:O3} (trivial bound; \lref{GlobalTheory.T.evenRankChain_params_trivialBound_rapidDecay}) the parameters of
elements of $\mathrm{CC}(s)$ are bounded by $\Nm w^{(\dim V_s-1)/2}$.
\item\label{pf:fs:O4} (upward lifts; \lref{GlobalTheory.T.thetaLiftUpward_cuspSpan_and_adjoint}, from $\sfname{DiscreteSeriesCover}$,
$\sfname{HarishChandraModules}$, $\sfname{CuspRapidDecay}$, $\sfname{CasselmanWallach}$) for $u\in\mathrm{CC}(s)$ and $K$-finite $v$,
$\mathrm{up}_p(P,u(v))$ lies in the span of the values of $\mathrm{CC}(s')$ on $K$-finite vectors, and the upward and downward lifts are
adjoint: $\langle y(w),\mathrm{up}_p(P,u(v))\rangle_{s'}=\langle\mathrm{io}(y)(\mathrm{It}(P,w)),u(v)\rangle_s$.
\item\label{pf:fs:O5} (transfer of parameters; \lref{GlobalTheory.T.evenThetaChainUpwardLift_heckeParam_transferred}, from
$\sfname{CuspRapidDecay}$ and $\sfname{HoweUnram}$) if $u$ has parameters $B$ and $\mathrm{up}_p(P,u(v))\ne0$, it has parameters
$\mathrm{step}_k(B)$.
\item\label{pf:fs:O6} (zeta integrals; \lref{GlobalTheory.T.evenThetaStep_realZeta_nonvanishing} from
$\sfname{DiscreteSeriesCover}$, $\sfname{HarishChandraModules}$, and
\lref{GlobalTheory.T.evenThetaStep_semilocalZeta_ne_zero_of_real}) the archimedean zeta integral of the Weil coefficient against
the matrix coefficient of a $K$-finite vector is non-zero and remains integrable after multiplication by the power
$\mathrm{ht}^{s_1-t'/2}$ of the height for some $s_1<t'/2$; hence the semilocal zeta integrals over the elements integral off $S$
are non-zero for every non-zero $u$.
\item\label{pf:fs:O7} (non-vanishing of upward lifts; \lref{GlobalTheory.T.upwardThetaLift_ne_zero_of_lValue_ne_zero}) if $u$ has
parameters $B$ and the partial $L$-function of $B$ twisted by $\chi_k$ has a non-zero value at the relevant point, then
$\mathrm{up}_p(P,u(v))\ne0$ for some $P$ and $K$-finite $v$. The input is the Rallis inner product formula in the form
$\sfname{HeckeForm}(D)$: the Petersson norm of an upward lift is the product of the partial $L$-value and the semilocal zeta integral
(\lref{GlobalTheory.Core.Rk.IP.heckeForm_of_doubling'}, from $\sfname{WeilRepresentationUnitaryPair}$, $\sfname{CuspRapidDecay}$,
$\sfname{WeilSiegelFormulaUnitary}$, $\sfname{SiegelWeilUnitaryConvergent}$, $\sfname{DoublingMethodLocalZetaIntegrals}$ and
$\sfname{UnramifiedDoublingGL}$).
\item\label{pf:fs:O8} (chain $L$-values; \lref{GlobalTheory.T.evenThetaChain_LValues_ne_zero_and_params_injective}) if $B$ satisfies
the trivial bound in rank $n$, the $L$-values of $\mathrm{chain}_k(B)$ twisted by $\chi_k$ are non-zero for all $k<m'$; and
$\mathrm{chain}_k(B)$ determines $B$ placewise.
\end{enumerate}

\begin{lemma}[One step up]\label{pf:fs:stepUp}
Let $p$ be a good step of index $k$ from $s$ to $s'$, and $u\in\mathrm{CC}(s)$ with parameters $B$ whose twisted $L$-value is
non-zero. Then there is $x\in\mathrm{CC}(s')$ with parameters $\mathrm{step}_k(B)$ and $\langle y,x\rangle_{s'}=\langle\mathrm{io}(y),u\rangle_s$ for all
$y\in\mathrm{CC}(s')$.
\formallabel{GlobalTheory.EvenOnto.stepUp}
\formallabel{GlobalTheory.T.even_theta_step_adjoint_nonzero_transferred_params}
\end{lemma}

\begin{proof}
$\mathrm{CC}(s')$ is finite-dimensional with a positive definite form \ref{pf:fs:O1}, so there is a unique $x$ with
$\langle y,x\rangle_{s'}=\langle\mathrm{io}(y),u\rangle_s$ for all $y$ (\lref{GlobalTheory.hermitian_adjoint_exists_unique}). By
\ref{pf:fs:O7} there are $P$ and a $K$-finite $v$ with $F=\mathrm{up}_p(P,u(v))\ne0$. By \ref{pf:fs:O4}, $F$ lies in the span of
values of $\mathrm{CC}(s')$, and for all $y$ and $K$-finite $w$,
$\langle y(w),F\rangle=\langle\mathrm{io}(y)(\mathrm{It}(P,w)),u(v)\rangle=\langle\mathrm{io}(y),u\rangle_s\langle\mathrm{It}(P,w),v\rangle
=\langle y,x\rangle_{s'}\langle\mathrm{It}(P,w),v\rangle$ by \ref{pf:fs:O1}. A function in the span of the values with these inner
products is a value of $x$: $F=x(w_0)$ for some $w_0$ (\lref{GlobalTheory.ChainOnto.value_of_adjoint_of_mem_span}). By
\ref{pf:fs:O5}, $F$ has parameters $\mathrm{step}_k(B)$, and by \ref{pf:fs:O2} so has $x$.
\end{proof}

\begin{proposition}[Surjectivity of chains of lifts]\label{pf:fs:ontoAll}\sloppy
$\sfname{HarishChandraFinite}$, $\sfname{CuspRapidDecay}$, $\sfname{DiscreteSeriesCover}$, $\sfname{HarishChandraModules}$,
$\sfname{CasselmanWallach}$, $\sfname{HoweUnram}$, $\sfname{WeilRepresentationUnitaryPair}$, $\sfname{WeilSiegelFormulaUnitary}$,
$\sfname{SiegelWeilUnitaryConvergent}$, $\sfname{DoublingMethodLocalZetaIntegrals}$ and $\sfname{UnramifiedDoublingGL}$ imply
$\fsP{OA}(D)$ for every $D$.
\formallabel{GlobalTheory.EvenOnto.even_surjective_of_representatives}
\formallabel{GlobalTheory.T.twist_even_surjective_of_representatives}
\formallabel{GlobalTheory.EvenOnto.ontoAll_of}
\formallabel{GlobalTheory.T.evenRankThetaChainLift_surjective_of_adjoint_nonzero}
\formallabel{GlobalTheory.EvenOnto.chainUp_of}
\end{proposition}

\begin{proof}
The hypotheses give \ref{pf:fs:O1}--\ref{pf:fs:O8}, hence \cref{pf:fs:stepUp}. Let $s_0$ reach $s$ at height $k$ by $\Theta$.
\emph{Chain step.} For every $u\in\mathrm{CC}(s_0)$ with parameters $B$ there is $x\in\mathrm{CC}(s)$ with parameters
$\mathrm{chain}_k(B)$ and $\langle y,x\rangle_s=\langle\Theta(y),u\rangle_{s_0}$ for all $y$: by induction along the chain, using
\cref{pf:fs:stepUp} at each step, whose $L$-value hypothesis holds by \ref{pf:fs:O8} since $B$ satisfies the trivial bound in rank
$\dim V_{s_0}=n$ \ref{pf:fs:O3} (this step also shows that all stages of the chain are
good). \emph{Linear algebra.} $\Theta\colon\mathrm{CC}(s)\to\mathrm{CC}(s_0)$ is a linear map into a finite-dimensional space with a
positive definite hermitian form, $\mathrm{CC}(s_0)$ is spanned by elements with parameters \ref{pf:fs:O2}, the elements of
$\mathrm{CC}(s_0)$ with parameters $B$ form, together with $0$, a subspace depending only on the class of $B$ off $S$
(\lref{GlobalTheory.T.multiplicity_vanishing_and_same_params_subspace}), the transfer $\mathrm{chain}_k$ is injective on classes
\ref{pf:fs:O8}, elements of $\mathrm{CC}(s)$ with distinct parameters are independent \ref{pf:fs:O2}, and every $u$ with parameters $B$
has an adjoint preimage $x$ with parameters $\mathrm{chain}_k(B)$ (chain step). These conditions imply that $\Theta$ is surjective
(\lref{GlobalTheory.T.surjective_of_adjoint_separates_labelled_parts}: if $u\ne0$ is orthogonal to the image of $\Theta$, decompose
it into components with parameters; the adjoint preimages $x_j$ of the components have distinct parameters, hence are independent,
while $\langle y,\sum_jx_j\rangle=\langle\Theta(y),u\rangle=0$ for all $y$ forces $\sum_jx_j=0$).
\end{proof}

\subsubsection{From chains of stages to the core column}

The core column of a grid is a chain of stages of this kind. The common data of a Weil datum $W$ are $\theta,\psi$ of $W$,
$S=S_0$, $n=\dim V_{\fsc_0}$ and $t'$ the length of the core segments (so that $\dim V_{\fsc_i}=n+it'$), $m'=m$,
$\chi_k=\chi_{m-k}^{-1}$ and the exponents given by the centres of the grid. For a unimodular lattice $r$ of $V_{\fsc_i}$ which is
\emph{a core lattice} (its frame lattice is $(\bar{\mathfrak A}\mathfrak A^{-1})^i\Lambda_{\fsc_i}$, $\mathfrak A$ an ideal of the host), a
core volume $v$ and a model $\mathcal M$ of $V_{\fsc_i}$, the \emph{grid stage} $s_i(r,v,\mathcal M)$ has space $V_{\fsc_i}$ with
lattice $r$, model $\mathcal M$, level $K_{\fsc_i}$ and the measure and domain of $v$.

\begin{definition}[Four statements on the core column]\label{pf:fs:corestatements}
\begin{enumerate}[label=(\arabic*)]
\item $\fsP{COM}$, $\forall_{\mathrm{gap}}$: the common data of $W$ are good.
\item $\fsP{STG}$, $\forall_{\mathrm{all}}$: for $i\le m$ there is a core lattice of $V_{\fsc_i}$, and every grid stage
$s_i(r,v,\mathcal M)$ with a core lattice $r$ is good of height $i$.
\item $\fsP{CHN}$, $\forall_{\mathrm{gap}}$: for every core volume $v$, core lattice $r_0$ at $\fsc_0$ and model $\mathcal M_0$ of
$V_{\fsc_0}$, and every $k\le m$, there are a core lattice $r$ and a model $\mathcal M$ at $\fsc_k$ and a chain by which
$s_0(r_0,v,\mathcal M_0)$ reaches $s_k(r,v,\mathcal M)$ with composite $\Theta$, such that for every $y\in\mathrm{CC}(s_k)$ and every
$K$-finite $x$, $\Theta(y)(x)$ lies in the span of the bottoms $\Psi_0$ of the families $(\Psi_i)_{i\le k}$ with $\Psi_i$ rapidly
decreasing, $\Psi_i=\fsdn_i(f_i,\Psi_{i+1})$ for depth vacua $f_i$ ($i<k$) and $\Psi_k=y(w)$ for some $w$.
\item $\fsP{TC}$, $\forall_{\mathrm{all}}$: for every model $\mathcal M$ of $V_{\fsc_m}$ with parameter $S^{\fsc}_m$, every core form
$y$ at $\fsc_m$ and every $w$, the function $y(w)$ is the scalar form $\Psi_\xi$ of some column form $\Psi$ and some $\xi$.
\end{enumerate}
\formal{GlobalTheory.CoreOnto.gridCom}
\formal{GlobalTheory.CoreOnto.IsCoreLat}
\formal{GlobalTheory.CoreOnto.gridStage}
\formal{GlobalTheory.CoreOnto.Below}
\formal{GlobalTheory.CoreOnto.ComOK}
\formal{GlobalTheory.CoreOnto.StageOK}
\formal{GlobalTheory.CoreOnto.ChainEx}
\formal{GlobalTheory.CoreOnto.TopCols}
\end{definition}

These hold: $\fsP{COM}$ by \lref{GlobalTheory.T.testGridCommonData_properties_hold}; $\fsP{STG}$ by
\lref{GlobalTheory.T.testGridStageLatticesProperties_proof}; $\fsP{CHN}$, given $\sfname{DiscreteSeriesCover}$,
$\sfname{HarishChandraModules}$, $\sfname{CuspRapidDecay}$, $\sfname{CasselmanWallach}$, $\fsP{DF}$ and $\fsP{CM}$, by \cref{pf:fs:chainEx} below
(the steps are built from the Weil representations of the core pairs, and their downward lifts are the lifts $\fsdn_i$ by depth vacua); and
$\fsP{TC}$ by \lref{GlobalTheory.T.definiteTopSpace_columnValues} (the top space $V_{\fsc_m}=V_{\fst_0}$ is totally definite).

\begin{proposition}[Reachability from chains]\label{pf:fs:coreOntoChain}
$\fsP{OA}$ (for all $D$), $\fsP{COM}$, $\fsP{STG}$, $\fsP{CHN}$ and $\fsP{TC}$ imply $\fsP{CO}$.
\formallabel{GlobalTheory.CoreOnto.coreOnto_of_chain}
\formallabel{GlobalTheory.T.coreChain_lifts_surjective_of_evenRank}
\end{proposition}

\begin{proof}
Fix a setting with $\sfname{Sizes}(D)$ and $c_{\mathrm{gap}}=1$, a core volume $v$, a model $\mathcal M_{\fsc}$ of $V_{\fsc_0}$, a core
form $y$ and $x\in\fsKF(\mathcal M_{\fsc})$. Changing the lattice of a hermitian space to another unimodular lattice changes neither
the multiplicity spaces, nor the $K$-finite tensors, nor the Cartan involutions, nor the rapid decay of functions
(\lref{GlobalTheory.T.selfDualLattice_changeOfLattice_laws}). By $\fsP{STG}$ choose a core lattice $r_0$ at $\fsc_0$; the stage
$s_0=s_0(r_0,v,\mathcal M_{\fsc})$ is good of height $0$, $y\in\mathrm{CC}(s_0)$ and $x$ is $K$-finite for its model. By $\fsP{CHN}$ with
$k=m$ there are $r$, $\mathcal M$ at $\fsc_m$ and a chain with composite $\Theta$; by $\fsP{COM}$ the common data are good, so
$\fsP{OA}$ shows that $\Theta$ is surjective: $y=\Theta(y')$ with $y'\in\mathrm{CC}(s_m)$. By $\fsP{CHN}$, $y(x)=\Theta(y')(x)$ lies in
the span of the bottoms $\Psi_0$ of families $(\Psi_i)_{i\le m}$ as there, with $\Psi_m=y'(w)$. By the lattice change, $y'$ is a core form
at $\fsc_m$, so by $\fsP{TC}$, $\Psi_m=\Psi'_\xi$ for a column form $\Psi'$. For $k=m$ the conditions on the family are those of a
column chain, whose top is a scalar form of a column form; so $\Psi_0\in\fsReach(v)$.
\end{proof}

\begin{theorem}[Core forms are reachable]\label{pf:fs:coreOnto}\sloppy
$\sfname{ThetaLocalInputs}$, $\sfname{HoweUnram}$ and $\sfname{SiegelWeilUnitaryConvergent}$ imply $\fsP{CO}$.
\formallabel{GlobalTheory.T.twist_core_surjective}
\formallabel{GlobalTheory.Twist.CoreOntoAssembly.twist_core_surjective_of_local_components}
\end{theorem}

\begin{proof}\sloppy
\Cref{pf:fs:ontoAll} (with the fields of $\sfname{ThetaLocalInputs}$ listed there, $\sfname{HoweUnram}$ and
$\sfname{SiegelWeilUnitaryConvergent}$) gives $\fsP{OA}$; $\fsP{COM}$, $\fsP{STG}$ and $\fsP{TC}$ hold; $\fsP{CHN}$ holds since $\fsP{DF}$
(\cref{pf:fs:depFixed}) and $\fsP{CM}$ (\cref{pf:fs:models}) hold. Apply \cref{pf:fs:coreOntoChain}.
\end{proof}

% ---------------------------------------------------------------------------
\subsection{Ascent of periods along the tower}\label{sec:pf:fs:ascent}

This subsection proves $\fsP{AS}$ from $\fsP{DF}$, and $\fsP{DF}$ itself. Fix a setting. The core column and the first row are
compared through a \emph{tower} of periods: Fourier--Jacobi periods at the odd nodes $(2k+1,0)$ and Bessel periods at the even
nodes $(2k+2,0)$, $0\le k<m/2$. Complex conjugation is applied throughout because the formal periods are defined with the
conjugate convention; $\bar f$ denotes the complex conjugate function.

\subsubsection{Frames, periods and the rungs}

\begin{definition}[Tower frames and periods]\label{pf:fs:frames}
A \emph{tower frame} $\mathcal F$ of $G$ consists of: a hermitian line $L_z$ with an orthogonal decomposition
$V_{\fsc_0}=Y_0\oplus L_z$; line vectors $\zeta_\rho$ of $-L_z$ ($\rho\in\N$) with links $\zeta_{\rho}\leadsto\zeta_{\rho+1}$ (a matrix
carrying $\zeta_\rho$ to $-\zeta_{\rho+1}$ and multiplying the norm by $\xi_0\bar\xi_0$, $\xi_0$ the imaginary generator of the host);
and for $2\rho+2\le m$: a decomposition $V_{(2\rho+2,0)}=H^\rho\oplus V^0_\rho\oplus H^\rho$ with $\rho$ hyperbolic planes, an orthogonal
decomposition $V^0_\rho=V_{\fsc_{2\rho+2}}\oplus(-L_z)$, a decomposition $V_{(2\rho+1,0)}=H^\rho\oplus V_{\fsc_{2\rho+1}}\oplus H^\rho$, and
measures with fundamental domains on the unipotent radicals of the two flags. With the core volume $v$ this gives:
\begin{enumerate}[label=(\alph*)]
\item the \emph{Bessel period} $\fsB_\rho(\varphi,f)$ of $f$ on $G_{V_{(2\rho+2,0)}}$ against $\varphi$ on $G_{V_{\fsc_{2\rho+2}}}$: the
integral over $[G_{V_{\fsc_{2\rho+2}}}]$ and over the unipotent radical of the flag of $f$ restricted along the embedding, against $\varphi$
and the character of the unipotent radical defined by $\psi$ and $\zeta_{\rho+1}$;
\item the \emph{Fourier--Jacobi period} $\mathcal J_\rho(\varphi,f;\Phi_0)$ of $f$ on $G_{V_{(2\rho+1,0)}}$ against $\varphi$ on
$G_{V_{\fsc_{2\rho+1}}}$ and a theta function $\Phi_0$ of a Weil datum $\omega_z$ for $(V_{\fsc_{2\rho+1}},-L_z)$: the analogous integral in
which the character is replaced by the theta series of $\omega_z$ attached to $\Phi_0$ and $\zeta_\rho$.
\end{enumerate}
For a column chain $\Psi$, a row-rapid $u$ and families of lattices $\Gamma^L_\rho$ of $\fsX(V_{\fsc_{2\rho+1}},-L_z)$, define, for
$2k+2\le m$: $\mathrm{NZJ}(k)$ holds if there are a Weil datum $\omega_z$ (with characters $(\chi',\eta)$ for some $\chi'$), $\rho,\Phi$ and a
depth vacuum $\Phi_0$ for $\Gamma^L_k$ with
\[
\mathcal J_k\bigl(\overline{\Psi_{2k+1}},\overline{\Theta^{(2k+1)}_{\rho,\Phi}u};\Phi_0\bigr)\neq0;
\]
$\mathrm{NZB}(k)$ holds if there are $\rho,\Phi$ with
$\fsB_k\bigl(\overline{\Psi_{2k+2}},\overline{\Theta^{(2k+2)}_{\rho,\Phi}u}\bigr)\neq0$. The frame has \emph{good ends} if the embedding of
$Y_0$ given by $V_{\fsc_0}=Y_0\oplus L_z$ is $\iota_0$, and at the
top ($2\rho+2=m$), for rapidly decreasing $\varphi_m$ on $G_{V_{\fsc_m}}$ and $f_m$ on $G_{V_{(m,0)}}$, $\fsB_\rho(\overline{\varphi_m},\overline{f_m})\ne0$
implies $\fsB(\varphi_m,f_m)\ne0$. It is \emph{good} for $(\Gamma^L_\rho)$ if the local seesaw data of every Bessel-to-Fourier--Jacobi and
every Fourier--Jacobi-to-Bessel step are available (a family $\mathcal S$ of depth vacua for the first-row pair containing a vector with
non-zero partial Fourier transform, invariant under the relevant Heisenberg translations and under the level, with the support
conditions of the formal definition, and compatible boxes of depth vacua), the lattices $\Gamma^L_\rho$ carry non-zero vacua, and the last
seesaw at the bottom has its lattice data.
\formal{GlobalTheory.Asc.Frames}
\formal{GlobalTheory.Asc.BesL.per}
\formal{GlobalTheory.Asc.JacL.per}
\formal{GlobalTheory.Asc.NZJ}
\formal{GlobalTheory.Asc.NZB}
\formal{GlobalTheory.Asc.Frames.Ends}
\formal{GlobalTheory.Asc.Frames.Good}
\formal{GlobalTheory.Asc.Loc}
\end{definition}

\begin{definition}[Statements on the tower]\label{pf:fs:towerstatements}
A functional $\ell$ on theta functions of a pair $(S,T)$ is \emph{tame} if it is additive and homogeneous on theta functions,
continuous and bounded along Heisenberg translates, and commutes with integrals of Heisenberg translates against integrable
densities. All statements are quantified by $\forall_{\mathrm{all}}$.
\begin{enumerate}[label=(\arabic*)]
\item $\fsP{TE}$ (towers exist): for every core volume $v$ there are a tower frame with good ends and lattices $(\Gamma^L_\rho)$ for which
it is good.
\item $\fsP{FT}$ (Fock vectors are total): for $i<m$, every tame functional $\ell$ on theta functions of $\omega_{\chi(i,0)}$ (in the
conjugate orientation) which vanishes on all $\overline{\mathrm{fock}(\phi_i(\kappa),P)}$, $\kappa\in\mathcal C_i$, $P$ a polynomial,
vanishes on every depth vacuum for $\chi(i,0)$.
\item $\fsP{PF}$ (partial lifts are fixed): for all $\rho,\Phi,u$ and $i<m$, $\Theta^{(i+1)}_{\rho,\Phi}u$ is right invariant under
$K_{(i+1,0)}$.
\item $\fsP{DD}$ (row data have depth): every row datum $\phi_i(\rho,\Phi)$ is a depth vacuum for $\chi(i,0)$.
\item $\fsP{R0}$, $\fsP{RB}$, $\fsP{RJ}$ (the rungs): for every core volume $v$, every tower frame with good ends, good for
$(\Gamma^L_\rho)$, every column chain $\Psi$ and every row-rapid $u$:
$\fsP{R0}$: $\fsper_0(\Psi_0,u)\ne0$ implies $\mathrm{NZJ}(0)$;
$\fsP{RB}$: $\mathrm{NZJ}(k)$ implies $\mathrm{NZB}(k)$;
$\fsP{RJ}$: $\mathrm{NZB}(k)$ implies $\mathrm{NZJ}(k+1)$ when $2k+4\le m$.
\end{enumerate}
\formal{GlobalTheory.Asc.Tame}
\formal{GlobalTheory.Asc.TowerEx}
\formal{GlobalTheory.Asc.FockTot}
\formal{GlobalTheory.Asc.PartFixed}
\formal{GlobalTheory.SepRows.DatumDep}
\formal{GlobalTheory.Asc.Rung0}
\formal{GlobalTheory.Asc.RungB}
\formal{GlobalTheory.Asc.RungJ}
\end{definition}

$\fsP{TE}$ follows from $\fsP{DF}$ (\cref{pf:fs:towerEx}: the frames are constructed from the decompositions of the grid, and the local seesaw
data are verified place by place), and $\fsP{FT}$ holds (\cref{pf:fs:fockTot}); both are proved in \cref{sec:pf:fs:aux}.

\begin{theorem}[Ascent from the rungs]\label{pf:fs:ascentRungs}
$\fsP{TE}$, $\fsP{R0}$, $\fsP{RB}$ and $\fsP{RJ}$ imply $\fsP{AS}$.
\formallabel{GlobalTheory.Asc.ascent_of_rungs}
\formallabel{GlobalTheory.T.besselPeriod_ascent_of_tower_rungs}
\end{theorem}

\begin{proof}
Fix a setting, a core volume $v$, a column chain $\Psi$ and a row-rapid $u$ with $\fsper_0(\Psi_0,u)\ne0$. By $\fsP{TE}$ choose a frame
with good ends, good for some $(\Gamma^L_\rho)$. We show $\mathrm{NZB}(k)$ for all $k$ with $2k+2\le m$, by induction on $k$: for $k=0$,
$\fsP{R0}$ gives $\mathrm{NZJ}(0)$ and $\fsP{RB}$ gives $\mathrm{NZB}(0)$; if $2k+4\le m$ and $\mathrm{NZB}(k)$ holds, $\fsP{RJ}$ gives
$\mathrm{NZJ}(k+1)$ and $\fsP{RB}$ gives $\mathrm{NZB}(k+1)$. Write $m=2r'$ with $r'\ge1$ and take $k=r'-1$, so $2k+2=m$: there are
$\rho,\Phi$ with $\fsB_k(\overline{\Psi_m},\overline{\Theta^{(m)}_{\rho,\Phi}u})\ne0$. As $\Psi_m$ and $\Theta^{(m)}_{\rho,\Phi}u=\Theta_{\rho,\Phi}u$
are rapidly decreasing ($\Psi$ is a column chain, $u$ is row-rapid), the good top end gives $\fsB(\Psi_m,\Theta_{\rho,\Phi}u)\ne0$.
\end{proof}

\subsubsection{The rungs}

The three rungs have the same structure: a seesaw identity transfers the non-vanishing of one period to the non-vanishing of the next
period of a theta lift by \emph{some} depth vacuum; the period of the lift is a tame functional of the Schwartz function; by $\fsP{FT}$ it
is non-zero on some conjugated Fock vector over a coset vector; and changing the row datum at the relevant index realises this Fock
vector inside the row lift. We use the following facts about the grid, all proved in the development.

\begin{enumerate}[label=(T\arabic*)]
\item\label{pf:fs:T1} (conjugation; \lref{GlobalTheory.T.gridConjugation_weilData_inverseCharacter}) For $i<m$: the conjugated
Weil representations $\bar\omega_{\chi(i,0)}$ (on $(V_{(i+1,0)},V_{(i,0)})$) and $\bar\omega_{\fsc(i)}$ (on $(V_{\fsc_{i+1}},V_{\fsc_i})$) are
Weil data, with characters $(\chi_{m-i}^{-1}\eta,\chi_{m-i}^{-1})$ for odd $i$ and $(\chi_{m-i}^{-1},\chi_{m-i}^{-1}\eta)$ for even $i$, and
$(\chi_{m-i}^{-1},\chi_{m-i}^{-1})$ respectively; $f$ is a depth vacuum for $\fsc(i)$ if and only if $\bar f$ is a vacuum of the
conjugated depth lattice; $\overline{\fsdn_i(f,F)}$ is the lift of $\bar F$ by $\bar f$ for $\bar\omega_{\fsc(i)}$; the conjugate of
$\Theta^{(i+1)}_{\rho,\Phi}u$ is the lift of $\overline{\Theta^{(i)}_{\rho,\Phi}u}$ by $\overline{\phi_i(\rho,\Phi)}$ for
$\bar\omega_{\chi(i,0)}$; $\Theta^{(i)}_{\rho,\Phi}u$ depends only on $(\rho_j,\Phi_j)_{j<i}$; and $(\nu_{(i,0)},\fsF_{(i,0)})$ is a Haar
measure with a fundamental domain.
\item\label{pf:fs:T2} (tameness; \lref{GlobalTheory.T.thetaLift_periods_isTameThetaFunctional}) For a rapidly decreasing $G$,
the Bessel periods and the Fourier--Jacobi periods (against a theta function $\Phi_0$) of the theta lifts of $G$ by a theta function $\Phi$
are tame functionals of $\Phi$.
\item\label{pf:fs:T3} (seesaw transfers) Fourier--Jacobi to Bessel (\cref{pf:fs:upBF}), Bessel to Fourier--Jacobi (\cref{pf:fs:upFB}) and the
last seesaw at the bottom (\lref{GlobalTheory.T.lastSeesaw_periodNonvanishing}): in the local situation provided by a good frame, if the period of the
lower level of the seesaw is non-zero for the lift of the core function by a box of depth vacua, then the period at the next level of the
theta lift of the row function by some depth vacuum $\Phi'$ from the family $\mathcal S$ of the frame is non-zero.
\item\label{pf:fs:T4} (links; \lref{GlobalTheory.T.linkedLineVectors_invertible_on_theta_functions}) Every theta function for
$(V_{\fsc_{2\rho+1}},-L_z)$ is of the form $w\mapsto\overline{\Xi(\kappa_\rho w)}$ for a theta function $\Xi$ of the linked pair.
\end{enumerate}

\begin{proposition}[Fourier--Jacobi to Bessel]\label{pf:fs:rungB}
$\fsP{FT}$ and $\fsP{PF}$ imply $\fsP{RB}$.
\formallabel{GlobalTheory.AscRungB.rungB}
\formallabel{GlobalTheory.T.fourierJacobi_to_bessel_nonvanishing_tower}
\end{proposition}

\begin{proof}
Assume $\mathrm{NZJ}(k)$: there are $\omega_z$, $\rho,\Phi$ and a depth vacuum $\Phi_0$ for $\Gamma^L_k$ with
$\mathcal J_k(\overline{\Psi_{2k+1}},\overline{\Theta^{(2k+1)}_{\rho,\Phi}u};\Phi_0)\ne0$. Put $i=2k+1<m$. As $\Psi$ is a column chain,
$\Psi_i=\fsdn_i(f,\Psi_{i+1})$ for a depth vacuum $f$, so by \ref{pf:fs:T1} $\overline{\Psi_i}$ is the lift of $\varphi=\overline{\Psi_{i+1}}$ by
$\bar f$, and $\bar f$ is a vacuum of the conjugated depth lattice. By \ref{pf:fs:T4}, $\Phi_0=\overline{\Xi_2(\kappa_k\,\cdot)}$ for a theta
function $\Xi_2$, and the box condition of the good frame shows that the box formed from $\bar f$ and $\Xi_2$ is a depth vacuum of the
lattice of the frame. Let $G'=\overline{\Theta^{(i)}_{\rho,\Phi}u}$: it is rapidly decreasing ($u$ is row-rapid), and right invariant under
$K_{(i,0)}$ by $\fsP{PF}$. The step of the frame at level $k$ (Bessel above, Fourier--Jacobi below), with the Weil datum
$\bar\omega_{\chi(i,0)}$ and the characters of \ref{pf:fs:T1}, satisfies the hypotheses of the Fourier--Jacobi-to-Bessel transfer
\ref{pf:fs:T3}; it gives $\Phi_1\in\mathcal S$ with $\fsB_k(\varphi,\Theta(G',\Phi_1))\ne0$, where $\Theta(G',\Phi)$ is the lift of $G'$ by
$\Phi$ for $\bar\omega_{\chi(i,0)}$ along $(i,0)\to(i+1,0)$.

Let $\ell(\Phi)=\fsB_k(\varphi,\Theta(G',\Phi))$; it is tame by \ref{pf:fs:T2}, and $\ell(\Phi_1)\ne0$ with $\Phi_1$ a depth vacuum for
$\chi(i,0)$ (elements of $\mathcal S$ are). Suppose that $\mathrm{NZB}(k)$ fails. For $\kappa\in\mathcal C_i$ and a polynomial $P$ let
$\rho'$ and $\Phi'$ be $\rho$ and $\Phi$ with the $i$-th entries replaced by $P$ and $\kappa$. By \ref{pf:fs:T1},
$\Theta^{(i)}_{\rho',\Phi'}u=\Theta^{(i)}_{\rho,\Phi}u$ and $\phi_i(\rho',\Phi')=\mathrm{fock}(\phi_i(\kappa),P)$, so
$\overline{\Theta^{(i+1)}_{\rho',\Phi'}u}=\Theta\bigl(G',\overline{\mathrm{fock}(\phi_i(\kappa),P)}\bigr)$, and the failure of $\mathrm{NZB}(k)$ gives
$\ell\bigl(\overline{\mathrm{fock}(\phi_i(\kappa),P)}\bigr)=0$. By $\fsP{FT}$, $\ell$ vanishes on all depth vacua, in particular
$\ell(\Phi_1)=0$, a contradiction.
\end{proof}

\begin{proposition}[Bessel to Fourier--Jacobi]\label{pf:fs:rungJ}
$\fsP{FT}$ and $\fsP{PF}$ imply $\fsP{RJ}$.
\formallabel{GlobalTheory.Asc.RungJac.rungJ}
\formallabel{GlobalTheory.T.bessel_to_fourierJacobi_nonvanishing_tower}
\end{proposition}

\begin{proof}
Assume $\mathrm{NZB}(k)$ with $\rho,\Phi$, and $2k+4\le m$; put $i=2k+2<m$ and $a=\chi_{m-i}$. By \ref{pf:fs:T1} ($i$ is even) the
characters of $\bar\omega_{\chi(i,0)}$ and $\bar\omega_{\fsc(i)}$ are $(a^{-1},a^{-1}\eta)$ and $(a^{-1},a^{-1})$; hence $\eta$ is a Hecke
character with the restriction required for the line, and there is a Weil datum $\omega_z$ for $(V_{\fsc_{i+1}},-L_z)$ with characters
$(a^{-1},\eta)$ (\lref{GlobalTheory.WeilDatumExistence.weilDatumExists}). The good frame
provides the local data of the step at level $k$ (Fourier--Jacobi above, Bessel below) and a non-zero vacuum $\Phi_0$ for $\Gamma^L_{k+1}$.
Write $\Psi_i=\fsdn_i(f,\Psi_{i+1})$; by \ref{pf:fs:T1} the Bessel period of $\mathrm{NZB}(k)$ is
$\fsB_k(\mathrm{lift}(\overline{\Psi_{i+1}},\bar f),G')$ with $G'=\overline{\Theta^{(i)}_{\rho,\Phi}u}$, which is rapidly decreasing and
$K_{(i,0)}$-invariant ($\fsP{PF}$). The Bessel-to-Fourier--Jacobi transfer \ref{pf:fs:T3} gives $\Phi'\in\mathcal S$ with
$\mathcal J_{k+1}(\overline{\Psi_{i+1}},\Theta(G',\Phi');\Phi_0)\ne0$. The functional $\ell(\Phi)=\mathcal J_{k+1}(\overline{\Psi_{i+1}},\Theta(G',\Phi);\Phi_0)$
is tame \ref{pf:fs:T2} and non-zero at the depth vacuum $\Phi'$, so by $\fsP{FT}$ there are $\kappa,P$ with
$\ell(\overline{\mathrm{fock}(\phi_i(\kappa),P)})\ne0$. With $\rho',\Phi'$ as in the proof of \cref{pf:fs:rungB},
$\overline{\Theta^{(i+1)}_{\rho',\Phi'}u}=\Theta(G',\overline{\mathrm{fock}(\phi_i(\kappa),P)})$, which gives $\mathrm{NZJ}(k+1)$ with
$(\omega_z,\rho',\Phi',\Phi_0)$.
\end{proof}

\begin{proposition}[The bottom rung]\label{pf:fs:rung0}
$\fsP{FT}$ implies $\fsP{R0}$.
\formallabel{GlobalTheory.AscBottom.rung0}
\formallabel{GlobalTheory.T.nonvanishing_thetaTower_bottom_rung}
\end{proposition}

\begin{proof}
Let $\fsper_0(\Psi_0,u)\ne0$ and write $\Psi_0=\fsdn_0(f,\Psi_1)$. Let $a=\chi_m$. By \ref{pf:fs:T1} (with $i=0$, which is even),
$\bar\omega_{\chi(0,0)}$ has characters $(a^{-1},a^{-1}\eta)$ and $\bar\omega_{\fsc(0)}$ has characters $(a^{-1},a^{-1})$; hence $\eta$ is a
Hecke character, its restriction to $\A_F^\times$ has the parity required for the line $-L_z$ (from $\dim V_{\fsc_0}=\dim Y_0+1$), and there is
a Weil datum $\omega_z$ for $(V_{\fsc_1},-L_z)$ with characters $(a,\eta)$ (\lref{GlobalTheory.WeilDatumExistence.weilDatumExists}).
Since the bottom end of the frame is $\iota_0$, the conjugate of $\fsper_0(\Psi_0,u)$ is the integral over $\fsF_{(0,0)}$ of the lift of
$\overline{\Psi_1}$ by $\bar f$, restricted to $Y_0$, against $\bar u$; it is non-zero. The last seesaw \ref{pf:fs:T3}, with the lattice data
of the good frame, gives a depth vacuum $\Phi'$ for $\chi(0,0)$ and a vacuum $\Phi_z$ for $\Gamma^L_0$ with
$\mathcal J_0(\overline{\Psi_1},\Theta(\bar u,\Phi');\Phi_z)\ne0$. The functional $\Phi\mapsto\mathcal J_0(\overline{\Psi_1},\Theta(\bar u,\Phi);\Phi_z)$ is
tame \ref{pf:fs:T2}, so by $\fsP{FT}$ there are $\kappa,P$ with a non-zero value at $\overline{\mathrm{fock}(\phi_0(\kappa),P)}$. With
$\rho=(P,0,\dots)$ and $\Phi=(\kappa,\ast,\dots)$, \ref{pf:fs:T1} gives
$\overline{\Theta^{(1)}_{\rho,\Phi}u}=\Theta(\bar u,\overline{\mathrm{fock}(\phi_0(\kappa),P)})$, which is $\mathrm{NZJ}(0)$.
\end{proof}

\subsubsection{Depth}

\begin{proposition}[Depth vacua are fixed]\label{pf:fs:depFixed}
$\fsP{DF}$ holds.
\formallabel{GlobalTheory.DepthFix.depFixed}
\formallabel{GlobalTheory.T.dependence_fixed}
\end{proposition}

\begin{proof}
Let $p=\chi(i,0)$ or $\fsc(i)$ with $i<m$. The two lattices of the pair are self-dual (the lattice of $S_p$ is the self-dual lattice of
the negative of a node in range), and both splitting characters of $\omega_p$ satisfy the depth bound of the grid
(\lref{GlobalTheory.DepthFix.depthLe_att}, deduced from the conductor conditions on $\chi_{m-i}$ in the grid datum and on $\eta$
in $L$). The depth lattice $\Gamma_p$ contains every element of the lattice of $\mathrm{Hom}$ whose components at the places $w(v)$ lie in
$\fsm\,\mathrm{Hom}$. In this situation a theta function fixed by the Heisenberg translations by $\Gamma_p$ is fixed by $\omega_p(k,k')$ for
$k,k'$ in the principal congruence subgroups modulo $\fsm$ (\lref{GlobalTheory.DepthFix.fix_pair}). Its proof is local, place by
place: at the split places $w(v)$, $v\in S_0$, the Weil representation acts by the splitting characters and by translations, the
translations $u\mapsto k'uk^{-1}-u$ lie in the depth lattice, and the splitting characters are trivial on the determinants of the
congruence elements by the depth bound; at the other places, a lattice vacuum is fixed by the stabilisers of the self-dual
lattices.
\end{proof}

\begin{proposition}[Partial lifts are fixed]\label{pf:fs:partFixed}
$\fsP{DD}$ holds, and $\fsP{DF}$ and $\fsP{DD}$ imply $\fsP{PF}$.
\formallabel{GlobalTheory.FirstRowDepth.datumDep}
\formallabel{GlobalTheory.T.firstRow_data_have_depth}
\formallabel{GlobalTheory.Asc.partFixed_of_depFixed}
\formallabel{GlobalTheory.T.gridLiftsFixedByCongruenceDepth_of_depthData}
\end{proposition}

\begin{proof}
$\fsP{DD}$: every coset vector $\phi_i(\kappa)$ is a non-zero depth vacuum and a Fock vacuum at every $\tau$
(\lref{GlobalTheory.T.coset_vectors_nonzero_depth_fock_vacua}). Fock vectors $\mathrm{fock}(\phi,P)$ over a theta function $\phi$ are
theta functions, and the Heisenberg translation by an element $u$ with trivial archimedean component commutes with the creation
operators, which act at the archimedean places (\lref{GlobalTheory.Hk3.FiniteHeis.heisO_pureTensorO_of_finite}); hence
$\mathrm{fock}(\phi_i(\kappa),P)$ is fixed by the translations by $\Gamma_{\chi(i,0)}$ (whose elements have trivial archimedean components),
in both conditions of the depth lattice.

$\fsP{PF}$: let $k\in K_{(i+1,0)}$. By definition,
$\Theta^{(i+1)}_{\rho,\Phi}u(gk)=\int_{\fsF_{(i,0)}}\theta_{\omega}((gk)^-,x;\phi)\,\Theta^{(i)}_{\rho,\Phi}u(x)\,d\nu(x)$ with $\omega=\omega_{\chi(i,0)}$ and
$\phi=\phi_i(\rho,\Phi)$, and $((gk)^-,x)=(g^-,x)(k^-,1)$. As $\phi$ is a depth vacuum ($\fsP{DD}$), $\fsP{DF}$ with $k'=k$ and the trivial
element in the small group gives $\omega(k^-,1)\phi=\phi$, so $\theta_\omega((gk)^-,x;\phi)=\theta_\omega(g^-,x;\phi)$ and the lift is
right $K_{(i+1,0)}$-invariant.
\end{proof}

\begin{theorem}[Ascent]\label{pf:fs:ascent}
$\fsP{DF}$ implies $\fsP{AS}$.
\formallabel{GlobalTheory.Asc.ascent_of_depFixed}
\formallabel{GlobalTheory.T.bessel_period_nonvanishing_ascent_of_depth_data}
\end{theorem}

\begin{proof}
$\fsP{FT}$ holds (\cref{pf:fs:fockTot}), and $\fsP{PF}$ follows from $\fsP{DF}$ and $\fsP{DD}$ (\cref{pf:fs:partFixed}). Hence $\fsP{R0}$, $\fsP{RB}$ and
$\fsP{RJ}$ hold (\cref{pf:fs:rung0,pf:fs:rungB,pf:fs:rungJ}), $\fsP{TE}$ follows from $\fsP{DF}$ (\cref{pf:fs:towerEx}), and \cref{pf:fs:ascentRungs} gives
$\fsP{AS}$.
\end{proof}

% ---------------------------------------------------------------------------
\subsection{The Bessel gap and the top separation}\label{sec:pf:fs:bur}

\subsubsection{Uniqueness of Bessel forms}

The statement $\fsP{BG}$ is deduced from a uniqueness theorem for Bessel maps proved elsewhere in the development. For integers
$k\ge1$, $r'\ge0$ and $n=2r'+k+1$, let $H_{k,r'}$ be the standard hermitian matrix of size $n$ with $r'$ hyperbolic planes on the
first and last $r'$ coordinates and the diagonal block $\mathrm{diag}(1,\dots,1,-1)$ of size $k+1$ in the middle; let $N_{k,r'}$ be the
standard unipotent radical with the generic character $\psi_{c}$ attached to $c\in(\C^\times)^{r'}$ and a vector $z$ of the middle block
with $\langle z,z\rangle<0$, and let $\Uni(k)$ act on the positive part of the middle block. A \emph{Bessel map} for a representation
$\sigma$ of $\Uni(H_{k,r'})$ and a Harish-Chandra type model $\pi$ of $\Uni(k)$ on a Hilbert space $W$ is a map $B\colon H\to W$, linear on
smooth vectors, equivariant for the stabiliser of the positive part, $(N_{k,r'},\psi_c)$-equivariant, and continuous for a Sobolev norm in
the sense of Casselman--Wallach.

\begin{littheorem}[{Uniqueness of Bessel maps in large gap; \lref{CompanionTheory.Bmu.besselMapUniqGap}}]\label{pf:fs:bmu}
Let $\sigma$ be a discrete series of $\Uni(H_{k,r'})$ with Harish-Chandra parameter $\lambda$ whose gap is at least
$h\ge120n^3$, let $z$ be as above, $c\in(\C^\times)^{r'}$ and $\pi$ a Harish-Chandra type model of an
irreducible representation of $\Uni(k)$. Then any two smooth Bessel maps $B_1,B_2$ are linearly dependent on the smooth vectors: there are
$(a,b)\ne(0,0)$ with $aB_1(v)+bB_2(v)=0$ for all smooth $v$.
\end{littheorem}

This is a special case of the archimedean uniqueness of Bessel models \cite{JiangSunZhu2010,SunZhu2012} (multiplicity at most one
for the Gan--Gross--Prasad pairs of unitary groups), restricted to discrete series with a large gap between the Harish-Chandra parameters
of opposite signs; the formal proof (about $5\,000$ lines, shared with other parts of the development) is independent of these
references.

\begin{proposition}[The Bessel gap]\label{pf:fs:burGapPos}
$\fsP{BG}$ holds.
\formallabel{UpperTheory.Lit.burGapPos}
\formallabel{UpperTheory.Lit.Bur.BurPos.burGapPos_of_bmu}
\end{proposition}

\begin{proof}
Let the data of $\fsP{BG}$ be given; by the frame conditions $n=2r'+m'+1$. If $W=0$ there is nothing to prove. Otherwise:

\emph{Normalisation.} There are invertible matrices $Q$ and $A$ such that $Q$ carries the standard form $H_{m',r'}$ to $J$, the standard
isotropic family to $e$ and a multiple $t\,e_{\mathrm{last}}$ ($t>0$) of a standard vector to $z$, and $A^*J_0A=1$ with $FA$ the standard
embedding of the positive block (\lref{UpperTheory.Lit.Bur.BurPos.frame_normal}, which uses $\langle z,z\rangle_J<0$ and the
positivity of $J_0$). Put $\sigma'(g)=\sigma(QgQ^{-1})$; then $\sigma'$ is a discrete series of $\Uni(H_{m',r'})$ with gap at least $h$
(\lref{UpperTheory.Lit.Bur.BurPos.isDSHpk_tr}), and $v$ is smooth for $\sigma'$ if and only if it is smooth for $\sigma$
(\lref{UpperTheory.Lit.Bur.smoothVec_conjugate_unitaryGroup_iff}).

\emph{The compact side.} The contragredient of $h\mapsto\pi_0(AhA^{-1})$ is a continuous irreducible representation of $\Uni(m')$
(\lref{UpperTheory.Lit.Bur.contragredient_conj_definiteUnitary_isContIrred}); it is unitarizable
(\lref{UpperTheory.Lit.Bur.contIrred_compactUnitaryGroup_rep_isUnitarizable}) and isomorphic to a Harish-Chandra type model $\rho'$
(\cref{pf:fs:hcType}). Let $T\colon W^*\to W_q$ be the resulting isomorphism, so
that $T((\pi_0(AhA^{-1})^{-1})^*\ell)=\rho'(h)T\ell$ for $h\in\Uni(m')$.

\emph{From Bessel forms to Bessel maps.} For a Bessel form $\beta$ let $B(x)=T(\beta(\cdot,x))$ on smooth vectors (extended by $0$).
Then $B$ is a smooth Bessel map for $\sigma'$, $\rho'$ and the generic character $\psi_c$ attached to the constants $c$ of $\chi$:
linearity is that of $\beta$; equivariance under the stabiliser follows from the invariance property of $\beta$, since an element of the
stabiliser of the positive block conjugated by $Q$ extends the corresponding element $AsA^{-1}$ of $\Uni(J_0)$
(\lref{UpperTheory.Lit.Bur.blockEmbed_conj_isUnitaryEmbedding}); the $(N,\psi_c)$-equivariance follows from the equivariance of
$\beta$, since $Q$ maps the standard unipotent radical into $N(J,e)$
(\lref{UpperTheory.Lit.Bur.isometry_conj_maps_unipotentRadical} and \lref{UpperTheory.Lit.Bur.mem_besselUnipotentRadical_of_mem_coordGroup})
and the characters correspond (\lref{UpperTheory.Lit.Bur.BurPos.chi_eq}); and the Sobolev continuity follows from that of the
finitely many linear forms $\beta(b_i^*,\cdot)$, $b_i$ a basis of $W$
(\lref{UpperTheory.Lit.Bur.casselmanWallachCont_finite_sum_sobolev_forms} and \lref{UpperTheory.Lit.Bur.sobolevCont_of_conj_gram_form}).

\emph{Conclusion.} Given Bessel forms $\beta_1,\beta_2$, \cref{pf:fs:bmu} applied to the associated Bessel maps gives $(a,b)\ne(0,0)$ with
$aB_1+bB_2=0$ on smooth vectors; since $T$ is injective, $a\beta_1+b\beta_2=0$. Hence any two Bessel forms are linearly dependent, and the
space of Bessel forms has dimension at most $1$.
\end{proof}

\subsubsection{The coefficient space and the top separation}

\begin{proposition}[Tensor decomposition of the dual coefficient space]\label{pf:fs:coefTensor}
$\fsP{CT}$ holds.
\formallabel{GlobalTheory.CoefTensorProof.coefTensor}
\formallabel{GlobalTheory.T.dualCoefSpace_isTensorOverRealPlaces}
\end{proposition}

\begin{proof}
Let $c$ be a coefficient datum of $V$ with its standing properties (in particular $W^c\ne0$). For $\tau\in\fsT$ and $g\in\Uni(V^\tau)$ let
$a_\tau(g)$ act on $(W^c)^*$ by $\ell\mapsto\ell\circ\rho^c(\iota_\tau(g))^{-1}$. The actions at distinct indices commute, because the
images $\iota_\tau(g)$ and $\iota_{\tau'}(h)$ commute in $G_V$ (\lref{GlobalTheory.PointAtPlaceComm.atTau_comm}), and $(W^c)^*$ is
irreducible under all of them together (\lref{GlobalTheory.T.dual_irreducible_coefficientSpace_realPlaces}). A finite-dimensional
space which is irreducible under a finite family of pairwise commuting group actions is isomorphic to the tensor product of irreducible
representations $W_\tau$ of the individual groups, compatibly with the actions
(\lref{GlobalTheory.T.exists_tensorDecomposition_commutingFamilies}, a form of the theorem that irreducible representations of a
product of groups are tensor products). Let $\pi_\tau(g)$ be the action of $a_\tau(g)$ transported to $W_\tau$ (and the identity for
$g\notin\Uni(V^\tau)$). Then $\pi_\tau$ is a representation ($a_\tau$ is multiplicative because $\iota_\tau$ is, and $\iota_\tau(1)=1$), it
is irreducible, and its matrix coefficients are continuous on $\Uni(V^\tau)$, being matrix coefficients of the contragredient of $\rho^c$
(\lref{GlobalTheory.T.contragredient_matrix_coefficient_continuousOn_unitary}). The compatibility of the isomorphism $e$ with the
actions is the equivariance asserted in $\fsP{CT}$.
\end{proof}

\begin{proposition}[Top separation]\label{pf:fs:topSep}
$\fsP{TS}$ holds (and the formal statement takes $\sfname{ThetaLocalInputs}$ as an unused hypothesis).
\formallabel{GlobalTheory.SepRows.SepProof.theta_lifts_separate_of_separating_data}
\formallabel{GlobalTheory.T.theta_lifts_separate_points}
\end{proposition}

\begin{proof}
Fix a setting, a model $\mathcal M_B$ of $V_{(m,0)}$, $F'\in\fsA(\mathcal M_B;\emptyset,1)$ and $\Psi\in\fsCols$. The values $F'(v)$ are
continuous left $\Gamma$-invariant functions (\lref{GlobalTheory.Core.Mi.SepPer.mult_mem_contLeft} for pure tensors, and by
linearity in general), and so are the scalar forms $\Psi_\xi$ (\lref{GlobalTheory.RowBplus.pairVec_cont} and \lref{GlobalTheory.RowBplus.pairVec_left});
$\xi\mapsto\Psi_\xi$ is linear. There is a compact set $\Omega$ meeting every rational orbit relevant to the Bessel pairing
(\lref{GlobalTheory.T.exists_compact_meets_besselTopRationalOrbits}), and for each scalar form $\Psi_\xi$ the Bessel pairing
$F\mapsto\fsB(\Psi_\xi,F)$ is additive in $F$ on continuous left-invariant functions, additive in the first argument, and bounded by
$C_\xi\sup_\Omega|F|$ (\lref{GlobalTheory.T.besselPeriodCorankOne_additive_bounded_compact}); it is also homogeneous in each
argument. Hence (i): $(\xi,v)\mapsto\fsB(\Psi_\xi,F'(v))$ is bilinear.

For (ii), fix $\tau$, all coordinates of $x$ and $v$ except those at $\tau$, and let $\beta_\tau(y,u)$ be the value at the updated
coordinates. It is bilinear by (i) and the multilinearity of $\otimes$ and of $e$. \emph{Continuity:} by the uniform bound for
multiplicity maps on compact sets (\lref{GlobalTheory.Core.Mi.SepPer.mult_bound}: there are a Sobolev form $L_w$ and a constant with
$|F'(\ldots,u,\ldots)(g)|\le C\,L_w(u)$ for $g\in\Omega$), the linear form $u\mapsto\beta_\tau(y,u)$ is bounded by a constant times a Sobolev
norm. \emph{Invariance:} if $g\in\Uni(V_{(m,0)}^\tau)$ extends $h\in\Uni(V_{\fst_0}^\tau)$, then
$\iota_\tau(g)=\iota_{\mathrm{top}}(\iota_\tau(h))$ (\lref{GlobalTheory.T.embRel_atRealPlace_of_unitaryEmbedding}); replacing $y$ by
$\pi_\tau(h)y$ replaces $\Psi_{e(\otimes x)}$ by its right translate by $\iota_\tau(h)$ (by the equivariance of $e$,
\lref{GlobalTheory.SepRows.SepProof.pairVec_comp_rho}), and replacing $u$ by $\pi_\tau(g)u$ replaces $F'(\ldots)$ by its right translate by
$\iota_\tau(g)$ (equivariance of multiplicity maps); simultaneous right translation by $\iota_\tau(h)$ and $\iota_{\mathrm{top}}(\iota_\tau(h))$
does not change the Bessel pairing (\lref{GlobalTheory.T.besselPeriodCorankOne_simultaneous_rightTranslate}).
\emph{Equivariance:} for $u_0$ in the unipotent radical at $\tau$, $\iota_\tau(u_0)$ lies in the adelic unipotent radical
(\lref{GlobalTheory.T.atRealPlace_mem_unipotentFlag_adelic}), and right translation of $F$ by such an element multiplies
$\fsB(\varphi,F)$ by the value of the character (\lref{GlobalTheory.T.besselPeriodCorankOne_rightTranslate_unipotent}). Hence
$\beta_\tau$ is a Bessel form at $\tau$.
\end{proof}

\subsubsection{The gap at the top of the row}

\begin{lemma}[Gap at the top]\label{pf:fs:gap}
Let $D$ carry size data with $\sfname{SizeRows}$, $L$ with $L.\mathrm{Is}(\mu_0)$, $G$ a grid datum with $c_{\mathrm{gap}}=1$ and
$\mathcal M_B$ a model of $V_{(m,0)}$ with parameter $S^{\mathrm{row}}_m$. Then $\mathcal M_B$ has large gap: at every $\tau$, $\lambda_\tau$ is
a Harish-Chandra parameter and the gap of $(\eps_\tau,\lambda_\tau)$ is at least $120(\dim V_{(m,0)})^3$.
\formallabel{UpperTheory.Lit.Bur.Gap.fit_gap}
\end{lemma}

\begin{proof}
That $\lambda_\tau$ is a Harish-Chandra parameter is the integrality statement used in \cref{pf:fs:models}. For the gap: $\dim V_{(m,0)}=r+m(r+6)$.
Under $\sfname{SizeRows}$, consecutive nodes of the Harish-Chandra data of $D$ satisfy $b_\tau(j+1)-b_\tau(j)>120(\dim V_{(m,0)})^3$
(\lref{UpperTheory.Lit.Bur.Gap.lowsOver_of_sizeRows}: an explicit arithmetic consequence of the inequalities of the parameter choice
and of $\sfname{SizeRows}$, through the quantity $\ell^*=mb/\mathrm{Den}$ of the size data). The inequality propagates to every pair of
entries of opposite signs of the grid parameter at $(m,0)$ (\lref{UpperTheory.Lit.Bur.Gap.gapOver_parD}), and is preserved by the
shift defining the row parameter
(\lref{UpperTheory.Lit.Bur.Gap.gapOver_shift}); for the enumeration of the model it gives the gap condition
(\lref{UpperTheory.Lit.Bur.Gap.gapAtLeast_of_gapOver}).
\end{proof}

% ---------------------------------------------------------------------------
\subsection{The archimedean input}\label{sec:pf:fs:real}

This subsection proves \cref{pf:fs:RIholds}. Notation is as in \cref{pf:fs:RI}: $\pi$ is a discrete series of $\Uni(J_\eps)$
($n+1$ variables) with parameter $\lambda$, $\sigma$ one of $\Uni(J_{\eps'})$ ($n$ variables) with parameter $\mu$, the data satisfy
$\mathrm{Par}(e,\eps,\lambda,\eps',\mu,B)$, $dh$ is the fixed Haar measure of $\Uni(J_{\eps'})$, $\mathcal U_\delta=\{h:\|h-1\|\le\delta\}$,
and $d_\sigma$ is the formal degree of $\sigma$ for $dh$. A \emph{branching embedding} is a continuous linear map $j\colon H'\to H$
with $j\sigma(h)=\pi(\tilde h)j$, $\|jw\|=\|w\|$, together with its adjoint $P$.

\begin{definition}[Two branching statements]\label{pf:fs:branching}
\begin{enumerate}[label=(\arabic*)]
\item $\fsP{PL}$ (large projections): for all data as above and every lowest $K$-type unit vector $v_0$ of $\pi$ there are a branching
embedding $(j,P)$ and $g_0\in\Uni(J_\eps)$ with $\|g_0\|\le(8(n+1)B)^{10(n+1)}$ and $1\le\mathrm{big}(100,n,B)\,\|P\pi(g_0)v_0\|$.
\item $\fsP{EM}$ (eigen maps): for all data as above, every formal degree $d$ of $\sigma$, every unit $w_0$ with integrable coefficient, every
branching embedding $(j,P)$ and every $v\in H$ with integrable restricted coefficient $h\mapsto\langle v,\pi(\tilde h)v\rangle$ and $Pv\ne0$,
there are $w\ne0$, $c\ge\|Pv\|^2/(2d)$ and an intertwiner $T\colon H\to H'$ for the restriction ($T\pi(\tilde h)=\sigma(h)T$) with $Tv=w$ and
$\int\overline{\langle v,\pi(\tilde h)v\rangle}\langle b,\sigma(h)w\rangle\,dh=c\langle b,w\rangle$ for all $b$.
\end{enumerate}
\formal{UpperTheory.Lit.Branch.IsEmb}
\formal{UpperTheory.Lit.Branch.ProjLarge}
\formal{UpperTheory.Lit.Branch.EigMap}
\end{definition}

\begin{theorem}[The archimedean input from branching]\label{pf:fs:realInputE}
$\fsP{PL}$ and $\fsP{EM}$ imply $\fsP{RI}$.
\formallabel{UpperTheory.Lit.Branch.realInputE}
\end{theorem}

\begin{proof}
We use the following estimates for discrete series, proved in the development.
\begin{enumerate}[label=(E\arabic*)]
\item\label{pf:fs:E1} (integral bounds; \lref{GlobalTheory.T.translateCoefficientIntegralBounds_of_unitary_laws}) There is $C_1>0$
such that for every lowest $K$-type unit $v_0$ and every $g_0$ with $\|g_0\|\le(8(n+1)B)^{10(n+1)}$, the vector $v=\pi(g_0)v_0$ satisfies:
for all $g$, $h\mapsto\langle v,\pi(g\tilde h)v\rangle$ is integrable with
$\int|\langle v,\pi(g\tilde h)v\rangle|\,dh\le\vol(\mathcal U_{1/4})\,\mathrm{big}(C_1,n,B)\,\mathrm{col}(g)^{-(2n+7)}$, and
$|\langle v,\pi(g)v\rangle|\le A\|g\|^{-20(n+1)^2}$ for some $A$. (It rests on the decay of lowest $K$-type coefficients, a volume bound for
balls and a boost decomposition for the norms of the embedded group.)
\item\label{pf:fs:E2} (derivative bounds; \lref{GlobalTheory.T.translateCoefficientDerivativeIntegralBounds_of_laws}) There is $C_3>0$
with the same hypotheses such that for $X$ in the Lie algebra of $\Uni(J_{\eps'})$ with $\|X\|\le1$,
$\int|\langle d\pi(\tilde X)v,\pi(\tilde h)v\rangle|\,dh\le\vol(\mathcal U_{1/4})\,\mathrm{big}(C_3,n,B)$, the integrand being integrable.
\item\label{pf:fs:E3} (volume and degree; \lref{UpperTheory.Lit.Branch.volDeg}) There is $C_2>0$ with
$\vol(\mathcal U_{1/4})\le\mathrm{big}(C_2,n,B)/d$ for every formal degree $d$ of $\sigma$.
\item\label{pf:fs:E4} (smoothness; \lref{UpperTheory.Lit.Branch.mapSmooth}, \lref{GlobalTheory.T.lowestKType_lieWord_normBound})
Lowest $K$-type vectors are smooth, and every intertwiner $T$ for the restriction maps smooth vectors of $\pi$ to smooth vectors of $\sigma$.
\item\label{pf:fs:E5} (degrees and lowest types; \lref{GlobalTheory.T.formalDegree_loweringWeight_ratio_le},
\lref{GlobalTheory.T.discreteSeries_lowestKType_unit_unitary_formalDegree}) $\sigma$ has a formal degree $d'$; $\pi$ and $\sigma$ have
lowest $K$-type unit vectors, and the coefficient of a lowest $K$-type unit vector of $\sigma$ is integrable (its decay exponent exceeds the
dimension because of the gap).
\end{enumerate}
Put $C=C_1+C_3+C_2+201$. Given the data, choose a lowest $K$-type unit vector $v_0$ of $\pi$. By $\fsP{PL}$ there are $(j,P)$ and $g_0$ with
$\|g_0\|\le(8(n+1)B)^{10(n+1)}$ and $1\le\mathrm{big}(100,n,B)\|Pv\|$, $v=\pi(g_0)v_0$; in particular $Pv\ne0$. By \ref{pf:fs:E1} with $g=1$, the
restricted coefficient of $v$ is integrable. Let $w_0$ be a lowest $K$-type unit vector of $\sigma$ (\ref{pf:fs:E5}). By $\fsP{EM}$ (with $d'$ and
$w_0$) there are $w\ne0$, $c\ge\|Pv\|^2/(2d')>0$ and $T$ with $Tv=w$ and
$\int\overline{\langle v,\pi(\tilde h)v\rangle}\langle b,\sigma(h)w\rangle\,dh=c\langle b,w\rangle$. By \ref{pf:fs:E4}, $v$ is smooth (a translate
of a smooth vector) and $w=Tv$ is smooth. Put $w'=w/\|w\|$, a smooth unit vector. Then
\[
\int\overline{\langle v,\pi(\tilde h)v\rangle}\,\langle b,\sigma(h)w'\rangle\,dh=\|w\|^{-1}c\,\langle b,w\rangle=c\,\langle w',w'\rangle\langle b,w'\rangle,
\]
which is (i). By \ref{pf:fs:E3}, $\vol(\mathcal U_{1/4})\le\mathrm{big}(C_2,n,B)/d'\le\mathrm{big}(C_2,n,B)\cdot2c/\|Pv\|^2
\le c\,\mathrm{big}(C_2+201,n,B)$, using $\|Pv\|^{-2}\le\mathrm{big}(200,n,B)$, $2\le\mathrm{big}(1,n,B)$ (as $B\ge1$) and the multiplicativity
$\mathrm{big}(C,n,B)\mathrm{big}(C',n,B)=\mathrm{big}(C+C',n,B)$. Inserting this into \ref{pf:fs:E1} and \ref{pf:fs:E2} and using the monotonicity
of $\mathrm{big}$ in $C$ gives (ii) and (iii) with the constant $C$; (iv) is the second part of \ref{pf:fs:E1}.
\end{proof}

\begin{proof}[Proof of \cref{pf:fs:RIholds}]
$\fsP{EM}$ holds (\lref{UpperTheory.Lit.Branch.eigMap_of}). It follows from two general statements about unitary representations
$\sigma,\rho$ of a unimodular locally compact group with Haar measure $dg$, $\sigma$ of formal degree $d$: if $w_0$ is a unit vector of
$\sigma$ with integrable coefficient, $i$ an isometric intertwiner from $\sigma$ to $\rho$ and $v$ a vector of $\rho$ with integrable
coefficient, then there are $w\ne0$ and $c\ge\|i^*v\|^2/(2d)$ with
$\int\overline{\langle v,\rho(g)v\rangle}\langle b,\sigma(g)w\rangle\,dg=c\langle b,w\rangle$ for all $b$
(\lref{GlobalTheory.CompanionTheory.RgHeA.eigA}); and if this relation holds with $c\ne0$, then $w=Tv$ for an intertwiner $T$ from
$\rho$ to $\sigma$ (\lref{GlobalTheory.CompanionTheory.RgHeA.TMapPf.tMap}). $\fsP{PL}$ holds
(\lref{UpperTheory.Lit.Branch.projLarge_of}). It follows from the existence of a Schur eigenprojection (an isometric branching
embedding $i$ and a smooth vector $\eta$ of $\pi$ with $\langle\eta,iw\rangle\ne0$ for which the orthogonality relation
$\int\overline{\langle w,\sigma(h)w\rangle}\langle\eta,\pi(\tilde h)a\rangle\,dh=\kappa\langle\eta,iw\rangle\langle iw,a\rangle$ holds,
$w$ a lowest $K$-type unit vector of $\sigma$), which is deduced from the occurrence of $\sigma$ in the restriction of $\pi$ under the
interlacing condition (\cref{pf:fs:occ}) and a statement on lowest $K$-type models
(\lref{PositivityTheory.RgHe.LowestKTypeModel.seed}, \lref{PositivityTheory.RgHe.A.Pf.seedB_of}); and from a lower bound for the
projection of a translate of a lowest $K$-type vector (\lref{PositivityTheory.RgHe.A.Pf.bigProjP}), which rests on the moment estimate
(\cref{pf:fs:moment}) and the analytic estimates for orbits (\cref{pf:fs:analytic}). \Cref{pf:fs:realInputE} then gives $\fsP{RI}$. The
occurrence theorem, the moment estimate and the analytic estimates are proved in \cref{sec:pf:fs:aux}; the remaining deductions (of the
eigenprojections, of the lower bound and of $\fsP{EM}$) are not reproduced here.
\end{proof}

The archimedean input is a statement about branching from $\Uni(p+1,q)$ to $\Uni(p,q)$ for discrete series with interlacing parameters,
in the spirit of the restriction theory of holomorphic and discrete series representations \cite{JakobsenVergne1979,Knapp1986}; we do not
know a reference for the explicit bounds (ii)--(iv).

% ---------------------------------------------------------------------------
\subsection{Proofs of auxiliary results}\label{sec:pf:fs:aux}

This subsection proves fifteen results which were used above through references to the formal development: the existence of
good tower frames and the totality of Fock vectors (used in \cref{sec:pf:fs:ascent}), the two seesaw transfers of
\ref{pf:fs:T3}, the existence of chains along the core column (used in \cref{sec:pf:fs:onto}), three statements on Poincar\'e
series, the grid parameters and the existence of packs (used in \cref{sec:pf:fs:periods}), and four statements on
representations of unitary groups (used in \cref{sec:pf:fs:bur,sec:pf:fs:real}). Deeper ingredients are again cited with their
statements.

\subsubsection{Vacua and tame functionals}

\begin{definition}[Vacua and annihilators]\label{pf:fs:vacua}
Let $(S,T)$ be a pair of hermitian spaces of $D$, $\psi$ a standard additive character and $\Gamma$ a subgroup of
$\fsX(S,T)$. For $w\in\fsX(S,T)$ let $\mathrm{heis}(w)$ be the Heisenberg translation by $w$ on functions on $\fsX(S,T)$, and
$\langle\,,\rangle$ the symplectic form of $\fsX(S,T)$.
\begin{enumerate}[label=(\alph*)]
\item A \emph{$\Gamma$-vacuum} is a theta function $f$ for $(S,T)$ with $\mathrm{heis}(u)f=f$ for all $u\in\Gamma$. (The depth vacua of
\cref{pf:fs:chain} are the $\Gamma_p$-vacua.)
\item The \emph{annihilator} $\mathrm{Ann}(\Gamma)$ is the set of $w$ with $\psi(\langle w,2u\rangle)=1$ for all $u\in\Gamma$ and
$\psi(\langle w,u\rangle)=1$ for every $u\in\Gamma$ which is a half element for a split idempotent of $\A_E$ (in the sense of the
formal definition).
\end{enumerate}
\formal{GlobalTheory.Asc.VacG}
\formal{GlobalTheory.Asc.Ann}
\end{definition}

The following irreducibility statement is the basic tool for the tower; it expresses that the $\Gamma$-vacua form an irreducible
module under the translations by the annihilator, in a form adapted to tame functionals (\cref{pf:fs:towerstatements}).

\begin{lemma}[Irreducibility of vacua; \lref{GlobalTheory.T.tame_functional_vacuum_irreducibility}]\label{pf:fs:irr}
Let $\psi$ be standard, $\ell$ a tame functional on theta functions of $(S,T)$ and $f_0$ a non-zero $\Gamma$-vacuum such that
$\ell(\mathrm{heis}(w)f_0)=0$ for all $w\in\mathrm{Ann}(\Gamma)$. Then $\ell(f)=0$ for every $\Gamma$-vacuum $f$.
\end{lemma}

\begin{proposition}[Fock vectors are total]\label{pf:fs:fockTot}
$\fsP{FT}$ holds.
\formallabel{GlobalTheory.T.fockConjugates_span_depth_data}
\formallabel{GlobalTheory.ConjDataTotal.fockTot}
\end{proposition}

\begin{proof}
Fix a setting, $i<m$ and a tame functional $\ell$ on theta functions of the pair $(V_{(i+1,0)},V_{(i,0)})$ which vanishes on
$\overline{\mathrm{fock}(\phi_i(\kappa),P)}$ for all $\kappa\in\mathcal C_i$ and all polynomials $P$. Write $\Gamma=\Gamma_{\chi(i,0)}$.
We use the following properties of the coset vectors (\lref{GlobalTheory.T.heisenbergCommutant_permutes_firstRowCosetVectors}):
for every $\kappa\in\mathcal C_i$,
\begin{enumerate}[label=(\alph*)]
\item $\phi_i(\kappa)\ne0$, and $\phi_i(\kappa)$ is a Fock vacuum at every $\tau$ for the complex structure defined by the Cartan
involutions of $W$ at the nodes $(i,0)$ and $(i+1,0)$;
\item $\overline{\phi_i(\kappa)}$ is a $\Gamma$-vacuum;
\item for every $w\in\mathrm{Ann}(\Gamma)$ with trivial archimedean component, $\mathrm{heis}(w)\overline{\phi_i(\kappa)}$ lies in the span of
the functions $\overline{\phi_i(\kappa')}$, $\kappa'\in\mathcal C_i$.
\end{enumerate}

\emph{Step 1: archimedean translates.} Let $\kappa\in\mathcal C_i$ and let $w$ have trivial finite component. The functional
$\bar\ell(f)=\overline{\ell(\bar f)}$ is again tame. By (a), the Fock vectors over $\phi_i(\kappa)$ for the frames of the Cartan
involutions of $W$ are the vectors $\mathrm{fock}(\phi_i(\kappa),P)$, and $\bar\ell$ vanishes on all of them by hypothesis. A tame
functional which vanishes on all Fock vectors over a Fock vacuum vanishes on all archimedean Heisenberg translates of this vacuum
(\lref{GlobalTheory.T.tameFunctional_kills_fockVacuum_realTranslates}). Hence $\bar\ell(\mathrm{heis}(w)\phi_i(\kappa))=0$, and since
complex conjugation commutes with Heisenberg translations for a standard $\psi$, $\ell(\mathrm{heis}(w)\overline{\phi_i(\kappa)})=0$.

\emph{Step 2: all translates of one coset vector.} Fix $\kappa_0\in\mathcal C_i$ and let $w\in\mathrm{Ann}(\Gamma)$. Write
$w=w_\infty+w_f$ with $w_\infty$ archimedean and $w_f$ finite. The elements of $\Gamma$ have trivial archimedean components, so
$\langle w,u\rangle=\langle w_f,u\rangle$ for $u\in\Gamma$ and $w_f\in\mathrm{Ann}(\Gamma)$. The symplectic pairing of an archimedean and a
finite element vanishes, so $\mathrm{heis}(w)=\mathrm{heis}(w_\infty)\mathrm{heis}(w_f)$. By (c), applied to $w_f$ (which has trivial archimedean
component),
$\mathrm{heis}(w_f)\overline{\phi_i(\kappa_0)}=\sum_{\kappa'}c_{\kappa'}\overline{\phi_i(\kappa')}$ for some $c_{\kappa'}\in\C$, hence
\[
\ell\bigl(\mathrm{heis}(w)\overline{\phi_i(\kappa_0)}\bigr)=\sum_{\kappa'}c_{\kappa'}\,\ell\bigl(\mathrm{heis}(w_\infty)\overline{\phi_i(\kappa')}\bigr)=0
\]
by Step 1, using the additivity and homogeneity of $\ell$ on theta functions (Heisenberg translates of theta functions are theta
functions).

\emph{Step 3.} By (a) and (b), $f_0=\overline{\phi_i(\kappa_0)}$ is a non-zero $\Gamma$-vacuum, and by Step 2, $\ell$ vanishes on
$\mathrm{heis}(w)f_0$ for all $w\in\mathrm{Ann}(\Gamma)$. By \cref{pf:fs:irr}, $\ell$ vanishes on every $\Gamma$-vacuum, that is on every depth
vacuum for $\chi(i,0)$.
\end{proof}

\subsubsection{Seesaw transfers}

We now prove the two transfers of \ref{pf:fs:T3}. They concern general Bessel and Fourier--Jacobi data, of which those of a tower frame
(\cref{pf:fs:frames}) are special cases.

\begin{definition}[Bessel and Fourier--Jacobi data, seesaw steps]\label{pf:fs:seesaw}
Let $\rho\in\N$.
\begin{enumerate}[label=(\alph*)]
\item A \emph{Bessel datum} $\mathcal B$ at level $\rho$ consists of a decomposition $V=H^\rho\oplus V_0\oplus H^\rho$ with $\rho$ hyperbolic planes,
an orthogonal decomposition $V_0=A\oplus L$ with a line $L$ and a line vector $\zeta$ of $L$, a Haar measure with a fundamental domain on
$G_A$ and a measure with a fundamental domain $\fsF_N$ on the unipotent radical $N$ of the flag of the isotropic family. Its \emph{Bessel
period} of $f$ on $G_V$ against $\varphi$ on $G_A$ is
\[
\fsB_{\mathcal B}(\varphi,f)=\int_{\fsF_A}\int_{\fsF_N}\varphi(h)\,f(n\,\iota_A(h))\,\overline{\psi_\zeta(n)}\,dn\,dh,
\]
where $\iota_A$ is the embedding of $G_A$ and $\psi_\zeta$ the character of $N$ attached to $\psi$ and $\zeta$.
\item A \emph{Fourier--Jacobi datum} $\mathcal J$ at level $\rho$ consists of a decomposition $W=H^\rho\oplus W_0\oplus H^\rho$, a Haar
measure with a fundamental domain on $G_{W_0}$ and a measure with a fundamental domain on the unipotent radical $N'$ of the flag. For a
line $L'$ with a line vector $\zeta'$ and a Weil datum $\omega_z$ for $(W_0,L')$, its \emph{Fourier--Jacobi period} of $f$ on $G_W$
against $\varphi$ on $G_{W_0}$ and a theta function $\Phi_0$ of $(W_0,L')$ is
\[
\mathcal J_{\mathcal J}(\varphi,f;\Phi_0)=\int_{\fsF_{W_0}}\int_{\fsF_{N'}}\varphi(h)\,f(n\,\iota_0(h))\,\overline{\vartheta(\Phi_0)(n,h)}\,
\overline{\psi_{N'}(n)}\,dn\,dh,
\]
where $\vartheta(\Phi_0)$ is the Jacobi theta function of $\Phi_0$ for $\omega_z$ and $\zeta'$, and $\psi_{N'}$ the character of the flag.
\item A \emph{seesaw step} $s$ from a pair of spaces $(X\supseteq X_0)$ to a pair $(Y\supseteq Y_0)$, both given with hyperbolic
decompositions, consists of Weil data $\omega$ for $(Y,X)$ with splitting characters $(\chi_Y,\chi_X)$ and $\omega_0$ for $(Y_0,X_0)$,
measures and fundamental domains, and the resulting theta lift $\mathrm{lift}_s(G,\Phi)$ of a function $G$ on $G_X$ to $G_Y$ by a theta
function $\Phi$ of $(Y,X)$. For a Bessel datum $\mathcal B$ and a Fourier--Jacobi datum $\mathcal J$ at the same level, a
\emph{Bessel-to-Fourier--Jacobi step} is a seesaw step from the spaces $(W\supseteq W_0)$ of $\mathcal J$ to the spaces $(V\supseteq V_0)$ of
$\mathcal B$; for a Fourier--Jacobi datum $\mathcal J$ at level $\sigma+1$ and a Bessel datum $\mathcal B$ at level $\sigma$, a
\emph{Fourier--Jacobi-to-Bessel step} is a seesaw step from the spaces of $\mathcal B$ to those of $\mathcal J$. In both cases the flat map
$\Phi\mapsto\Phi^\flat$ of the step (\cref{pf:fs:loc}) is formed with the matrices prescribed by the type of the step (built from the
isotropic families and the Bessel vector of $\zeta$).
\item For an orthogonal decomposition $V_0=A\oplus L$ and functions $\Xi_A$ on $\fsX(A,T)$, $\Xi_L$ on $\fsX(L,T)$, the \emph{box}
$\Xi_A\boxtimes\Xi_L$ is the function $w\mapsto\Xi_A(w_A)\Xi_L(w_L)$ on $\fsX(V_0,T)$, $(w_A,w_L)$ the components of $w$; the box
on $\fsX(S,V_0)$ of functions on $\fsX(S,A)$ and $\fsX(S,L)$ is defined in the same way.
\end{enumerate}
\formal{GlobalTheory.Asc.BesL}
\formal{GlobalTheory.Asc.JacL}
\formal{GlobalTheory.Asc.BesL.per}
\formal{GlobalTheory.Asc.JacL.per}
\formal{GlobalTheory.Asc.StBF}
\formal{GlobalTheory.Asc.StFB}
\formal{GlobalTheory.MixModel.OrthDec.boxS}
\formal{GlobalTheory.MixModel.OrthDec.boxT}
\end{definition}

For a seesaw step $s$ from $(X\supseteq X_0)$ to $(Y\supseteq Y_0)$, with $\rho$ the number of hyperbolic planes of $X$, the formal development
defines a \emph{flat map} $\Phi\mapsto\Phi^\flat$ from theta functions of $(Y,X)$ to functions on $\fsX(Y_0,X_0)$: a partial Fourier transform
of $\Phi$ in the mixed model of the step, formed with two matrices $x$ (with $\rho$ columns, in $X$) and $y$ attached to the step. Let $R_x$ be
the group of $g\in G_X$ with $gx=xu$ for some upper unitriangular $u\in\GL_\rho(\A_E)$.

\begin{definition}[Local seesaw data]\label{pf:fs:loc}
Let $s$ be a seesaw step from $(X\supseteq X_0)$ to $(Y\supseteq Y_0)$ with flat map $\Phi\mapsto\Phi^\flat$, $\mathcal S$ a set of functions on
$\fsX(Y,X)$, $\Gamma_0$ a subgroup of $\fsX(Y_0,X_0)$ and $K$ a subgroup of $G_X$. Then $(\mathcal S,\Gamma_0,K)$ are \emph{local seesaw data}
for $s$ if: every $\Phi\in\mathcal S$ is a theta function; $\mathcal S$ is stable under the Heisenberg translations by the elements with trivial
finite component and by the images in $\fsX(Y,X)$ of the elements of $\mathrm{Ann}(\Gamma_0)$; $\Phi^\flat$ is a $\Gamma_0$-vacuum for every
$\Phi\in\mathcal S$; $K$ is a compact open subgroup of the finite part of $G_X$ and $\omega(1,k)\Phi=\Phi$ for $k\in K$ and $\Phi\in\mathcal S$; if
$(\omega(1,h)\Phi)^\flat\ne0$ for some $\Phi\in\mathcal S$, then $h\in R_x\,G_{X,\infty}\,K$; and $\Phi^\flat\ne0$ for some $\Phi\in\mathcal S$.
\formal{GlobalTheory.Asc.Loc}
\end{definition}

In both transfers, the theta lift of $\varphi$ by a theta function $\Xi$ for a Weil datum $\omega'$ is denoted $\mathrm{lift}_{\omega'}(\varphi,\Xi)$, and
a \emph{link} from $(L,\zeta)$ to a line $(L',\zeta')$ is a matrix $\tau$ with $\tau\zeta'=-\zeta$ and
$\langle\zeta,\zeta\rangle=\xi_0\bar\xi_0\langle\zeta',\zeta'\rangle$; it induces a map $\kappa$ from $\fsX(W_0,L')$ to $\fsX(L,W_0)$
(\ref{pf:fs:T4} is the statement that this map is invertible on theta functions for the links of a tower frame).

\begin{proposition}[Fourier--Jacobi to Bessel]\label{pf:fs:upBF}
Let $\mathcal B$ and $\mathcal J$ be a Bessel and a Fourier--Jacobi datum at level $\rho$, with $V_0=A\oplus L$, $s$ a
Bessel-to-Fourier--Jacobi step between them and $(\mathcal S,\Gamma_0,K)$ local seesaw data for $s$. Let $(L',\zeta')$ be a line with a
link from $(L,\zeta)$, inducing $\kappa$; let $\chi_A,\chi_z,\chi'$ be characters with $\chi_V=\chi_A\chi_z$, $\omega_a$ a Weil datum for
$(A,W_0)$ with characters $(\chi_A,\chi_W)$ and $\omega_z$ a Weil datum for $(W_0,L')$ with characters $(\chi',\chi_z)$. Let $\varphi$ on
$G_A$ and $G$ on $G_W$ be rapidly decreasing, $G$ right invariant under $K$, and let $\Xi_1,\Xi_2$ be theta functions of $(A,W_0)$ and
$(L,W_0)$ such that the box $\Xi_1\boxtimes\Xi_2$ is a $\Gamma_0$-vacuum. If
\[
\mathcal J_{\mathcal J}\bigl(\mathrm{lift}_{\omega_a}(\varphi,\Xi_1),G;\,w\mapsto\overline{\Xi_2(\kappa w)}\bigr)\neq0,
\]
then $\fsB_{\mathcal B}(\varphi,\mathrm{lift}_s(G,\Phi))\ne0$ for some $\Phi\in\mathcal S$.
\formallabel{GlobalTheory.T.fourierJacobi_period_gives_bessel_period_of_lift}
\formallabel{GlobalTheory.UpBessel.upBF}
\end{proposition}

\begin{proof}
Suppose that $\fsB_{\mathcal B}(\varphi,\mathrm{lift}_s(G,\Phi))=0$ for all $\Phi\in\mathcal S$. For a theta function $\Xi$ of $(V_0,W_0)$ let
$\ell(\Xi)=\mathcal J^{\mathrm{all}}(\varphi,G;\Xi)$ be the Fourier--Jacobi period of $G$ formed with the Fourier--Jacobi kernel of $\Xi$ for
$\omega_0$ and integrated against $\varphi$ over $[G_A]$ (the period over all data of the step).

(1) \emph{Descent of vanishing.} The seesaw identity of the step shows that $\ell(\Phi^\flat)=0$ for all $\Phi\in\mathcal S$
(\lref{GlobalTheory.T.vanishing_descends_bessel_to_fourierJacobi}: for local seesaw data and rapidly decreasing, $K$-invariant $G$, the
vanishing of all Bessel periods $\fsB_{\mathcal B}(\varphi,\mathrm{lift}_s(G,\Phi))$, $\Phi\in\mathcal S$, implies the vanishing of
$\mathcal J^{\mathrm{all}}(\varphi,G;\Phi^\flat)$ for all $\Phi\in\mathcal S$).

(2) \emph{Tameness.} $\ell$ is a tame functional on theta functions of $(V_0,W_0)$
(\lref{GlobalTheory.T.fourierJacobiPeriod_all_data_isTameThetaFunctional}).

(3) \emph{Irreducibility.} Choose $\Phi_1\in\mathcal S$ with $\Phi_1^\flat\ne0$; $\Phi_1^\flat$ is a $\Gamma_0$-vacuum. Let $w\in\mathrm{Ann}(\Gamma_0)$
and let $\hat w$ be its image in $\fsX(V,W)$. The flat map intertwines the Heisenberg translations: $(\mathrm{heis}(\hat w)\Phi_1)^\flat=
\mathrm{heis}(w)\Phi_1^\flat$ (\lref{GlobalTheory.FlatHeis.flatAt_heisO}). As $\mathrm{heis}(\hat w)\Phi_1\in\mathcal S$, (1) gives
$\ell(\mathrm{heis}(w)\Phi_1^\flat)=0$. By \cref{pf:fs:irr}, $\ell$ vanishes on all $\Gamma_0$-vacua, in particular $\ell(\Xi_1\boxtimes\Xi_2)=0$.

(4) \emph{Product formula.} For the box of $\Xi_1$ and $\Xi_2$,
\[
\ell(\Xi_1\boxtimes\Xi_2)=\mathcal J_{\mathcal J}\bigl(\mathrm{lift}_{\omega_a}(\varphi,\Xi_1),G;\,w\mapsto\overline{\Xi_2(\kappa w)}\bigr)
\]
(\lref{GlobalTheory.T.periodAllData_products_eq_fourierJacobiPeriod}: under the hypotheses on the characters and Weil data of the
proposition, the period over all data of a box $\Xi_1\boxtimes\Xi_2$ of theta functions is this Fourier--Jacobi period). This contradicts the
hypothesis.
\end{proof}

\begin{proposition}[Bessel to Fourier--Jacobi]\label{pf:fs:upFB}
Let $\mathcal J$ be a Fourier--Jacobi datum at level $\sigma+1$ and $\mathcal B$ a Bessel datum at level $\sigma$, with $V_0=A\oplus L$ and line
vector $\zeta$, $s$ a Fourier--Jacobi-to-Bessel step from the spaces $V\supseteq V_0$ of $\mathcal B$ to the spaces $W\supseteq W_0$ of $\mathcal J$,
with characters $\chi_W$ (attached to $W$) and $\chi_V$ (attached to $V$), and $(\mathcal S,\Gamma_0,K)$ local seesaw data for $s$. Let
$\chi_A,\chi_z$ be characters with $\chi_V=\chi_A\chi_z$, $\omega_b$ a Weil datum for $(W_0,A)$ with characters $(\chi_W,\chi_A)$ and $\omega_z$
a Weil datum for $(W_0,L)$ with characters $(\chi_W,\chi_z)$. Let $\varphi$ on $G_{W_0}$ and $G$ on $G_V$ be rapidly decreasing, $G$ right
invariant under $K$, $\Xi_1$ a theta function of $(W_0,A)$ and $\Phi_0\ne0$ a theta function of $(W_0,L)$ such that the box
$\Xi_1\boxtimes\Phi_0$ is a $\Gamma_0$-vacuum. If $\fsB_{\mathcal B}(\mathrm{lift}_{\omega_b}(\varphi,\Xi_1),G)\ne0$, then
$\mathcal J_{\mathcal J}(\varphi,\mathrm{lift}_s(G,\Phi);\Phi_0)\ne0$ for some $\Phi\in\mathcal S$, the Fourier--Jacobi period being formed with
$(L,\zeta)$ and $\omega_z$.
\formallabel{GlobalTheory.T.bessel_period_gives_fourierJacobi_period_of_lift}
\formallabel{GlobalTheory.UpFromBessel.upFB}
\end{proposition}

\begin{proof}
Suppose that $\mathcal J_{\mathcal J}(\varphi,\mathrm{lift}_s(G,\Phi);\Phi_0)=0$ for all $\Phi\in\mathcal S$. Choose a measure with a fundamental
domain on $\fsX(W_0,L)$ (\lref{GlobalTheory.Unf.VolExist.measures_exist}). For a theta function $\Xi$ of $(W_0,V_0)$ let
$\mathrm{hd}(\Xi)$ be the \emph{handed-down function} on $G_A$,
\[
\mathrm{hd}(\Xi)(h_1)=\int_{\fsF_{W_0}}\varphi(h_0)\,c\bigl(\omega_0(h_0,\iota_A(h_1))\Xi,\ \omega_z(h_0,1)\Phi_0\bigr)(0)\,dh_0,
\]
where $c$ is the contraction of a function on $\fsX(W_0,V_0)$ with a function on $\fsX(W_0,L)$ along the $L$-component (an integral over the
chosen fundamental domain), and let $\ell(\Xi)=\fsB_{\mathcal B}(\mathrm{hd}(\Xi),G)$.

(1) \emph{Descent of vanishing.} $\ell(\Phi^\flat)=0$ for all $\Phi\in\mathcal S$
(\lref{GlobalTheory.T.vanishing_descends_fourierJacobi_to_bessel}: for local seesaw data, rapidly decreasing $\varphi$, rapidly decreasing
$K$-invariant $G$ and a theta function $\Phi_0$, the vanishing of the Fourier--Jacobi periods of all lifts $\mathrm{lift}_s(G,\Phi)$,
$\Phi\in\mathcal S$, implies the vanishing of the Bessel periods of $G$ against the handed-down functions of all $\Phi^\flat$).

(2) \emph{Tameness.} $\ell$ is tame on theta functions of $(W_0,V_0)$
(\lref{GlobalTheory.T.besselPeriod_handedDown_isTameThetaFunctional}).

(3) \emph{Irreducibility.} As in the proof of \cref{pf:fs:upBF}, with $\Phi\in\mathcal S$ such that $\Phi^\flat\ne0$, the intertwining of the
flat map with the core Heisenberg translations (\lref{GlobalTheory.FlatHeis.flatAt_heisO}) and the stability of $\mathcal S$ under them
show that $\ell$ vanishes on $\mathrm{heis}(w)\Phi^\flat$ for $w\in\mathrm{Ann}(\Gamma_0)$, so by \cref{pf:fs:irr} it vanishes on the
$\Gamma_0$-vacuum $\Xi_1\boxtimes\Phi_0$.

(4) \emph{Product formula.} The handed-down function of a box is a multiple of a theta lift:
\[
\mathrm{hd}(\Xi_1\boxtimes\Phi_0)=\langle\Phi_0,\Phi_0\rangle\,\mathrm{lift}_{\omega_b}(\varphi,\Xi_1)
\]
(\lref{GlobalTheory.T.besselHandedFunction_product_eq_coreLift_mul_pairing}: under the hypotheses on the characters and Weil data of the
proposition, the handed-down function, formed with $\Phi_0$, of a box of $\Xi_1$ with a theta function of $(W_0,L)$ is the theta pairing of
that function with $\Phi_0$ times $\mathrm{lift}_{\omega_b}(\varphi,\Xi_1)$; here we only use the box with $\Phi_0$ itself), where
$\langle\Phi_0,\Phi_0\rangle=\int|\Phi_0|^2$ over the chosen fundamental domain. Hence
$0=\ell(\Xi_1\boxtimes\Phi_0)=\langle\Phi_0,\Phi_0\rangle\,\fsB_{\mathcal B}(\mathrm{lift}_{\omega_b}(\varphi,\Xi_1),G)$. Since $\Phi_0$ is a non-zero
theta function, $\langle\Phi_0,\Phi_0\rangle\ne0$: theta functions are continuous and bounded, so $|\Phi_0|^2$ is integrable on the fundamental
domain, and if its integral vanished, $\Phi_0$ would vanish almost everywhere on the fundamental domain, hence (as $|\Phi_0|$ is invariant
under the translations by rational points) almost everywhere, hence everywhere by continuity. So
$\fsB_{\mathcal B}(\mathrm{lift}_{\omega_b}(\varphi,\Xi_1),G)=0$, a contradiction.
\end{proof}

\subsubsection{Tower frames}

Fix a setting. By \cref{pf:fs:grid}(d), for $1\le i\le m$ the grid embeds $V_{\fsc_i}$ into $V_{(i,0)}$ with a complement of dimension
$i-1$ whose Gram matrix in a given basis $(a_0,\dots,a_{i-2})$ is antidiagonal of sign $-1$ (the \emph{column decomposition} at $i$);
for $i=2\rho+2$ the middle vector $a_\rho$ is anisotropic. Recall that $z_\ell$ is the basis vector of $\ell$ in $V_{\fsc_0}=Y_0\oplus\ell$.

\begin{definition}[Standard frames]\label{pf:fs:std}
A tower frame $\mathcal F$ (\cref{pf:fs:frames}) is \emph{standard} if:
\begin{enumerate}[label=(\alph*)]
\item the embedding of $Y_0$ into $V_{\fsc_0}$ given by $V_{\fsc_0}=Y_0\oplus L_z$ has the matrix of the grid decomposition
$V_{\fsc_0}=Y_0\oplus\ell$, and the embedding of $L_z$ maps $\zeta_0$ to $z_\ell$;
\item $\zeta_\rho=(-\xi_0)^\rho\zeta_0$ for all $\rho$;
\item for $2\rho+2\le m$, the two isotropic families of the decomposition $V_{(2\rho+2,0)}=H^\rho\oplus V^0_\rho\oplus H^\rho$ are non-zero
multiples of $a_0,\dots,a_{\rho-1}$ and of $a_{2\rho},\dots,a_{\rho+1}$ (the first ones being $(-\xi_0)^{j+1}a_j$), the composite embedding
of $V_{\fsc_{2\rho+2}}$ into $V_{(2\rho+2,0)}$ is that of the grid, and the composite embedding of $-L_z$ maps $\zeta_0$ to $a_\rho$;
\item for $2\rho+2\le m$, the decomposition $V_{(2\rho+1,0)}=H^\rho\oplus V_{\fsc_{2\rho+1}}\oplus H^\rho$ has the same relation to the column
decomposition at $2\rho+1$, and its embedding of $V_{\fsc_{2\rho+1}}$ is that of the grid.
\end{enumerate}
The \emph{line lattices} of $\mathcal F$ are the following subgroups $\Gamma^L_\rho$ of $\fsX(V_{\fsc_{2\rho+1}},-L_z)$: $\Gamma^L_0$ is the
set of $w$ such that the element of $\fsX(V_{\fsc_1},V_{\fsc_0})$ with components $0$ along $Y_0$ and $w$ along $L_z$ lies in the core depth
lattice $\Gamma_{\fsc(0)}$; for $\rho\ge1$, $\Gamma^L_\rho$ is the set of $w$ carrying the lattice $\Lambda_{\fsc_{2\rho+1}}$ into the closure of
$\mathfrak a_\rho\zeta_0$, where $\mathfrak a_\rho$ is the ideal of the host generated by an explicit element (a product of powers of the
generators of the type primes of $D$ and of a shift element of the grid).
\formal{GlobalTheory.AscFin.Std}
\formal{GlobalTheory.AscFin.lineLat}
\end{definition}

\begin{lemma}[Standard frames exist]\label{pf:fs:stdEx}
In every setting there is a standard tower frame.
\formallabel{GlobalTheory.T.exists_standardAuxiliaryFrames_allRowGrids}
\formallabel{GlobalTheory.AscFin.exists_std_frames}
\end{lemma}

\begin{proof}
The tower frame built from the decompositions of the grid (with $L_z=\ell$, the hyperbolic parts spanned by the multiples of the
complement vectors prescribed in \cref{pf:fs:std}(c),(d), and $V^0_\rho$ the orthogonal sum of $V_{\fsc_{2\rho+2}}$ and the line of $a_\rho$)
is standard (\lref{GlobalTheory.AscFin.auxiliaryFrames_standard}, a verification on the matrices of the decompositions).
\end{proof}

\begin{lemma}[Standard frames have good ends]\label{pf:fs:ends}
Every standard tower frame has good ends, for every core volume $v$.
\formallabel{GlobalTheory.T.standardAuxiliaryFrames_fit_grid_ends}
\formallabel{GlobalTheory.AscFin.EndsPf.ends}
\end{lemma}

\begin{proof}
\emph{Bottom.} By \cref{pf:fs:std}(a) the embeddings of $Y_0$ given by the frame and by the grid have the same matrix, so they coincide:
the embedding of $Y_0$ given by $V_{\fsc_0}=Y_0\oplus L_z$ is $\iota_0$.

\emph{Top.} Let $2\rho+2=m$. By \cref{pf:fs:std}(b),(c) the line vector of the Bessel period $\fsB_\rho$ is
$\zeta_{\rho+1}=(-\xi_0)^{\rho+1}\zeta_0$, which the frame maps to $(-\xi_0)^{\rho+1}a_\rho$; the hyperbolic part is attached to the column
decomposition at $m$ and the embedding of $V_{\fsc_m}=V_{\fst_0}$ is that of the grid. For such data the non-vanishing of the frame
Bessel period of $(\overline{\varphi_m},\overline{f_m})$, for rapidly decreasing $\varphi_m$ on $G_{V_{\fsc_m}}$ and $f_m$ on $G_{V_{(m,0)}}$,
implies $\fsB(\varphi_m,f_m)\ne0$ (\lref{GlobalTheory.AscFin.EndsPf.top_gen}: under these conditions the unipotent radical, the embedding of
$G_{V_{\fsc_m}}$ and the character of the frame period coincide with those of the Bessel datum of the grid (\cref{pf:fs:grid}); after the change
of measures, which multiplies the period by a constant, the frame Bessel period of $(\overline{\varphi_m},\overline{f_m})$ is the complex
conjugate of a multiple of $\fsB(\varphi_m,f_m)$).
\end{proof}

The verification that standard frames are good is local, step by step. For a seesaw step $s$ from $(X\supseteq X_0)$ to $(Y\supseteq Y_0)$
(\cref{pf:fs:seesaw}) and a lattice $\Gamma$ of $\fsX(Y,X)$, the formal development defines families $p$ of \emph{admissible pieces} of the step
(local data at the places of $S_0$, which may \emph{split} $\Gamma$), and for each such family a \emph{supported class} $\mathcal S(p)$ of theta
functions of $(Y,X)$, built from the mixed models of the step, and a lattice $\Gamma_0(p)$ of $\fsX(Y_0,X_0)$. It also defines
\emph{unfolding conditions} on the elements of $G_X$.
\begin{lemma}[Supported classes give local data]\label{pf:fs:locCls}
Let a seesaw step be given with $\psi$ standard, $x$ the image of a rational matrix $x_0$ with linearly independent columns, admissible
pieces $p$ splitting a lattice $\Gamma$ of $\fsX(Y,X)$, and a compact open subgroup $K$ of the finite part of $G_X$. Assume that
$\omega(1,k)f=f$ for every $\Gamma$-vacuum $f$ and $k\in K$, and that every $h$ such that $h$ and all $h\kappa$, $\kappa\in K$, satisfy the
unfolding conditions of the step lies in $R_x\,G_{X,\infty}\,K$. Then $(\mathcal S(p),\Gamma_0(p),K)$ are local seesaw data, and every element
of $\mathcal S(p)$ is a $\Gamma$-vacuum.
\formallabel{GlobalTheory.AscFin.GoodStd.loc_cls}
\end{lemma}

\begin{proof}
The elements of $\mathcal S(p)$ are theta functions by definition, and they are $\Gamma$-vacua
(\lref{GlobalTheory.T.doubleMixedModelSupportedClass_subset_splitLatticeVacua}: the supported class lies in the vacua of a lattice split
by the pieces). The class is stable under the archimedean and the core Heisenberg translations, and the flats of its elements are
$\Gamma_0(p)$-vacua (\lref{GlobalTheory.T.doubleMixedModelClass_stable_flat_vacua}). The invariance under $K$ is the hypothesis, applied to
the $\Gamma$-vacua of the class. Some element has a non-zero flat
(\lref{GlobalTheory.T.exists_unfolding_flatMap_ne_zero_fullRank}, which uses that $x_0$ has full rank). For the support condition, let
$\Phi\in\mathcal S(p)$ with $(\omega(1,h)\Phi)^\flat\ne0$. Then $h$ satisfies the unfolding conditions
(\lref{GlobalTheory.T.unfoldingConditions_of_flatMap_weilAction_ne_zero}: the non-vanishing of the flat of a Weil translate of an element
of the class implies the unfolding conditions). For $\kappa\in K$, $\omega(1,h\kappa)\Phi=\omega(1,h)\omega(1,\kappa)\Phi=\omega(1,h)\Phi$, so
$h\kappa$ satisfies them as well; by hypothesis $h\in R_x\,G_{X,\infty}\,K$.
\end{proof}

\begin{proposition}[Standard frames are good]\label{pf:fs:stdGood}
Assume $\fsP{DF}$. In every setting, for every core volume $v$, every standard tower frame is good for its line lattices
$(\Gamma^L_\rho)$.
\formallabel{GlobalTheory.T.standardAuxiliaryFrames_good_row_grids}
\formallabel{GlobalTheory.AscFin.GoodStd.std_good}
\end{proposition}

\begin{proof}
Let $\mathcal F$ be standard. The goodness of $\mathcal F$ (\cref{pf:fs:frames}) consists of four conditions.

\emph{Lattice data.} The lattices $\Gamma^L_\rho$ carry non-zero vacua, and the last seesaw at the bottom has its lattice data
(\lref{GlobalTheory.T.last_seesaw_lattices_nonzero_line_vacua}: for a standard frame, the annihilator of the core depth lattice
$\Gamma_{\fsc(0)}$ is carried into the annihilators of $\Gamma_{\chi(0,0)}$ and $\Gamma^L_0$ by the decomposition $V_{\fsc_0}=Y_0\oplus L_z$, and
there are non-zero vacua $f_1$ for $\Gamma_{\chi(0,0)}$ and $f_2$ for $\Gamma^L_0$ whose box is a $\Gamma_{\fsc(0)}$-vacuum).

We use three general facts about the grid. (G1) For every node $(i,0)$ with $i\le m$, the group $K_{(i,0)}$ is a compact open subgroup of
the finite part of $G_{V_{(i,0)}}$; and for a subgroup $R$ of $G_{V_{(i,0)}}$, an element $h$ lies in $R\,G_\infty\,K_{(i,0)}$ as soon as its
component at every finite place outside $S_0$ lies in the local coset of $R$ and the lattice there, and its component at $S_0$ is the
product of an element of $R$ and an element of $K_{(i,0)}$, both supported at $S_0$
(\lref{GlobalTheory.T.depthCongruenceSubgroup_compactOpen_places_to_adeles}). (G2) The stabiliser of the lattice at a
place outside $S_0$ lies in the level (\lref{GlobalTheory.T.latticeStabilizerAtPlace_mem_depthCongruence_gridNode}). (G3) Given
$\fsP{DF}$, the conjugated Weil representations of the first row satisfy the conjugation dictionary of \ref{pf:fs:T1}; in particular,
for $i<m$, every depth vacuum $f$ for the conjugated pair $\bar\omega_{\chi(i,0)}$ satisfies $\bar\omega_{\chi(i,0)}(1,k)f=f$ for
$k\in K_{(i,0)}$ (\lref{GlobalTheory.T.weil_datum_conjugation_dictionary_row_grid}).

\emph{Bessel-to-Fourier--Jacobi steps.} Let $2\rho+2\le m$. The Bessel datum of $\mathcal F$ at $\rho$ has $V=V_{(2\rho+2,0)}$, $V_0=V^0_\rho$,
$A=V_{\fsc_{2\rho+2}}$, $L=-L_z$ and $\zeta=\zeta_{\rho+1}$, and the Fourier--Jacobi datum at $\rho$ has $W=V_{(2\rho+1,0)}$ and
$W_0=V_{\fsc_{2\rho+1}}$ (with the measures of $v$ and of $\mathcal F$). Let $s$ be a Bessel-to-Fourier--Jacobi step for these data whose Weil
representation is $\bar\omega_{\chi(2\rho+1,0)}$. Its matrix $x$ is the image of the rational
isotropic family of the Fourier--Jacobi decomposition, which has linearly independent columns (\lref{GlobalTheory.Unf.base_point}). The
mixed-model pieces of the step are admissible and split the depth lattice $\Gamma_{\chi(2\rho+1,0)}$
(\lref{GlobalTheory.T.besselToFourierJacobi_mixedPieces_admissible}). We apply \cref{pf:fs:locCls} with $K=K_{(2\rho+1,0)}$ and
$\Gamma=\Gamma_{\chi(2\rho+1,0)}$: $K$ is compact open by (G1), and the invariance hypothesis is (G3). For the support hypothesis, let
$h$ and all $h\kappa$ satisfy the unfolding conditions; by (G1) it suffices to decompose the local components. The non-vanishing condition
localises at every place (\lref{GlobalTheory.T.flatMapNonvanishingCondition_localise}). At a place $w$ outside $S_0$, the frame is in good
position (\lref{GlobalTheory.T.goodPosition_and_incAdj_mem_latO_iff_mem_intAt}), and an element satisfying the localised condition, with
all its translates by the local stabiliser (which lies in the level by (G2) and commutes with the components at other places), lies in a
single coset (\lref{GlobalTheory.T.goodPlace_oneCoset_of_besselLevel}); at the places of $S_0$ the same holds by
\lref{GlobalTheory.T.besselLevel_oneCoset_at_exceptionalPlaces}. \Cref{pf:fs:locCls} gives local seesaw data
$(\mathcal S,\Gamma_0,K_{(2\rho+1,0)})$ consisting of depth vacua for $\chi(2\rho+1,0)$. Finally, if $\Xi_1$ is a depth vacuum for $\fsc(2\rho+1)$ and
$\Xi_2$ a theta function such that $w\mapsto\overline{\Xi_2(\kappa_\rho w)}$ is a $\Gamma^L_\rho$-vacuum, then the box of $\Xi_1$ and $\Xi_2$ along
$V^0_\rho=V_{\fsc_{2\rho+2}}\oplus(-L_z)$ is a $\Gamma_0$-vacuum (\lref{GlobalTheory.T.boxDepthData_mul_lineVacua_besselLevel}).

\emph{Fourier--Jacobi-to-Bessel steps.} Let $2\sigma+4\le m$ and let $s$ be a Fourier--Jacobi-to-Bessel step for the Fourier--Jacobi datum of
$\mathcal F$ at $\sigma+1$ and the Bessel datum at $\sigma$, whose Weil representation is $\bar\omega_{\chi(2\sigma+2,0)}$. Its matrix $x$ is the
image of a rational matrix with independent columns (a direct verification with the decompositions of $\mathcal F$), and the pieces of the step
are admissible and split
$\Gamma_{\chi(2\sigma+2,0)}$ (\lref{GlobalTheory.T.fourierJacobi_step_pieces_admissible_splits}). The argument is the same, with
$K=K_{(2\sigma+2,0)}$: at a place outside $S_0$ the successive levels of the standard frame are in good position
(\lref{GlobalTheory.T.standard_frames_successive_levels}), the grid lattices are self-dual and split compatibly
(\lref{GlobalTheory.T.grid_lattices_compact_selfDual_splitting}), and the one-coset property holds
(\lref{GlobalTheory.T.fourierJacobiLevel_oneCoset_at_goodPlace}); at the places of $S_0$ it holds by
\lref{GlobalTheory.T.one_coset_decomposition_statement}. The boxes of depth vacua for $\fsc(2\sigma+2)$ with $\Gamma^L_{\sigma+1}$-vacua are
$\Gamma_0$-vacua (\lref{GlobalTheory.T.fourierJacobi_depthDatum_mul_lineVacuum_isVacuum}).
\end{proof}

\begin{theorem}[Towers exist]\label{pf:fs:towerEx}
$\fsP{DF}$ implies $\fsP{TE}$.
\formallabel{GlobalTheory.T.existsGoodGridFrames_of_depthFixed}
\formallabel{GlobalTheory.AscFin.towerEx_of_depFixed}
\end{theorem}

\begin{proof}
Given a setting and a core volume $v$, take a standard tower frame (\cref{pf:fs:stdEx}). It has good ends (\cref{pf:fs:ends}) and, by
$\fsP{DF}$, it is good for its line lattices (\cref{pf:fs:stdGood}).
\end{proof}

\subsubsection{Chains along the core column}

For $k\le m$ and a function $F$ on $G_{V_{\fsc_k}}$ let $\mathrm{Bel}_k(F)$ be the set of the bottoms $\Psi_0$ of the families $(\Psi_i)_{i\le k}$
of functions on $G_{V_{\fsc_i}}$ such that $\Psi_i$ is rapidly decreasing for $i\le k$, $\Psi_i=\fsdn_i(f_i,\Psi_{i+1})$ for a depth vacuum
$f_i$ for $\fsc(i)$ ($i<k$), and $\Psi_k=F$. Thus $\fsP{CHN}$ (\cref{pf:fs:corestatements}) asks for $\Theta(y)(x)\in\mathrm{span}\bigcup_w
\mathrm{Bel}_k(y(w))$.

\begin{proposition}[Chains along the core column]\label{pf:fs:chainEx}
Assume $\sfname{DiscreteSeriesCover}$, $\sfname{HarishChandraModules}$, $\sfname{CuspRapidDecay}$, $\sfname{CasselmanWallach}$, $\fsP{DF}$
and $\fsP{CM}$. Then $\fsP{CHN}$ holds.
\formallabel{GlobalTheory.T.core_chain_composites_exist_lift_combinations}
\formallabel{GlobalTheory.ChainLifts.chainEx}
\end{proposition}

\begin{proof}
Fix a setting with $\sfname{Sizes}(D)$ and $c_{\mathrm{gap}}=1$, a core volume $v$, a core lattice $r_0$ at $\fsc_0$ and a model $\mathcal M_0$
of $V_{\fsc_0}$; write $s_i(r,\mathcal M)$ for the grid stage $s_i(r,v,\mathcal M)$. We use: $\fsP{COM}$ and $\fsP{STG}$ (cited in
\cref{sec:pf:fs:onto}); the existence of models at the core nodes ($\fsP{CM}$); and the following facts.
\begin{enumerate}[label=(C\arabic*)]
\item (pre-steps; \lref{GlobalTheory.T.core_chain_pre_steps_exist_of_depth_congruence}, using $\fsP{DF}$) For $k<m$, core lattices $r,r'$ at
$\fsc_k,\fsc_{k+1}$ and models $\mathcal M,\mathcal M'$ there are a Weil datum $\omega$ for $(V_{\fsc_k},V_{\fsc_{k+1}})$ and a vacuum vector $f$
which satisfy the conditions of a good step of index $k$ from $s_k(r,\mathcal M)$ to $s_{k+1}(r',\mathcal M')$ not involving the maps
$\mathrm{It}$ and $\mathrm{io}$; the theta kernel of $\omega$ at $(g,h)$ applied to a theta function $\phi$ equals the theta kernel of the grid Weil
datum $\omega_{\fsc(k)}$ at $(h,g)$ (with $h$ viewed in $G_{-V_{\fsc_{k+1}}}$) applied to the exchanged function $\mathrm{ex}(\phi)$ on
$\fsX(-V_{\fsc_{k+1}},V_{\fsc_k})$; and for every polynomial $P$ the exchanged Fock vector $\mathrm{ex}(\mathrm{pt}_f(P))$ is a depth vacuum for
$\fsc(k)$.
\item (steps; \lref{GlobalTheory.T.evenRankThetaChain_weilDatum_extends_to_step}, using the four fields of $\sfname{ThetaLocalInputs}$
assumed) Such a pair $(\omega,f)$ between good stages extends to a good step $p$ of index $k$ with $\omega_p=\omega$ and $f_p=f$.
\item (rapid decay; \lref{GlobalTheory.T.multiplicity_values_continuous_bounded_hecke_conj} together with the lattice change
\lref{GlobalTheory.T.selfDualLattice_changeOfLattice_laws}) For $i\le m$, a core lattice $r$, a model $\mathcal M$ and
$y\in\mathrm{CC}(s_i(r,\mathcal M))$, every value $y(w)$ is rapidly decreasing.
\item (exchange; \lref{GlobalTheory.T.adjointExchange_laws_thetaHeisenberg}) A function on $\fsX(V_{\fsc_k},V_{\fsc_{k+1}})$ is a theta
function if and only if its exchange is.
\item (onto; \lref{GlobalTheory.ItOnto.kfinSpan_le_span_it}, using $\sfname{DiscreteSeriesCover}$ and $\sfname{HarishChandraModules}$) For a
good step $p$ from $s$ to $s'$, every $K$-finite vector of the model of $s$ lies in the span of the vectors $\mathrm{It}_p(P,w')$ with $P$ a
polynomial and $w'$ a $K$-finite vector of the model of $s'$.
\end{enumerate}
We also use the following extension property, which is immediate from the definition: if $f$ is a depth vacuum for $\fsc(k)$, $F$ is rapidly
decreasing on $G_{V_{\fsc_{k+1}}}$ and $F_0\in\mathrm{Bel}_k(\fsdn_k(f,F))$, then $F_0\in\mathrm{Bel}_{k+1}(F)$ (extend the family by
$\Psi_{k+1}=F$).

\emph{Claim.} For every $k\le m$ there are a core lattice $r$ and a model $\mathcal M$ at $\fsc_k$, a chain by which $s_0=s_0(r_0,\mathcal M_0)$
reaches $s_k(r,\mathcal M)$ with composite $\Theta$, a set $I$ and linear maps $N_\iota$ ($\iota\in I$) from the smooth vectors of $\mathcal M$
to those of $\mathcal M_0$ such that
\begin{enumerate}[label=(\roman*)]
\item $\Theta(y)(N_\iota w)\in\mathrm{Bel}_k(y(w))$ for all $\iota$, $y\in\mathrm{CC}(s_k(r,\mathcal M))$ and $w$;
\item every $K$-finite vector of $\mathcal M_0$ lies in the span of the vectors $N_\iota w$ with $w$ a $K$-finite vector of $\mathcal M$.
\end{enumerate}

For $k=0$ take $r=r_0$, $\mathcal M=\mathcal M_0$, the trivial chain with $\Theta=\mathrm{id}$, a one-point set $I$ and $N=\mathrm{id}$; (i)
holds with the family $(y(w))$, which is rapidly decreasing by (C3), and (ii) is trivial. Suppose the claim holds for $k<m$, with
$(r,\mathcal M,\Theta,I,N)$. By $\fsP{STG}$ and $\fsP{CM}$ choose a core lattice $r'$ and a model $\mathcal M'$ at $\fsc_{k+1}$; the stages
$s=s_k(r,\mathcal M)$ and $s'=s_{k+1}(r',\mathcal M')$ are good of heights $k$ and $k+1$. Take $(\omega,f)$ by (C1) and extend it to a good step $p$
by (C2). Then $s_0$ reaches $s'$ by $\Theta\circ\mathrm{io}_p$. Put $I'=I\times\{\text{polynomials}\}$ and $N'_{(\iota,P)}=N_\iota\circ\mathrm{It}_p(P,\cdot)$.

(i) Let $y\in\mathrm{CC}(s')$ and $w$. By the induction hypothesis applied to $\mathrm{io}_p(y)\in\mathrm{CC}(s)$ and the vector $\mathrm{It}_p(P,w)$,
$\Theta(\mathrm{io}_p(y))(N_\iota\mathrm{It}_p(P,w))\in\mathrm{Bel}_k\bigl(\mathrm{io}_p(y)(\mathrm{It}_p(P,w))\bigr)$. By the defining property of
$\mathrm{io}_p$, $\mathrm{io}_p(y)(\mathrm{It}_p(P,w))$ is the theta lift of $y(w)$ by $\mathrm{pt}_f(P)$ for $\omega$, and by the kernel identity of
(C1) (with (C4) to see that $\mathrm{pt}_f(P)$ is a theta function) this is
\[
\mathrm{io}_p(y)(\mathrm{It}_p(P,w))=\fsdn_k\bigl(\mathrm{ex}(\mathrm{pt}_f(P)),\,y(w)\bigr).
\]
Since $\mathrm{ex}(\mathrm{pt}_f(P))$ is a depth vacuum (C1) and $y(w)$ is rapidly decreasing (C3), the extension property gives
$\Theta(\mathrm{io}_p(y))(N'_{(\iota,P)}w)\in\mathrm{Bel}_{k+1}(y(w))$.

(ii) A $K$-finite vector of $\mathcal M_0$ lies in the span of the $N_\iota w$ with $w$ $K$-finite for $\mathcal M$; by (C5) each such $w$ lies in
the span of the $\mathrm{It}_p(P,w')$ with $w'$ $K$-finite for $\mathcal M'$; by linearity of $N_\iota$ the vector lies in the span of the
$N'_{(\iota,P)}w'$.

\emph{Conclusion.} Given $k\le m$, take the data of the claim. For $y\in\mathrm{CC}(s_k(r,\mathcal M))$ and $x$ $K$-finite for $\mathcal M_0$, write
$x=\sum_jc_jN_{\iota_j}w_j$ by (ii); by linearity of $\Theta(y)$ and (i), $\Theta(y)(x)=\sum_jc_j\Theta(y)(N_{\iota_j}w_j)$ lies in the span of
$\bigcup_w\mathrm{Bel}_k(y(w))$. This is $\fsP{CHN}$.
\end{proof}

\subsubsection{Poincar\'e series}

We keep the notation of \cref{sec:pf:fs:periods}. An \emph{isotropic frame} of a hermitian space $V$ is a family of $r\ge1$ linearly
independent vectors of $E^{\dim V}$ spanning a totally isotropic subspace; $N_C\subseteq G_V$ denotes the unipotent radical attached to an
isotropic frame $C$ (the group generated by the elements which act trivially on the span of $C$ and on its orthogonal modulo it), and
$N_C(F)$ its rational points. A function $\phi$ on $G_V$ is \emph{cuspidal} if for every isotropic frame $C$, every Haar measure $\mu$ on $N_C$
and every fundamental domain $\fsF$ of $N_C(F)$ in $N_C$, $\int_{\fsF}\phi(ng)\,d\mu(n)=0$ for all $g\in G_V$.

\begin{lemma}[Poincar\'e series are cuspidal]\label{pf:fs:poinCusp}
Let $\mathcal M$ be a model of $V$, $\Lambda$ a self-dual lattice, $\fsm\ne0$ an ideal and $K$ the principal congruence subgroup of $\Lambda$
modulo $\fsm$. Let $v=(v_\tau)$ be a family of $K$-finite vectors such that for every $\tau$ and every $K$-finite $a\in H_\tau$ the function
$\langle v_\tau,\pi_\tau(g)a\rangle(1+\|g\|)^{2\dim V^2+1}$ is bounded on $\Uni(V^\tau)$. Assume $\fsT\ne\emptyset$. Then for every $c\in G_V$ and
every family $a$ of $K$-finite vectors, the Poincar\'e series $P_K(v,c,a)$ is cuspidal.
\formallabel{UpperTheory.Lit.Cusp.poin_isCuspFun}
\end{lemma}

\begin{proof}
Let $C$ be an isotropic frame, $N=N_C$, $\mu$ a Haar measure on $N$, $\fsF$ a fundamental domain of $N(F)$ and $g\in G_V$. Put
$f(y)=\mathrm{coef}_K(v,a)(cyg)$, so that $P_K(v,c,a)(ng)=\sum_{\gamma\in\Gamma_V}f(\gamma n)$; the coefficient is measurable, $K$ being open.

\emph{Unfolding.} There is a compact set $K_c\subseteq N$ whose translates by $N(F)$ cover $N$
(\lref{GlobalTheory.Mem.UnipCocompact.radicals_cocompact}). The partial sums of $\sum_\gamma|\mathrm{coef}_K(v,a)(c\gamma y)|$ are bounded
uniformly for $y$ in the compact set $K_cg$, by a constant $B$ (\lref{GlobalTheory.T.poincareSeries_partialSums_bounded_compact}, which
uses the decay of the coefficients). Hence $\sum_\gamma|f(\gamma n)|\le B$ for $n\in K_c$, and by the covering,
\[
\sum_{q\in\Gamma_V/N(F)}\int_N|f(qn)|\,d\mu(n)\le\int_{K_c}\sum_{\gamma\in\Gamma_V}|f(\gamma n)|\,d\mu(n)\le B\,\mu(K_c)<\infty
\]
(\lref{GlobalTheory.T.lintegral_unfolding_le_of_translate_cover}). The unfolding lemma for the closed subgroup $N$ and the discrete group
$\Gamma_V$ (\lref{GlobalTheory.T.integrableOn_and_hasSum_tsum_subgroup}) then shows that $n\mapsto P_K(v,c,a)(ng)$ is integrable on $\fsF$ and
\[
\int_{\fsF}P_K(v,c,a)(ng)\,d\mu(n)=\sum_{q\in\Gamma_V/N(F)}\int_Nf(qn)\,d\mu(n),
\]
$q$ running over representatives of the cosets.

\emph{The terms vanish.} Fix $\tau\in\fsT$. The group $N$ is the product of two commuting closed subgroups $P$ and $Q$ with $P\cap Q=1$ ($P$ the
elements of $N$ supported at the place of $\tau$, $Q$ those with trivial component there), and for every Haar measure $\nu$ on $P$, every
$q'\in Q$ and all $x,y\in G_V$,
\[
\int_P\mathrm{coef}_K(v,a)(x\,p\,q'\,y)\,d\nu(p)=0
\]
(\lref{UpperTheory.Lit.Cusp.fibre_integral_eq_zero}: as $p$ varies only the factor at $\tau$ of the coefficient varies, and the integral of
a matrix coefficient of a square-integrable irreducible representation of $\Uni(V^\tau)$, with decaying coefficients, along the unipotent
radical of an isotropic subspace vanishes, \lref{UpperTheory.Lit.Cusp.integral_coef_eq_zero}, by the Schur orthogonality relations).
A function on $N$ whose integrals along all the fibres $Pq'$ vanish has integral $0$ over $N$
(\lref{GlobalTheory.T.integral_haar_eq_zero_of_fibre_integrals_eq_zero_alt}, Fubini's theorem for $N=PQ$). Hence every term
$\int_Nf(qn)\,d\mu(n)$ vanishes, and so does the constant term.
\end{proof}

\begin{lemma}[Cusp forms are small at short points]\label{pf:fs:shortLine}
Let $V$ be a hermitian space with $\dim V\ge1$, $A$ a frame of $V$, $\Lambda$ a self-dual lattice and $g,q\ge1$ integers such that at every $\tau$
the smaller of the numbers of positive and of negative signs of $A_\tau$ is at most $q$. Let $\phi$ be a cuspidal automorphic function on $G_V$,
right invariant under the principal congruence subgroup of $\Lambda$ modulo $g$, and $c$-Lipschitz for $A$ ($c\ge0$). If $0<s\le\frac12$ and
$x\in G_V$ is not $s$-long for $(A,\Lambda)$, then
\[
|\phi(x)|\le c\,(\dim V)^2f\,g\,\mathrm{rd}_E\,(2s)^{1/q}.
\]
\formallabel{GlobalTheory.T.norm_cuspForm_le_of_shortLine}
\formallabel{GlobalTheory.ShortLineCusp.norm_le_of_not_long}
\end{lemma}

\begin{proof}
\emph{A model.} There is an isomorphism of $G_V$ with the adelic unitary group of a standard model of $(V,\Lambda)$ (a hermitian lattice over
$E$ with a frame at every archimedean place whose signs are those of $A$ and whose frame matrices are the inverses of those of $A$;
\lref{GlobalTheory.T.exists_instanceFrames_of_diagonalFrame_selfDual}), with $f$ archimedean places and root discriminant $\mathrm{rd}_E$.
It carries rational points to rational points, and for every rational isotropic subspace $W$ of the model there is an isotropic frame
$C_0$ of $V$ whose radical $N_{C_0}$ is carried into the unipotent radical of $W$ (\lref{GlobalTheory.BridgeDict.bridge_dict}); it is
compatible with the lattices $\Lambda_x$, with the norms $\mathrm{nrm}_\tau(A,x;\cdot)$ of rational vectors and with the principal congruence
subgroups of level $g$ (routine verifications with the frames).
We identify $x$ with the corresponding point of the model.

\emph{Step 1: the least size.} The model attaches to $x$ a \emph{least size} $\lambda(x)\ge0$ such that
$\lambda(x)^f\sqrt{\Nm\mathfrak a}\le\prod_\tau\mathrm{nrm}_\tau(A,x;y)$ for every non-zero ideal $\mathfrak a$ and every non-zero $y\in\mathfrak a\Lambda_x$
(\lref{GlobalTheory.T.leastSize_pow_mul_sqrt_absNorm_le_prod_norms}). If $\lambda(x)\ge s$, these inequalities give
$s^f\sqrt{\Nm\mathfrak a}\le\prod_\tau\mathrm{nrm}_\tau(A,x;y)$ for all such $\mathfrak a,y$, that is, $x$ is $s$-long. Hence $\lambda(x)<s\le\frac12$.

\emph{Step 2: an isotropic reduction.} Since $\lambda(x)<\frac12$ and the number of signs is bounded by $q$, there is a rational isotropic
subspace $W$ of the model such that every element $m$ of its unipotent radical satisfies $mx=\gamma\,x\,u\,k$ with $\gamma$ rational, $k$ in the
principal congruence subgroup of level $g$, and $u$ archimedean with $u_v=\exp(A_v)\exp(Z_v)$, $A_v,Z_v$ in the Lie algebra at $v$, and
\[
\Bigl(\sum_v\|A_v\|^2\Bigr)^{1/2}+\Bigl(\sum_v\|Z_v\|^2\Bigr)^{1/2}\le\tfrac14(\dim V)^2\sqrt f\,g\,\mathrm{rd}_E\,(2\lambda(x))^{1/q}
\]
(\lref{GlobalTheory.T.automorphicPoint_exists_isotropic_reduction_of_small_size}).

\emph{Step 3: the variation along the radical.} Let $m\in N_{C_0}$, the radical corresponding to $W$. By Step 2, automorphy and the
invariance under the level, $\phi(mx)=\phi(xu)$. The Lipschitz property along the two exponentials gives
$|\phi(xu)-\phi(x)|\le c\sum_v(\|A_v\|+\|Z_v\|)$ (\lref{GlobalTheory.T.lipschitz_bound_real_unipotent_action}), and by the
Cauchy--Schwarz inequality over the $f$ places and Step 2,
\[
|\phi(mx)-\phi(x)|\le c\sqrt f\Bigl(\bigl(\textstyle\sum_v\|A_v\|^2\bigr)^{1/2}+\bigl(\sum_v\|Z_v\|^2\bigr)^{1/2}\Bigr)\le
c\,(\dim V)^2f\,g\,\mathrm{rd}_E\,(2s)^{1/q},
\]
using $\lambda(x)<s$.

\emph{Step 4.} Let $\mathfrak N=N_{C_0}(F)\backslash N_{C_0}$; it is compact and carries an invariant probability measure $dn$, and by cuspidality
$\int_{\mathfrak N}\phi(nx)\,dn=0$. Hence $\phi(x)=\int_{\mathfrak N}(\phi(x)-\phi(nx))\,dn$ and $|\phi(x)|\le\sup_{m\in N_{C_0}}|\phi(mx)-\phi(x)|$, which by
Step 3 gives the assertion.
\end{proof}

\begin{lemma}[The identity term dominates]\label{pf:fs:idDom}
Let $V=Y\oplus\ell$ be an orthogonal extension of $Y$ by a line with basis vector $z_\ell$, $\iota\colon G_Y\to G_V$ the embedding, $\mathcal M_0$ and
$\mathcal M$ models of $Y$ and $V$, and $p=(A,B,v_0,g_0,w)$ a pack for $(\mathcal M_0,\mathcal M)$ whose frames are adapted to the decomposition;
put $v_\tau=\pi_\tau((g_0)_\tau)v_{0,\tau}$. Let $Q\in\R$, $\Lambda$ a self-dual lattice of $Y$, $\Lambda_V$ the lattice of $V$ built on $\Lambda$ and
$z_\ell$, $\fsm\ne0$, $K$ the principal congruence subgroup of $\Lambda_V$ modulo $\fsm$, and $(\nu,\fsF)$ a Haar measure on $G_Y$ with a
fundamental domain for $\Gamma_Y$. Let $\phi$ be a continuous left $\Gamma_Y$-invariant function on $G_Y$ with $|\phi|\le b$, let $x\in G_Y$ with
$|\phi(x)|>b/2$, $0\le s\le1$, and suppose that $x$ is $s$-long for $(A,\Lambda)$. Suppose that there is $c\ne0$ with the following
\emph{adelic clauses}:
\begin{enumerate}[label=(A\arabic*)]
\item $\int\mathrm{coef}_K(v,v)(\iota(x)^{-1}\iota(h))\,\phi(h)\,d\nu(h)=c\,\phi(x)$;
\item for $\gamma\in\Gamma_V$ and $z=\gamma z_\ell$, the function $h\mapsto\mathrm{coef}_K(v,v)(\iota(x)^{-1}\gamma\iota(h))$ is integrable,
\[
\Bigl(\int|\mathrm{coef}_K(v,v)(\iota(x)^{-1}\gamma\iota(h))|\,d\nu(h)\Bigr)\Bigl(\prod_\tau\mathrm{nrm}_\tau(B,\iota x;z)\Bigr)^{2\dim V+5}\le e^Q|c|,
\]
and it vanishes identically unless $z-z_\ell\in\fsm\Lambda_{V,\iota x}$.
\end{enumerate}
If $2\exp(Q+(6\dim V+9)f)\le s^f\sqrt{\Nm\fsm}$ and $1<s^f\sqrt{\Nm\fsm}$, then
\[
\int_{\fsF}P_K\bigl(v,\iota(x)^{-1},v\bigr)(\iota h)\,\phi(h)\,d\nu(h)\neq0 .
\]
\formallabel{GlobalTheory.T.poincareSeries_pairing_ne_zero_of_large_level}
\formallabel{GlobalTheory.IdDom.identity_dominates}
\end{lemma}

\begin{proof}
The embedding $\iota$ is continuous and injective, $\iota^{-1}(\Gamma_V)=\Gamma_Y$, and an element of $G_V$ lies in $\iota(G_Y)$ if and only if it
fixes $z_\ell$ (\lref{GlobalTheory.T.exists_unitary_embedding_stabilizer}). Put $f(g)=\mathrm{coef}_K(v,v)(\iota(x)^{-1}g)$; then
$P_K(v,\iota(x)^{-1},v)(\iota h)=\sum_{\gamma\in\Gamma_V}f(\gamma\iota(h))$.

\emph{Cosets and vectors.} For a coset $q\in\Gamma_V/\iota(\Gamma_Y)$ with representative $\gamma_q$ put $z_q=\gamma_qz_\ell$. If $z_q=z_{q'}$, then
$\gamma_q^{-1}\gamma_{q'}$ fixes $z_\ell$, hence lies in $\iota(G_Y)\cap\Gamma_V=\iota(\Gamma_Y)$, and $q=q'$; the coset $q_1$ of $1$ is the only one
with $z_q=z_\ell$. Put $\beta_q=b\int|f(\gamma_q\iota(h))|\,d\nu(h)$, so that $|\int f(\gamma_q\iota(h))\phi(h)\,d\nu(h)|\le\beta_q$.

\emph{Lengths.} Since $x$ is $s$-long for $(A,\Lambda)$ and the frames are adapted, $\iota(x)$ is $s$-long for $(B,\Lambda_V)$,
$\mathrm{nrm}_\tau(B,\iota x;z_\ell)=1$, and $\mathrm{nrm}_\tau(B,\iota x;z)\ge1$ for every $z=\gamma z_\ell$ with $\gamma\in\Gamma_V$
(\lref{GlobalTheory.T.long_element_embedding_norm}). By (A2), $\beta_q\le b\,e^Q|c|$ for every $q$. For $q\ne q_1$ put
$w_q=z_q-z_\ell$; by the triangle inequality and $\mathrm{nrm}_\tau(z_q)\ge1$, $(1+\mathrm{nrm}_\tau(B,\iota x;w_q))/3\le\mathrm{nrm}_\tau(B,\iota x;z_q)$,
so (A2) gives
\[
\beta_q\le b\,e^Q|c|\prod_\tau\Bigl(\frac{1+\mathrm{nrm}_\tau(B,\iota x;w_q)}{3}\Bigr)^{-(2\dim V+5)},
\]
and $\beta_q=0$ unless $w_q\in L:=\fsm\Lambda_{V,\iota x}$.

\emph{The tail.} Let $R=s^f\sqrt{\Nm\fsm}>1$. As $\iota(x)$ is $s$-long, every non-zero $y\in L$ satisfies $\prod_\tau\mathrm{nrm}_\tau(B,\iota x;y)\ge R$.
For such lattices the weighted lattice sum is bounded: for every finite set of non-zero vectors of $L$,
\[
\sum_y\prod_\tau\Bigl(\frac{1+\mathrm{nrm}_\tau(B,\iota x;y)}{3}\Bigr)^{-(2\dim V+5)}\le\frac{e^{(6\dim V+9)f}}{R}
\]
(\lref{GlobalTheory.T.latticeSum_longVectors_decayWeight_le}). The vectors $w_q$, $q\ne q_1$, are distinct and non-zero, so every finite
partial sum of $\sum_{q\ne q_1}\beta_q$ is at most $B_t:=b\,e^Q|c|\,e^{(6\dim V+9)f}/R\le b|c|/2$, by the first hypothesis on $R$. In particular
$\sum_q\beta_q<\infty$.

\emph{Unfolding and conclusion.} By the summability, the unfolding lemma (\lref{GlobalTheory.T.integrableOn_and_hasSum_tsum_subgroup})
applies to $f$ and $\phi$ (left $\Gamma_Y$-invariant), and the integral of the statement equals
$\sum_q\int_{G_Y}f(\gamma_q\iota(h))\phi(h)\,d\nu(h)$. For $q=q_1$ the representative is $\iota(h_0)$ with $h_0\in\Gamma_Y$, and by the
$\Gamma_Y$-invariance of $\phi$, the invariance of $\nu$ and (A1) the term equals $c\,\phi(x)$. If the integral vanished, then
\[
|c|\,|\phi(x)|=\Bigl|\sum_{q\ne q_1}\int f(\gamma_q\iota(h))\phi(h)\,d\nu(h)\Bigr|\le\sum_{q\ne q_1}\beta_q\le\frac{b|c|}{2},
\]
contradicting $|\phi(x)|>b/2$ and $c\ne0$.
\end{proof}

\subsubsection{Grid parameters and packs}

\begin{proposition}[Grid parameters at large scale]\label{pf:fs:gridPar}
$\forall_{\mathrm{ev}}$ (in fact with threshold $x_0=0$): for every $\tau$, every enumeration $(\eps,\lambda)$ of the parameter $S^{\fsc}_0$ at $\tau$
and every enumeration $(\eps',\mu)$ of the dual of $S^Y$ at $\tau$, there is a bijection $e$ with $\mathrm{Par}(e,\eps,\lambda,\eps',\mu,b^2)$
(\cref{pf:fs:RI}); moreover $0\le\log b$ and $\log(8(r+1)b^2)\le6\log b$.
\formallabel{GlobalTheory.T.branchingHypotheses_gridParameters_largeScale}
\formallabel{GlobalTheory.SepPer.realInputPar}
\end{proposition}

\begin{proof}
Let $D$ satisfy $\sfname{SizeRows}$ (at any scale). The archimedean parameters of the grid at $\tau$ are described by numerical data: even
integers $k\ge2$ and $m\ge2$ with $k+m=r+1$, integers $v_i$ and $b_l$, a rational number $a_0$ and a gap $g_a$. Under $\sfname{SizeRows}$ they
satisfy (\lref{GlobalTheory.T.chiGridNumericalData_realPlace_gap_size_bounds})
\begin{gather*}
v_{i+1}+2\le v_i\ (2\le i<k),\qquad 2\le v_k+b_1,\qquad b_l+1000(r+1)^3\le b_{l+1}\ (1\le l<m),\\
1000(r+1)^3\le g_a,\quad 0\le v_k,\quad 0<b_1,\quad a_0\le b^2,\quad b_m\le b^2,\quad 1\le b,\quad 8(r+1)\le b^4 .
\end{gather*}
The parameter $S^{\fsc}_0$ at $\tau$ is the parameter of the core node $0$ built from these data, and the dual of $S^Y$ is the dual parameter at
$(0,0)$ built from them (the twisted parameter at $j=0$ is the grid parameter at $(0,0)$, \lref{GlobalTheory.ChiGrid.parV_zero}, duality
is an involution, and the shift at $(0,0)$ is trivial). From the displayed inequalities, the enumerations satisfy
$\mathrm{Par}(e,\eps,\lambda,\eps',\mu,b^2)$ for a suitable $e$ (\lref{GlobalTheory.SepPer.exists_par_of_gridData_bounds}: the numbers of
negative signs are half of the sizes by evenness, which gives a bijection $e$ matching the signs; the parameters are Harish-Chandra parameters,
interlaced also after lowering one entry, with consecutive entries at distance at least $2$ and gaps at least $1000(r+1)^3$, and bounded by
$\max(a_0,b_m)\le b^2$). Finally $\log b\ge0$ as $b\ge1$, and $8(r+1)b^2\le b^4b^2=b^6$ gives $\log(8(r+1)b^2)\le6\log b$.
\end{proof}

\begin{definition}[Local pack data]\label{pf:fs:locPack}
Let $V=Y\oplus\ell$ be as in \cref{pf:fs:idDom}, $\mathcal M_0,\mathcal M$ models of $Y,V$, $\tau\in\fsT$ and $q\in\R$. \emph{Local pack data at $\tau$
with loss $q$} consist of: signs $\eps_A$ and an invertible matrix $A$ with $A^*J_{\eps_A}A=Y^\tau$; signs $\eps_B$ and an invertible matrix $B$
with $B^*J_{\eps_B}B=V^\tau$, adapted in the sense that $|B(\iota y+az_\ell)|^2=|Ay|^2+|a|^2$; $g_0\in\Uni(V^\tau)$; a $K$-finite $v_0\in H_\tau(\mathcal M)$
with $\langle v_0,\pi_\tau(g)v_0\rangle(1+\|g\|)^{2\dim V^2+1}$ bounded; a smooth $w\ne0$ in $H_\tau(\mathcal M_0)$; and $c\ne0$; such that, with
$m(g)=\langle\pi_\tau(g_0)v_0,\pi_\tau(g)\pi_\tau(g_0)v_0\rangle$, $u\mapsto\tilde u$ the identification of $\Uni(J_{\eps^0})$ with $\Uni(Y^\tau)$ given
by the frame of $\mathcal M_0$ at $\tau$ ($\eps^0$ the signs of $\mathcal M_0$ at $\tau$), and $du$ the Haar measure of $\Uni(J_{\eps^0})$:
\begin{enumerate}[label=(\roman*)]
\item $\int m(\iota\tilde u)\,\langle b,\pi^0_\tau(\tilde u)w\rangle\,du=c\langle b,w\rangle$ for all $b$;
\item for $g\in\Uni(V^\tau)$, $u\mapsto m(g\iota\tilde u)$ is integrable and
$\bigl(\int|m(g\iota\tilde u)|\,du\bigr)\,|Bgz_\ell|^{2\dim V+5}\le e^q|c|$;
\item for $X$ in the Lie algebra of $\Uni(Y^\tau)$ and $s\in\R$, $u\mapsto|m(\iota(\exp(sX)\tilde u))-m(\iota\tilde u)|$ is integrable with integral at
most $e^q|c|\,|s|\,\|AXA^{-1}\|$.
\end{enumerate}
\formal{GlobalTheory.PackLoc.LocOK}
\end{definition}

\begin{proposition}[Local pack data from the archimedean input]\label{pf:fs:locOK}
Assume $\fsP{RI}$. There is $C\ge0$ (namely the constant of $\fsP{RI}$) such that $\forall_{\mathrm{ev}}$: for all models $\mathcal M_0$ of $Y_0$ with parameter $S^Y$ and $\mathcal M_{\fsc}$ of
$V_{\fsc_0}$ with parameter $S^{\fsc}_0$ and every $\tau$, there are local pack data at $\tau$ with loss $6C(r+1)^4\log b$.
\formallabel{UpperTheory.Lit.Branch.locOK_of_eigen}
\end{proposition}

\begin{proof}
\emph{Transport.} The statement $\fsP{RI}$ concerns the standard forms $J_\eps$ of sizes $n+1$ and $n$ with a bijection
$e\colon\{\ast\}\sqcup\{1,\dots,n\}\to\{1,\dots,n+1\}$; it transports, with the same constant, to index sets $\{1,\dots,k\}$ and $\{1,\dots,k+1\}$
related by an arbitrary bijection, and $\mathrm{Par}$ provides the sign conditions on $(\eps,\eps')$ (a routine reindexing).

Given $\eps$ and $C_F\ge6510$, take the threshold of \cref{pf:fs:gridPar}. In a setting at larger scale let $\mathcal M_0,\mathcal M_{\fsc}$ and $\tau$
be given; $\dim Y_0=r$ and $\dim V_{\fsc_0}=r+1$. Let $(\eps,\lambda)$ be the enumeration of the parameter of $\mathcal M_{\fsc}$ at $\tau$ and $(\eps',\mu)$
that of the conjugate model $\overline{\mathcal M_0}$ of $\mathcal M_0$ (a model of $Y_0$ with the dual parameter, on the conjugate Hilbert spaces). By
\cref{pf:fs:gridPar}, $\mathrm{Par}(e,\eps,\lambda,\eps',\mu,b^2)$ holds for some $e$, and $\log(8(r+1)b^2)\le6\log b$. Apply the transported $\fsP{RI}$ to
$\pi=\pi_\tau(\mathcal M_{\fsc})$, $\sigma=\pi_\tau(\overline{\mathcal M_0})$ and $B=b^2$: it gives a lowest $K$-type unit vector $v_0$, $g'_0$, a smooth unit
vector $w'$ and $c_1>0$ with (i)--(iv) of \cref{pf:fs:RI}, for $v=\pi(g'_0)v_0$.

\emph{Frames and constants.} Let $A$ be the frame of $\overline{\mathcal M_0}$ at $\tau$. There is a frame $B$ of $V^\tau$, adapted to the decomposition
and extending $A$, which maps $z_\ell$ to the basis vector of index $e(\ast)$ and conjugates the embedding $\iota$ into the embedding
$h\mapsto\tilde h$ of \cref{pf:fs:RI} (\lref{GlobalTheory.T.exists_signature_matrix_extension}, using the sign conditions). The Haar measure of
$\Uni(J_{\eps'})$ fixed in $\fsP{RI}$ is $\lambda$ times the Haar measure obtained by reversing coordinates, for some $\lambda>0$
(\lref{GlobalTheory.T.unitary_haar_coordinate_reversal_positive_multiple}). The matrix coefficients of $\mathcal M_{\fsc}$ at points of $\Uni(V^\tau)$
are those of $\pi$ in the frame $B$ (\lref{GlobalTheory.T.translated_matrix_coefficient_in_adapted_frame}), and those of $\overline{\mathcal M_0}$ are
computed in the frame $A$ (\lref{GlobalTheory.T.conjugate_model_frame_matrix_facts_at_place}). Take $g_0=B^{-1}g'_0F_\tau$, $F_\tau$ the frame of
$\mathcal M_{\fsc}$ at $\tau$, $w$ the conjugate of $w'$ (a smooth non-zero vector of $\mathcal M_0$,
\lref{GlobalTheory.T.conj_smooth_vector_conjugate_model_mem_smooth}), and $c=\lambda c_1\ne0$. With these identifications, $m(\iota\tilde u)$
is the restricted coefficient $\langle v,\pi(\tilde h)v\rangle$ of $\fsP{RI}$ at the element $h$ obtained from $u$ by reversing the coordinates, and
$m(g\iota\tilde u)$ is $\langle v,\pi(BgB^{-1}\tilde h)v\rangle$.

\emph{The clauses.} $v_0$ is $K$-finite (lowest $K$-type vectors are), and its coefficients decay with exponent $2\dim V^2+1$ by (iv) of $\fsP{RI}$
(\lref{GlobalTheory.T.matrix_coefficient_decay_in_frame_of_translate_decay}). Clause (i) follows from (i) of $\fsP{RI}$ after the change of Haar
measure and complex conjugation (a direct computation with the identifications above). Clause (ii) follows from (ii) of $\fsP{RI}$, since $\mathrm{col}(BgB^{-1})$
is $|Bgz_\ell|$ and $\lambda c_1\,\mathrm{big}(C,r,b^2)\le e^q\lambda c_1$ for $q=6C(r+1)^4\log b$, because
\[
\mathrm{big}(C,r,b^2)=\exp\bigl(C(r+1)^4\log(8(r+1)b^2)\bigr)\le\exp\bigl(6C(r+1)^4\log b\bigr)
\]
(\lref{GlobalTheory.T.transversal_coefficient_decay_at_real_place}). Clause (iii) follows from (iii) of $\fsP{RI}$, the integrability of the
coefficients of smooth vectors (\lref{GlobalTheory.T.smooth_vector_matrixCoeff_integrable}) and the same bound for $\mathrm{big}$
(\lref{GlobalTheory.T.matrix_coefficient_lipschitz_at_real_place}).
\end{proof}

\begin{proposition}[Packs exist]\label{pf:fs:packEx}
$\fsP{RI}$ implies $\fsP{PX}$.
\formallabel{UpperTheory.Lit.Branch.packEx_of_eigen}
\end{proposition}

\begin{proof}
Let $C$ be the constant of \cref{pf:fs:locOK}. Given $\eps$ and $C_F\ge6510$, take its threshold. In a setting at larger scale and for models
$\mathcal M_0,\mathcal M_{\fsc}$, choose local pack data at every $\tau$ with loss $q=6C(r+1)^4\log b$. Local data at all $\tau$ with a common loss $q$
assemble into a $(fq)$-good pack, formed from the local frames, vectors and points; the clauses of \cref{pf:fs:pack} follow from the local
clauses by taking products over $\tau$ (\lref{GlobalTheory.T.realVectors_allPlaces_of_localData}). Since
$fq=f\cdot6C(r+1)^4\log b=\mathrm{loss}(D,6C,b)$, the pack is $\mathrm{loss}(D,6C,b)$-good, and $\fsP{PX}$ holds with the constant $6C$.
\end{proof}

\subsubsection{Branching of discrete series}

We prove three ingredients of the proof of \cref{pf:fs:RIholds}. Throughout, $n\ge1$, $e\colon\{\ast\}\sqcup\{1,\dots,n\}\to\{1,\dots,n+1\}$ is a
bijection, $h\mapsto\tilde h$ is the embedding of $\Uni(J_{\eps'})$ into $\Uni(J_\eps)$ of \cref{pf:fs:RI}, $z=e(\ast)$, $\|g\|$ is the
Hilbert--Schmidt norm, $\mathrm{col}(g)$ the norm of the column $z$ of $g$, and $dg$ the Haar measure of $\Uni(J_\eps)$ fixed in the formal
development. A \emph{discrete series} of $\Uni(J_\eps)$ with parameter $\lambda$ is a representation with the properties of the formal
definition (irreducible, unitary, square integrable, with a lowest $K$-type determined by $\lambda$); $e_s$ is the $s$-th unit vector.

\begin{definition}[Branching conditions and occurrence]\label{pf:fs:gp}
Let $\eps\in\{\pm1\}^{n+1}$, $\lambda\in\R^{n+1}$, $\eps'\in\{\pm1\}^n$, $\mu\in\R^n$.
\begin{enumerate}[label=(\alph*)]
\item The \emph{branching conditions} $\mathrm{Bc}(e,\eps,\lambda,\eps',\mu)$ are: $\eps_z=+1$ and $\eps_{e(i)}=\eps'_i$; $n\ge2$; $\eps$ has a negative
entry and at most as many negative as positive entries; $\lambda$ and $\mu$ are Harish-Chandra parameters (strictly decreasing, with
$\lambda_a\in\frac n2+\Z$ and $\mu_j\in\frac{n-1}2+\Z$); and the gaps are large: $|\lambda_a-\lambda_b|\ge200(n+1)^3$ whenever $\eps_a\ne\eps_b$, and
$|\mu_i-\mu_j|\ge200(n+1)^3$ whenever $\eps'_i\ne\eps'_j$. For $B\ge0$, $\mathrm{Bc}_B$ adds the bounds $|\lambda_a|\le B$, $|\mu_j|\le B$.
\item $(\eps,\lambda)$ and $(\eps',\mu)$ are \emph{interlaced} if for any two consecutive values of the combined family $(\lambda_a)\cup(\mu_j)$
the signs agree exactly when the two values belong to different families.
\item \emph{Occurrence} holds at $(\lambda,\mu)$ (for given $e,\eps,\eps'$) if for every discrete series $\pi$ of $\Uni(J_\eps)$ with parameter
$\lambda$ on $H$ and every discrete series $\sigma$ of $\Uni(J_{\eps'})$ with parameter $\mu$ on $H'$ there is a non-zero continuous linear map
$j\colon H'\to H$ with $j\sigma(h)=\pi(\tilde h)j$ for $h\in\Uni(J_{\eps'})$.
\end{enumerate}
\formal{PositivityTheory.RgHe.Gp}
\formal{GlobalTheory.CompanionTheory.PairingAnalysis.Par}
\formal{GlobalTheory.CompanionTheory.PairingAnalysis.HeWord}
\formal{PositivityTheory.RgHe.OccAt}
\end{definition}

Note that the values of $\lambda$ and $\mu$ lie in different classes modulo $\Z$, so that the combined family has distinct values. Consider the
following three statements, each for all data satisfying the branching conditions.
\begin{enumerate}
\item[(L)] If $\mathrm{Bc}$ holds for $(\lambda-e_s,\mu)$ and $\lambda_s-\frac12\ne\mu_j$ for all $j$, then occurrence at $(\lambda,\mu)$ implies
occurrence at $(\lambda-e_s,\mu)$.
\item[(M)] If $\mathrm{Bc}$ holds for $(\lambda,\mu-e_j)$ and $\mu_j-\frac12\ne\lambda_s$ for all $s$, then occurrence at $(\lambda,\mu)$ implies
occurrence at $(\lambda,\mu-e_j)$.
\item[(B)] If $(\eps,\lambda)$ and $(\eps',\mu)$ are interlaced, there is a function $a$ from the index set
$\{1,\dots,n+1\}\sqcup\{1,\dots,n\}$ to $\N$, monotone with respect to the values (a smaller value never gets a larger $a$), such that
occurrence holds at $(\lambda+a|_\lambda,\mu+a|_\mu)$.
\end{enumerate}

\begin{lemma}[The walk]\label{pf:fs:walk}
(L), (M) and (B) imply that occurrence holds at $(\lambda,\mu)$ for all data satisfying $\mathrm{Bc}$ for which $(\eps,\lambda)$ and $(\eps',\mu)$ are
interlaced.
\formallabel{PositivityTheory.RgHe.Dec.walk}
\formallabel{PositivityTheory.RgHe.Dec.occurrence_of_walk}
\end{lemma}

\begin{proof}
For a monotone $b$ write $(\lambda^b,\mu^b)=(\lambda+b|_\lambda,\mu+b|_\mu)$. Such a shift preserves the branching conditions: the signs are
unchanged, the shifts are integers, the order is preserved (if $\lambda_b<\lambda_a$ then $b(\lambda_b)\le b(\lambda_a)$), and differences between
entries can only grow in absolute value. We show by induction on $k=\sum b$: \emph{if $b$ is monotone and occurrence holds at $(\lambda^b,\mu^b)$,
then it holds at $(\lambda,\mu)$.} For $k=0$, $b=0$. Let $k>0$ and choose, among the indices $x$ with $b(x)>0$, one whose value is minimal. Let $b'$ be
$b$ with $b(x)$ lowered by $1$. Then $b'$ is monotone: if $y\ne x$ has a value smaller than that of $x$ (values are distinct), then $b(y)=0$ by
the minimality of $x$, so $b'(y)=0\le b'(x)$. Suppose $x$ is the index $s$ of $\lambda$. The data at $b$ and at $b'$ satisfy $\mathrm{Bc}$, and
$\lambda^b_s-\frac12\ne\mu^b_j$ for all $j$: if $\mu_j<\lambda_s$, then $b(\mu_j)=0$ by minimality, so $\mu^b_j=\mu_j<\lambda_s<\lambda_s+b(s)-\frac12$;
if $\mu_j>\lambda_s$, then $b(\mu_j)\ge b(s)$ and $\mu^b_j>\lambda^b_s$. Since $(\lambda^{b'},\mu^{b'})=(\lambda^b-e_s,\mu^b)$, (L) gives occurrence at
$(\lambda^{b'},\mu^{b'})$, and the induction hypothesis (for $\sum b'=k-1$) gives occurrence at $(\lambda,\mu)$. The case of an index of $\mu$ is the
same with (M). Finally, (B) provides a monotone $a$ with occurrence at $(\lambda^a,\mu^a)$, and the claim with $b=a$ gives occurrence at
$(\lambda,\mu)$.
\end{proof}

We use four further statements of the development, each for all data satisfying $\mathrm{Bc}$, discrete series $\pi$ (parameter $\lambda$, on
$H$) and $\sigma$ (parameter $\mu$, on $H'$) and a Haar measure $dh$ on $\Uni(J_{\eps'})$.
\begin{enumerate}[label=(W\arabic*)]
\item\label{pf:fs:W1} (models; \lref{PositivityTheory.RgHe.LowestKTypeModel.models}) Occurrence holds at $(\lambda,\mu)$ as soon as a non-zero
intertwiner exists for one pair of discrete series with these parameters.
\item\label{pf:fs:W2} (integrals; \lref{PositivityTheory.RgHe.LowestKTypeModel.ofInt}) If $w$ is a lowest $K$-type unit vector of $\sigma$ and
$\int\overline{\langle w,\sigma(h)w\rangle}\langle b,\pi(\tilde h)a\rangle\,dh\ne0$ for some $a,b\in H$, there is a non-zero continuous intertwiner from
$\sigma$ to the restriction of $\pi$.
\item\label{pf:fs:W3} (Schur eigenprojections; \lref{PositivityTheory.RgHe.LowestKTypeModel.seed}) Suppose a non-zero intertwiner exists.
Let $\kappa>0$ be such that $\int\langle w_1,\sigma(h)v_1\rangle\overline{\langle w_2,\sigma(h)v_2\rangle}\,dh=\kappa\langle v_2,v_1\rangle\langle w_1,w_2\rangle$
(the integrands being integrable) and let $w$ be a lowest $K$-type unit vector of $\sigma$. Then the coefficient of $w$ is integrable, and there
are an isometric intertwiner $i\colon H'\to H$ and a smooth vector $\eta\in H$ with $\langle\eta,iw\rangle\ne0$ and
\begin{equation}\label{eq:pf:fs:eigproj}
\int\overline{\langle w,\sigma(h)w\rangle}\,\langle\eta,\pi(\tilde h)a\rangle\,dh=\kappa\,\langle\eta,iw\rangle\langle iw,a\rangle\qquad(a\in H).
\end{equation}
\item\label{pf:fs:W4} (decomposition; \lref{PositivityTheory.RgHe.Dec.decomposition_existence_holds}) Let $\pi$ be a discrete series of
$\Uni(J_\eps)$ ($m$ variables, both signs present) with Harish-Chandra parameter $\lambda$ and gap at least $200m^3$. Then the smooth vectors of
$\pi\otimes\C^m$ (the space $H^m$ on which $g$ acts by $\pi(g)$ and the standard representation) are the direct sum of components $p_s(x)$,
$s\in S$, where $S$ is the set of $s$ for which $\lambda-e_s$ is a Harish-Chandra parameter; the projections $p_s$ are equivariant; for each
$s\in S$ there are a discrete series $\pi_s$ with parameter $\lambda-e_s$ and an equivariant map $\mathrm{pr}_s$ from $H^m$ to its space, and a
constant $c_s$, such that the hermitian form $\sum_a\eps_a\langle y_a,x_a\rangle$ of $H^m$ satisfies
$\langle y,p_s(x)\rangle_\eps=c_s\langle\mathrm{pr}_s(y),\mathrm{pr}_s(x)\rangle$; and $p_s(x)$ is an eigenvector of a Casimir-type operator with eigenvalue
$\lambda_s+\frac{m-1}2$.
\end{enumerate}

\begin{lemma}[Lowering $\lambda$]\label{pf:fs:stepL}
\ref{pf:fs:W1}--\ref{pf:fs:W4} imply (L).
\formallabel{PositivityTheory.RgHe.Sl.stepL}
\formallabel{PositivityTheory.RgHe.Sl.stepL_of}
\end{lemma}

\begin{proof}
Let the data satisfy $\mathrm{Bc}$ for $(\lambda,\mu)$ and $(\lambda-e_s,\mu)$, with $\lambda_s-\frac12\ne\mu_j$ for all $j$, and suppose that occurrence
holds at $(\lambda,\mu)$. There are discrete series $\pi$ (parameter $\lambda$, on $H$) and $\sigma$ (parameter $\mu$, on $H'$)
(\lref{PositivityTheory.RgHe.Sl.exists_ds}, using the gap conditions), hence a non-zero intertwiner. Let $w$ be a lowest $K$-type unit vector of
$\sigma$ (the lowest $K$-type of a discrete series is non-zero) and $\kappa>0$ the constant of the Schur orthogonality relations of the
square-integrable representation $\sigma$ for the fixed Haar measure. By \ref{pf:fs:W3} there are $i$ and $\eta$ with
$\langle\eta,iw\rangle\ne0$ and \eqref{eq:pf:fs:eigproj}. Apply \ref{pf:fs:W4} to $\pi$ ($m=n+1$); since $\lambda-e_s$ is a Harish-Chandra parameter,
$s\in S$. The adjoint $P=i^*$ is a contraction intertwining the restriction of $\pi$ with $\sigma$, as $i$ is an isometric intertwiner between
unitary representations.

Let $y_0=p_s(e_z\otimes\eta)_z$ be the $z$-component of the $s$-component of $e_z\otimes\eta$. The eigenvalue relation of \ref{pf:fs:W4} and the
intertwining property of $P$ give the \emph{corner identity}
\[
\Bigl(\prod_{j\ne s}(\lambda_j-\lambda_s)\Bigr)P(y_0)=\Bigl(\prod_{j=1}^n\bigl(\mu_j-\lambda_s+\tfrac12\bigr)\Bigr)P(\eta)
\]
(\lref{PositivityTheory.RgHe.Sl.corner_val}). Pairing with $w$ (note $\langle iw,y\rangle=\langle w,Py\rangle$) gives
$d\,\langle iw,y_0\rangle=\gamma\,\langle iw,\eta\rangle$ with $d=\prod_{j\ne s}(\lambda_j-\lambda_s)$ and $\gamma=\prod_j(\mu_j-\lambda_s+\frac12)$. By hypothesis
$\gamma\ne0$, and $\langle iw,\eta\rangle\ne0$; hence $\langle iw,y_0\rangle\ne0$.

Let $b=\mathrm{pr}_s(e_z\otimes\eta)$. The form identity of \ref{pf:fs:W4} (with $y=e_z\otimes\eta$, $\eps_z=+1$) and the equivariance of $p_s$ and
$\mathrm{pr}_s$ give, for every $h$, $c_s\langle b,\pi_s(\tilde h)b\rangle=\langle\eta,\pi(\tilde h)y_0\rangle$, since $\tilde h$ fixes $e_z$. Multiplying by
$\overline{\langle w,\sigma(h)w\rangle}$, integrating and using \eqref{eq:pf:fs:eigproj} with $a=y_0$,
\[
c_s\int\overline{\langle w,\sigma(h)w\rangle}\,\langle b,\pi_s(\tilde h)b\rangle\,dh=\kappa\,\langle\eta,iw\rangle\langle iw,y_0\rangle\neq0 .
\]
Hence the integral on the left is non-zero, and \ref{pf:fs:W2}, applied to the data $(\lambda-e_s,\mu)$ and the discrete series $\pi_s$ and
$\sigma$, gives a non-zero intertwiner from $\sigma$ to the restriction of $\pi_s$. By \ref{pf:fs:W1}, occurrence holds at $(\lambda-e_s,\mu)$.
\end{proof}

\begin{theorem}[Occurrence]\label{pf:fs:occ}
If $\mathrm{Bc}(e,\eps,\lambda,\eps',\mu)$ holds and $(\eps,\lambda)$, $(\eps',\mu)$ are interlaced, then occurrence holds at $(\lambda,\mu)$: every
discrete series $\sigma$ of $\Uni(J_{\eps'})$ with parameter $\mu$ admits a non-zero continuous intertwiner into the restriction to
$\Uni(J_{\eps'})$ of every discrete series of $\Uni(J_\eps)$ with parameter $\lambda$.
\formallabel{PositivityTheory.RgHe.occ}
\end{theorem}

\begin{proof}
The statements \ref{pf:fs:W1}--\ref{pf:fs:W4} hold, so (L) holds by \cref{pf:fs:stepL}. The statement (M) is deduced from \ref{pf:fs:W1}--\ref{pf:fs:W4} by
the parallel argument for the lowering of an entry of $\mu$ (\lref{PositivityTheory.RgHe.SM.stepM}). The base statement (B) holds
(\lref{PositivityTheory.RgHe3.base}, using \ref{pf:fs:W1} and a frame statement \lref{PositivityTheory.RgHe.FrameF.frame}: after shifting
the parameters monotonically so that consecutive values are far apart, occurrence is established by an asymptotic argument with trial vectors in
an explicit frame). \Cref{pf:fs:walk} concludes.
\end{proof}

\begin{proposition}[Moment estimate]\label{pf:fs:moment}
Let $\mathrm{Bc}_B(e,\eps,\lambda,\eps',\mu)$ hold, let $\pi$ and $\sigma$ be discrete series with parameters $\lambda$ and $\mu$, and $dh$ a Haar measure on
$\Uni(J_{\eps'})$. Let $P\colon H\to H'$ be a continuous linear map with $\|P\|\le1$ and $P\pi(\tilde h)=\sigma(h)P$, $w\in H'$ a unit vector with
$\|P^*w\|=1$ and $\kappa>0$, such that: $h\mapsto\langle w,\sigma(h)w\rangle$ is integrable; $\int\overline{\langle w,\sigma(h)w\rangle}\,
\overline{\langle w,\sigma(h^{-1}x)w\rangle}\,dh=\kappa\,\overline{\langle w,\sigma(x)w\rangle}$ for all $x$; and there is a smooth $\eta\in H$ with
$\langle\eta,P^*w\rangle\ne0$ and \eqref{eq:pf:fs:eigproj} with $i w$ replaced by $P^*w$. Then for every lowest $K$-type unit vector $v$ of $\pi$, with
$\xi=P^*w$ and $T=8(n+1)B$,
\[
\int\mathrm{col}(g^{-1})^2\,|\langle\xi,\pi(g)v\rangle|^2\,dg\le T^{8(n+1)}\int|\langle\xi,\pi(g)v\rangle|^2\,dg .
\]
\formallabel{PositivityTheory.RgHe.A.Pf.momentP}
\formallabel{PositivityTheory.UnitaryDecomposition.moment_higher_estimate_holds}
\formallabel{PositivityTheory.UnitaryDecomposition.momentHeA_of}
\end{proposition}

\begin{proof}
We use \ref{pf:fs:W4} in the form of a decomposition $(S,p_s,\mathrm{pr}_s,\pi_s,c_s)$ of the smooth vectors of $\pi\otimes\C^{n+1}$
(\lref{PositivityTheory.UnitaryDecomposition.DecB.unitary_decomposition_exists}), and the following facts.
\begin{enumerate}[label=(\alph*)]
\item (comatrix bounds; \lref{PositivityTheory.UnitaryDecomposition.Com.comatrix_word_bound}) For $s\in S$: $\|p_s(e_l\otimes v)_c\|\le T^{4(n+1)}\|v\|$
for lowest $K$-type vectors $v$ and all $l,c$; and there is $\gamma_s\in\C$ with $|\gamma_s|\le T^{n+1}$ such that
$P(p_s(e_z\otimes\eta')_z)=\gamma_sP(\eta')$ for every smooth $\eta'$ and every contraction $P$ intertwining the restriction of $\pi$ with $\sigma$.
\item (smoothness; \lref{CompanionTheory.Lime.T3.kFiniteWordBound_indefUnitary}) Lowest $K$-type vectors are smooth, with explicit bounds for
the Lie algebra action.
\item (formal degrees; \lref{CompanionTheory.D.Upq.Coef.fdeg_gap}, \lref{PositivityTheory.UnitaryDecomposition.fdeg_shift}) There is
$c_0>0$, depending on $(n,\eps)$, such that every discrete series with Harish-Chandra parameter $\lambda'$ and gap at least $120(n+1)^3$ has formal
degree $d(\lambda')=c_0\,\mathrm{fd}(\lambda')$ for $dg$, where $\mathrm{fd}$ is Harish-Chandra's polynomial; and
$\mathrm{fd}(\lambda)\le2^{2(n+1)}\mathrm{fd}(\lambda-e_s)$ when $\lambda-e_s$ is a Harish-Chandra parameter.
\item (Schur orthogonality) For a discrete series of formal degree $d$, $\int|\langle y,\pi(g)x\rangle|^2\,dg=\|x\|^2\|y\|^2/d$, the integrand
being integrable.
\end{enumerate}
Let $d=d(\lambda)$ and $d_s=d(\lambda-e_s)$; the gap of $\lambda-e_s$ is at least that of $\lambda$ minus $1$. By (c), $1/d_s\le T^{n+1}/d$, as
$2^{2(n+1)}\le T^{n+1}$ ($B\ge1$ follows from the gap conditions). By (d), $\int|\langle\xi,\pi(g)v\rangle|^2dg=1/d$.

\emph{Step 1: an average.} Let $b_s=\mathrm{pr}_s(e_z\otimes\eta)$. Using the convolution identity for $w$, averaging against the coefficient of $w$
produces a vector $b'_s$ with $\langle b'_s,y\rangle=\int\overline{\langle w,\sigma(h)w\rangle}\langle b_s,\pi_s(\tilde h)y\rangle\,dh$ for all $y$ and
$\langle b'_s,b'_s\rangle=\kappa\langle b'_s,b_s\rangle$ (\lref{PositivityTheory.UnitaryDecomposition.idem_avg}).

\emph{Step 2: the components of $\xi$.} For smooth $x\in H^{n+1}$,
\begin{equation}\label{eq:pf:fs:momentmain}
\kappa\,\langle\eta,\xi\rangle\,\langle\xi,p_s(x)_z\rangle=c_s\,\langle b'_s,\mathrm{pr}_s(x)\rangle .
\end{equation}
Indeed, by \eqref{eq:pf:fs:eigproj} with $a=p_s(x)_z$, the left side is $\int\overline{\langle w,\sigma(h)w\rangle}\langle\eta,\pi(\tilde h)p_s(x)_z\rangle\,dh$; since
$\tilde h$ fixes $e_z$, $\pi(\tilde h)p_s(x)_z$ is the $z$-component of $p_s$ applied to the translate of $x$ by $\tilde h$, and the form identity of
\ref{pf:fs:W4} with $y=e_z\otimes\eta$ (note $\eps_z=+1$) turns the integrand into $c_s\,\overline{\langle w,\sigma(h)w\rangle}\langle b_s,\pi_s(\tilde h)\mathrm{pr}_s(x)\rangle$.
Taking $x=e_z\otimes\eta$ in \eqref{eq:pf:fs:momentmain} and using (a) (with $\langle\xi,y\rangle=\langle w,Py\rangle$) gives
$\kappa\langle\eta,\xi\rangle\gamma_s\langle\xi,\eta\rangle=c_s\langle b'_s,b_s\rangle$, hence by Step 1
\[
|c_s|\,\|b'_s\|^2=\kappa^2|\langle\eta,\xi\rangle|^2|\gamma_s| .
\]
Taking $x=y=e_l\otimes v$ in the form identity and using (a): $|c_s|\,\|\mathrm{pr}_s(e_l\otimes v)\|^2\le T^{4(n+1)}$.

\emph{Step 3: one component.} Let $x_l(g)$ be the translate of $e_l\otimes v$ by $g$. By \eqref{eq:pf:fs:momentmain} and the equivariance of
$\mathrm{pr}_s$,
\[
|\langle\xi,p_s(x_l(g))_z\rangle|^2=\Bigl(\frac{|c_s|}{\kappa|\langle\eta,\xi\rangle|}\Bigr)^2|\langle b'_s,\pi_s(g)\mathrm{pr}_s(e_l\otimes v)\rangle|^2,
\]
and by (d) for $\pi_s$ and Step 2,
\begin{align*}
\int|\langle\xi,p_s(x_l(g))_z\rangle|^2dg&=\frac{|c_s|^2\|b'_s\|^2\|\mathrm{pr}_s(e_l\otimes v)\|^2}{\kappa^2|\langle\eta,\xi\rangle|^2d_s}\\
&=\frac{|\gamma_s|\,|c_s|\,\|\mathrm{pr}_s(e_l\otimes v)\|^2}{d_s}\le\frac{T^{n+1}\,T^{4(n+1)}\,T^{n+1}}{d}.
\end{align*}

\emph{Step 4: summation.} The $z$-component of $x_l(g)$ is $(g^{-1})_{lz}\pi(g)v$, and $\sum_{s\in S}p_s=\mathrm{id}$ on smooth vectors, so
$(g^{-1})_{lz}\langle\xi,\pi(g)v\rangle=\sum_{s\in S}\langle\xi,p_s(x_l(g))_z\rangle$. By the Cauchy--Schwarz inequality and $|S|\le n+1\le T$,
\[
\mathrm{col}(g^{-1})^2|\langle\xi,\pi(g)v\rangle|^2=\sum_l|(g^{-1})_{lz}|^2|\langle\xi,\pi(g)v\rangle|^2\le T\sum_l\sum_{s\in S}|\langle\xi,p_s(x_l(g))_z\rangle|^2 .
\]
Integrating with Step 3 and summing over at most $n+1\le T$ values of $l$ and of $s$,
\[
\int\mathrm{col}(g^{-1})^2|\langle\xi,\pi(g)v\rangle|^2dg\le T\cdot T\cdot T\cdot\frac{T^{6(n+1)}}{d}=T^{6(n+1)+3}\int|\langle\xi,\pi(g)v\rangle|^2dg,
\]
and $T^{6(n+1)+3}\le T^{8(n+1)}$ since $T\ge1$ and $n\ge2$.
\end{proof}

\begin{theorem}[Analytic estimates]\label{pf:fs:analytic}
Let $\mu_\eps$ denote a Haar measure on $\Uni(J_\eps)$ and $\mathcal U_\delta=\{g:\|g-1\|\le\delta\}$.
\begin{enumerate}[label=(\alph*)]
\item (volumes) For $m\ge1$, $0<\delta\le2$ and $R\ge0$: $\mu_\eps\{g:\|g\|\le R\}\le(4m(\sqrt m+1)R/\delta)^{2m^2}\mu_\eps(\mathcal U_\delta)$.
\item (decay) Let $1\le m\le N$, $\pi$ a discrete series of $\Uni(J_\eps)$ ($m$ variables) with Harish-Chandra parameter $\lambda$ and gap at least
$200N^3$, $v$ a lowest $K$-type unit vector and $\|g_0\|^2\le Y$. Then for all $g$
\[
|\langle\pi(g_0)v,\pi(g)\pi(g_0)v\rangle|\le\exp(6N^2\log m+12N^2)\,Y^{12N^2}\,\|g\|^{-12N^2}.
\]
\item (integrals) For $m\ge1$, $a\ge m^2+1$ and continuous $\phi$ with $|\phi(g)|\le C\|g\|^{-2a}$: $\phi$ is integrable and
$\int|\phi|\,d\mu_\eps\le C\cdot2(16m^2)^{2m^2}\mu_\eps(\mathcal U_1)$.
\item (restricted integrals) Let $n\ge2$, $\eps_z=+1$, $\eps_{e(i)}=\eps'_i$, $C\ge0$, $0<\delta\le1$ and $\phi$ with $|\phi(g)|\le C\|g\|^{-6(n+1)^2}$ on
$\Uni(J_\eps)$. Then for $g\in\Uni(J_\eps)$,
$\int|\phi(g\tilde h)|\,d\mu_{\eps'}(h)\le C(100(n+1)^2/\delta)^{10(n+1)^2}\mathrm{col}(g)^{-(2n+7)}\mu_{\eps'}(\mathcal U_\delta)$.
\item (formal degrees) Let $m\ge2$ and $\pi$ a discrete series of $\Uni(J_\eps)$ with Harish-Chandra parameter $\lambda$, $|\lambda_a|\le B$, gap at least
$100m^3$ and formal degree $d$ for $\mu_\eps$. Then $d\,\mu_\eps(\mathcal U_\delta)\le4$ for $\delta=(100m(B+1))^{-2}$.
\item (Lipschitz) Let $n\ge2$, $\eps,\eps'$ as in (d), $\pi$ a discrete series of $\Uni(J_\eps)$ with Harish-Chandra parameter $\lambda$, $|\lambda_a|\le B$,
gap at least $200(n+1)^3$, $v$ a lowest $K$-type unit vector, $\|g_0\|^2\le Y$, $0<\delta\le1$ and $X$ in the Lie algebra of $\Uni(J_{\eps'})$.
Then, with $f(h)=\langle\pi(g_0)v,\pi(\tilde h)\pi(g_0)v\rangle$,
\[
\int|f(\exp(-X)h)-f(h)|\,d\mu_{\eps'}(h)\le(8(n+1)BY/\delta)^{20(n+1)^2}\,\|X\|\,\mu_{\eps'}(\mathcal U_\delta).
\]
\end{enumerate}
\formallabel{CompanionTheory.Lime.T3.analytic_estimates}
\end{theorem}

\begin{proof}
Elements of $\Uni(J_\eps)$ satisfy $\|g\|\ge1$ and $\|g^{-1}\|=\|g\|$ ($g^{-1}=J_\eps g^*J_\eps$).

(a) Let $\lambda_0$ be Lebesgue measure on the $2m^2$-dimensional real space of complex $m\times m$ matrices, $O$ the ball of radius
$s=\delta/(4m)$ around $1$ for the maximum of the absolute values of the entries (so $\|x-1\|\le\delta/4$ on $O$), $S_R=\{g:\|g\|\le R\}$ and
$T=\{(g,x):g\in S_R,\ g^{-1}x\in O\}$. For fixed $g$, the slice $\{x:(g,x)\in T\}$ is $gO$, and left multiplication by $g$ has determinant $1$ on
this real space; so $(\mu_\eps\times\lambda_0)(T)=\mu_\eps(S_R)\lambda_0(O)$. For fixed $x$, if the slice $\{g:(g,x)\in T\}$ contains some $h$, then all entries
of $x=h(h^{-1}x)$ are bounded by $\rho=R(\sqrt m+1)$, and every $g$ in the slice satisfies $\|h^{-1}g-1\|\le\delta$ (as $(h^{-1}g)(g^{-1}x)=h^{-1}x$ with
$\|g^{-1}x-1\|,\|h^{-1}x-1\|\le\delta/4\le\frac12$); so the slice lies in $h\,\mathcal U_\delta$ and has measure at most $\mu_\eps(\mathcal U_\delta)$. Hence
$(\mu_\eps\times\lambda_0)(T)\le\mu_\eps(\mathcal U_\delta)\lambda_0(\text{ball of radius }\rho)$, and dividing by $\lambda_0(O)$ gives the bound with
$(\rho/s)^{2m^2}=(4m(\sqrt m+1)R/\delta)^{2m^2}$.

(b) Lowest $K$-type coefficients satisfy $|\langle v,\pi(g)v\rangle|\le(e\sqrt m)^A\|g\|^{-A}$ for $1\le A\le12N^2$
(\lref{CompanionTheory.Lime.T3.dec_low}: in the compact case $\|g\|=\sqrt m$; otherwise write $g=k\exp(X)$ with $k$ in the maximal compact
subgroup and $X$ hermitian, so that $\|g\|\le\sqrt me^{\|X\|}$, and use the exponential decay of lowest $K$-type coefficients of discrete series with
large gap, \lref{ReciprocityTheory.ChartRestriction.DS.coef_dec_of_isDS}). Since
$\langle\pi(g_0)v,\pi(g)\pi(g_0)v\rangle=\langle v,\pi(g_0^{-1}gg_0)v\rangle$ and $\|g\|\le\|g_0\|^2\|g_0^{-1}gg_0\|$, the case $A=12N^2$ gives the bound,
with $(e\sqrt m)^{12N^2}=\exp(6N^2\log m+12N^2)$ and $\|g_0\|^{24N^2}\le Y^{12N^2}$.

(c) As $\|g\|\ge1$ and $2a\ge2m^2+1$, $|\phi(g)|\le C\|g\|^{-(2m^2+1)}$. For $A\ge2m^2+1$ and $0<\delta\le2$,
$\int\|g\|^{-A}d\mu_\eps\le2(8m(\sqrt m+1)/\delta)^{2m^2}\mu_\eps(\mathcal U_\delta)$ (\lref{CompanionTheory.Lime.T3.int_ratio}, obtained from (a) by
decomposing $\Uni(J_\eps)$ into the shells $2^k\le\|g\|<2^{k+1}$ and summing a geometric series). With $\delta=1$ and $\sqrt m+1\le2m$ this gives the
integrability and the bound; for general $0<\delta\le1$ the same argument gives $\int|\phi|\le C\cdot2(16m^2/\delta)^{2m^2}\mu_\eps(\mathcal U_\delta)$.

(d) Let $a=2n+7$, $b=2n^2+1$ and $y=\mathrm{col}(g)^2$. Choose $h_1$ with $\|g\tilde h_1\|\le2\inf_h\|g\tilde h\|$. For all $h$, with $t=\|g(h_1h)^\sim\|$:
$\|g\tilde h_1\|\le2t$; $\|h\|\le\|\tilde h\|\le\|g\tilde h_1\|\,t$ (as $\tilde h=(g\tilde h_1)^{-1}g(h_1h)^\sim$); and $\sqrt y\le\|g\tilde h_1\|$, because $\tilde h_1$ fixes
$e_z$, so the column $z$ of $g\tilde h_1$ is that of $g$. Hence
$t^{-(a+2b)}\le(2/\|g\tilde h_1\|)^{a+b}(\|g\tilde h_1\|/\|h\|)^b\le2^{a+b}y^{-a/2}\|h\|^{-b}$, and since $a+2b\le6(n+1)^2$ and $t\ge1$,
$|\phi(g(h_1h)^\sim)|\le C2^{a+b}y^{-a/2}\|h\|^{-b}$. By left invariance of $\mu_{\eps'}$ and (c) on $\Uni(J_{\eps'})$ ($n$ variables, exponent $b$),
\[
\int|\phi(g\tilde h)|\,d\mu_{\eps'}(h)=\int|\phi(g(h_1h)^\sim)|\,d\mu_{\eps'}(h)\le C2^{a+b}\mathrm{col}(g)^{-a}\cdot2(16n^2/\delta)^{2n^2}\mu_{\eps'}(\mathcal U_\delta),
\]
and $2^{a+b+1}(16n^2/\delta)^{2n^2}\le(100(n+1)^2/\delta)^{10(n+1)^2}$.

(e) Let $v$ be a unit vector of the lowest $K$-type. It is smooth, and $\|d\pi(X)v\|\le8\sqrt m(B+2)$ for $\|X\|\le1$
(\lref{CompanionTheory.Lime.T3.low_bounds}); a smooth vector with this property satisfies $\|\pi(g)v-v\|\le7\delta\cdot8\sqrt m(B+2)$ for
$g\in\mathcal U_\delta$, $\delta\le\frac12$ (\lref{CompanionTheory.Lime.T3.disp_le}, by integrating the derivative along a curve from $1$ to $g$
obtained from the Cayley transform). For $\delta=(100m(B+1))^{-2}$ the right side is at most $\frac12$, so $|\langle v,\pi(g)v\rangle-1|\le\frac12$ and
$|\langle v,\pi(g)v\rangle|^2\ge\frac14$ on $\mathcal U_\delta$. By the definition of the formal degree, $\int|\langle v,\pi(g)v\rangle|^2d\mu_\eps=1/d$, hence
$\mu_\eps(\mathcal U_\delta)/4\le1/d$.

(f) Put $r=\|X\|$, $X'=X/r$ (if $r>0$), $c(t)=\exp(tX')$, $b=2n^2+1$, $T=8(n+1)BY/\delta$, $P=e\sqrt{n+1}\,Y$, $Q=\sqrt n\,e^2$ and $D=(PQ)^{2b}$. The
gap and the bounds give $n+1\le4B$, $T\ge18$ and $PQ\le T^2$, so $D\le T^{4b}$. By the bounds used in (b), with $A=2b$, and $\|k\|\le\|\tilde k\|$,
$|f(k)|\le P^{2b}\|k\|^{-2b}$. Note $\exp(-X)h=c(r)^{-1}h$.

If $r\ge1$: $|f(c(r)^{-1}h)-f(h)|\le F(c(r)^{-1}h)+F(h)$ with $F(k)=D\|k\|^{-b}$, and by left invariance and the bound of (c),
$\int|f(c(r)^{-1}h)-f(h)|\,dh\le2D\cdot2(16n^2/\delta)^{2n^2}\mu_{\eps'}(\mathcal U_\delta)\le T^{20(n+1)^2}r\,\mu_{\eps'}(\mathcal U_\delta)$.

If $r<1$: consider the one-parameter group $d(t)=g_0^{-1}c(t)^\sim g_0$ with generator $Z=g_0^{-1}\tilde X'g_0$, $\|Z\|\le\|g_0\|^2\le Y$. By
\lref{CompanionTheory.Lime.T3.low_bounds}, $\|d\pi(Z)^2v\|\le\|Z\|^2(8\sqrt{n+1}(B+2))^2\le T^{4b}$. For fixed $h$ let
$u=\pi(g_0^{-1}\tilde hg_0)v$; then $\langle\pi(d(t))v,u\rangle=f(c(t)^{-1}h)$, and for $t\in[0,2]$, $\|c(t)\|\le Q$ gives $\|h\|\le Q\|c(t)^{-1}h\|$ and
$|f(c(t)^{-1}h)|\le D\|h\|^{-2b}$. The interpolation inequality for a smooth vector along a one-parameter group
(\lref{CompanionTheory.Lime.T3.coef_interp}: if $|\langle\pi(d(t))x,u\rangle|\le\eta$ on $[0,2]$ and $\|d\pi(Z)^2x\|\le M$, then
$|\langle\pi(d(a))x,u\rangle-\langle x,u\rangle|\le(2\eta/\tau+M\tau)a$ for $0<\tau\le1$ and $a\in[0,1]$), applied with $\tau=\|h\|^{-b}$, $a=r$, gives
$|f(c(r)^{-1}h)-f(h)|\le(2D+T^{4b})r\|h\|^{-b}\le3T^{4b}r\|h\|^{-b}$. Integrating with (c) on $\Uni(J_{\eps'})$,
$\int|f(c(r)^{-1}h)-f(h)|\,dh\le3T^{4b}r\cdot2(16n^2/\delta)^{2n^2}\mu_{\eps'}(\mathcal U_\delta)\le T^{20(n+1)^2}r\,\mu_{\eps'}(\mathcal U_\delta)$, using
$16n^2/\delta\le T^2$.
\end{proof}

\subsubsection{Representations of the compact unitary group}

A \emph{Harish-Chandra type model} with parameter $u\in\R^k$ is a unitary representation $\rho'$ of $\Uni(k)$ on a non-zero finite-dimensional
Hilbert space which is a $K$-type with highest weight $\mathrm{lowK}(u)$, the lowest $K$-type attached to the Harish-Chandra parameter $u$, in the sense
of the formal definition.

\begin{lemma}[Highest weight types]\label{pf:fs:hcType}
Every irreducible continuous unitary representation $\rho$ of $\Uni(k)$ on a non-zero finite-dimensional Hilbert space $W_p$ is isomorphic to a
Harish-Chandra type model: there are $u\in\R^k$, a Harish-Chandra type model $\rho'$ with parameter $u$ on $W_q$ and a linear isomorphism
$T\colon W_p\to W_q$ with $T\rho(h)=\rho'(h)T$ for $h\in\Uni(k)$.
\formallabel{UpperTheory.Lit.Bur.irreducible_unitary_rep_iso_harishChandraType}
\formallabel{UpperTheory.Lit.Bur.HWSpanType.exists_hcType_of_irred}
\end{lemma}

\begin{proof}
This is the theorem of the highest weight. The representation is smooth, and its differential extends to a representation $E$ of
$\mathfrak{gl}_k(\C)$ on $W_p$ satisfying the commutation relations, the elements of the Lie algebra of $\Uni(k)$ acting skew-adjointly
(\lref{UpperTheory.Lit.Bur.HWSpanType.Derived.exists_derived_action}). Since $W_p\ne0$, there are a weight $\mu$ and a non-zero vector
$x$ of weight $\mu$ annihilated by the raising operators (\lref{UpperTheory.Lit.Bur.KShift.exists_highest}). Put $u=\mu+\rho_k$, $\rho_k$ the
half sum of the positive roots; then $\mathrm{lowK}(u)=\mu$ (\lref{GlobalTheory.KTypeCompare.lowK_allPos}). The subspace $S$ generated by $x$ under
the action of the compact Lie algebra is stable under $\rho$ (\lref{UpperTheory.Lit.Bur.HWSpanType.Wt.stable_of_derived}), and the restriction
of $\rho$ to $S$ is a Harish-Chandra type model with parameter $u$ (\lref{UpperTheory.Lit.Bur.HWSpanType.exists_type_of_highest}: the space
generated by a highest weight vector of weight $\mathrm{lowK}(u)$ is a $K$-type of highest weight $\mathrm{lowK}(u)$). As $S$ contains $x\ne0$ and
$\rho$ is irreducible, $S=W_p$; take $W_q=S$, $\rho'$ the restriction of $\rho$ and $T$ the identity.
\end{proof}

% ---------------------------------------------------------------------------
\subsection{Remarks on the formal statements}\label{sec:pf:fs:remarks}

\begin{remark}[Unused hypotheses]\label{pf:fs:rem:unused}
\begin{enumerate}[label=(\alph*)]
\item In $\sfname{FunctionalsSeparate}$ the hypotheses $\sfname{ESerTie}(D)$ and the group record of the real datum $R$ are not used
by the proof (\cref{pf:fs:besselSep}). Of the conditions defining a Weil datum over $R$, only two are used: the family $\theta$ of $W$ is
that of $R$, and the Cartan involution of $W$ at $(0,0)$ is that of the model $\mathcal M_Y(0)$.
\item The formal statement of $\fsP{TS}$ (\cref{pf:fs:topSep}) takes $\sfname{ThetaLocalInputs}$ as a hypothesis which its proof does not
use.
\item \Cref{pf:fs:mainRI} takes the archimedean input $\fsP{RI}$ as a hypothesis; it is discharged by \cref{pf:fs:RIholds}, which is
unconditional. Thus \cref{pf:fs:main} has exactly the three hypotheses of \cref{prop:families:fs}.
\end{enumerate}
\end{remark}

\begin{remark}[What the hypotheses are used for]\label{pf:fs:rem:name}
The formal name of \cref{pf:fs:main} mentions Howe's unramified statement and the Siegel--Weil formula, but the proof uses many fields
of $\sfname{ThetaLocalInputs}$ as well (existence of discrete series of $\Uni(p,q)$, the discrete series cover, Harish-Chandra modules
and finiteness, Casselman--Wallach globalisation, rapid decay of cusp forms, the Weil representation of unitary dual pairs, the
Weil--Siegel formula, the doubling method and the unramified doubling integral) and the archimedean input proved in
\cref{sec:pf:fs:real}. $\sfname{HoweUnram}$ and $\sfname{SiegelWeilUnitaryConvergent}$ enter only through the reachability statement
$\fsP{CO}$: the first in the transfer of Hecke parameters by upward lifts \ref{pf:fs:O5}, the second in the Rallis inner product formula
\ref{pf:fs:O7}.
\end{remark}

\begin{remark}[Thresholds and quantifiers]\label{pf:fs:rem:thresholds}
\begin{enumerate}[label=(\alph*)]
\item The constant $C_F$ is the depth constant $c_{\mathrm{dep}}$ of the grid. The hypothesis $C_F\ge6510$ is used only in the proof of
$\fsP{DP}$ (\lref{GlobalTheory.T.gridDepth_pays_every_loss}). The statement $\fsP{SG}$ holds with threshold $x_0=0$: the size
conditions of rows imply $\sfname{Sizes}(D)$ at every scale. Likewise the statement on the grid parameters (\cref{pf:fs:gridPar}), and hence
the statements on local pack data (\cref{pf:fs:locOK}) and on packs (\cref{pf:fs:packEx}), hold with threshold $x_0=0$ and do not use
$C_F\ge6510$.
\item In $\sfname{FunctionalsSeparate}$ the pair $(\rho,\xi)$ is chosen independently of $t$, while $(\Psi,\Phi)$ depend on $t$. When
$\fsE_R=0$ the statement holds with $\xi=0$; when $\fsE_R\ne0$ any separating $\xi$ is non-zero.
\item The statement $\fsP{PN}$ quantifies over all models of $Y_0$ with parameter $S^Y$, without the condition on the Cartan involution
imposed in the other statements on rows.
\item The uniqueness statement $\fsP{BG}$ is the case of a Bessel frame with $\langle z,z\rangle<0$ and $J_0$ positive definite. The
development also states the variant allowing the opposite signs ($\langle z,z\rangle>0$, $J_0$ negative definite); only the first is used here.
\end{enumerate}
\end{remark}

